\chapter{The Arithmetic--Geometric Mean}

\begin{center}\itshape
We have established that the arithmetic--geometric mean between $1$ and $\sqrt2$ is $\pi/\varpi$ to the eleventh decimal place; the
demonstration of this fact will surely open up an entirely new field of analysis.\\
\upshape --- Carl Friedrich Gauss, mathematical diary, 30 May 1799
\end{center}

Take two positive numbers, and replace them by their arithmetic mean and their geometric mean. Repeat. The two numbers approach
each other very fast: the number of correct digits doubles at every step. Their common limit, the \emph{arithmetic--geometric mean}, looks like
a curiosity. Gauss discovered at the age of twenty-one that it computes an integral, the one that gives the length of an ellipse, and through it the
number $\pi$. This chapter proves his theorem with a substitution found by Newman in 1985, computes the perimeter of an ellipse, proves
Legendre's relation between elliptic integrals, and combines everything into the algorithm of Brent and Salamin (1976), which computes $\pi$
with the number of correct digits doubling at every step. It closes the volume: the Gamma function, Wallis's product, Stirling, the lemniscate of
Chapter~7 and the fast multiplication of Chapter~9 all meet here.

\begin{think}
Start from $1$ and $\sqrt2$. After one step you have $1.2071$ and $1.1892$; after four, both are $1.19814023473559\ldots$ How many steps would you
expect to need for $1000$ digits? Gauss compared this number with $2\int_0^1\frac{\mathrm dt}{\sqrt{1-t^4}}=2.62205755\ldots$, the length of a half-lemniscate
(Chapter~7, Exercise~\ref{ex:lemniscate}). Compute $\pi$ divided by it.
\end{think}

\section{The arithmetic--geometric mean}

Let $a\ge b>0$. Define $a_0=a$, $b_0=b$ and
\[
\begin{gathered}a_{n+1}=\frac{a_n+b_n}2,\qq b_{n+1}=\sqrt{a_nb_n},\end{gathered}
\]
and $c_n=\sqrt{a_n^2-b_n^2}\ge0$.

\begin{proposition}\label{prop:agm}
\begin{enumerate}[label=(\alph*)]
\item $b\le b_n\le b_{n+1}\le a_{n+1}\le a_n\le a$ for every $n$, and $(a_n)$, $(b_n)$ have a common limit $M(a,b)$, the \emph{arithmetic--geometric mean} of $a$ and $b$.
\item $c_{n+1}=\frac{a_n-b_n}2$ and $c_{n+1}=\frac{c_n^2}{4a_{n+1}}$. Moreover $0\le a_{n+1}-b_{n+1}\le\frac{(a_n-b_n)^2}{8b}$.
\item $M(\lambda a,\lambda b)=\lambda M(a,b)$ for $\lambda>0$, and $M(a,b)=M\bigl(\frac{a+b}2,\sqrt{ab}\bigr)$.
\end{enumerate}
\end{proposition}
\begin{proof}
(a) The arithmetic--geometric inequality $\sqrt{a_nb_n}\le\frac{a_n+b_n}2$ gives $b_{n+1}\le a_{n+1}$, and since $b_n\le a_n$, both means lie between $b_n$ and $a_n$.
So $(a_n)$ decreases, $(b_n)$ increases, and both are bounded; they converge, to $A$ and $B$ say. Passing to the limit in $a_{n+1}=\frac{a_n+b_n}2$ gives $A=\frac{A+B}2$, so $A=B$.

(b) $c_{n+1}^2=\frac{(a_n+b_n)^2}4-a_nb_n=\frac{(a_n-b_n)^2}4$, and $\frac{a_n-b_n}2=\frac{a_n^2-b_n^2}{2(a_n+b_n)}=\frac{c_n^2}{4a_{n+1}}$. Finally
$a_{n+1}-b_{n+1}=\frac{(\sqrt{a_n}-\sqrt{b_n})^2}2=\frac{(a_n-b_n)^2}{2(\sqrt{a_n}+\sqrt{b_n})^2}$, and $(\sqrt{a_n}+\sqrt{b_n})^2\ge4b_n\ge4b$.

(c) The sequences built from $(\lambda a,\lambda b)$ are $(\lambda a_n,\lambda b_n)$. The sequences built from $(a_1,b_1)$ are those built from $(a,b)$, shifted by one.
\end{proof}

Part (b) is the reason for the speed. Once $a_n-b_n$ is small compared with $8b$, the next difference is of the order of its square: if $a_n-b_n\approx10^{-k}$,
then $a_{n+1}-b_{n+1}\approx10^{-2k}$. For $M(1,\sqrt2)$ the differences are $0.41$, $0.018$, $3.3\cdot10^{-5}$, $1.2\cdot10^{-10}$, $1.4\cdot10^{-21}$, $2\cdot10^{-43}$:
each step doubles the number of correct digits, and ten steps give a thousand digits. Before that, while $a_n$ is much larger than $b_n$, the
difference only halves at each step (Figure~\ref{fig:agm}).

\begin{figure}[ht]
\centering
\fig{ch12-agm}
\caption{The AGM iteration from $a_0=1$, $b_0=0.01$. The differences $a_n-b_n$ are $0.99$, $0.405$, $0.078$, $2.9\cdot10^{-3}$, $4.0\cdot10^{-6}$, $7.5\cdot10^{-12}$: slow
while the two numbers have different sizes, then quadratic. The limit is $M(1,0.01)=0.2621668872\ldots$}
\label{fig:agm}
\end{figure}

\begin{think}
Why does squaring the error double the number of correct digits? Newton's method for $\sqrt2$ has the same property; what do the two
algorithms have in common? (Look at $a_{n+1}-b_{n+1}$ as a function of $a_n-b_n$ near $0$.)
\end{think}

\section{Gauss's theorem}

For $a,b>0$ let
\[
I(a,b)=\int_0^{\pi/2}\frac{\mathrm d\theta}{\sqrt{a^2\cos^2\theta+b^2\sin^2\theta}}.
\]
The substitution $\theta\mapsto\frac\pi2-\theta$ shows $I(a,b)=I(b,a)$, $I(\lambda a,\lambda b)=\frac1\lambda I(a,b)$, and $I(a,a)=\frac\pi{2a}$.

\begin{theorem}[Gauss, 1799]\label{thm:gaussagm}
For $a,b>0$,
\[
I(a,b)=\frac\pi{2M(a,b)}.
\]
\end{theorem}
\begin{proof}
\emph{Step 1: an algebraic form.} Substitute $t=b\tan\theta$, a bijection from $\left[0,\frac\pi2\right[$ onto $[0,\infty[$, with $\mathrm dt=b(1+\tan^2\theta)\dd\theta$.
Then $a^2\cos^2\theta+b^2\sin^2\theta=\cos^2\theta\,(a^2+t^2)$ and $1+\tan^2\theta=\frac1{\cos^2\theta}=\frac{b^2+t^2}{b^2}$, so
\[
\begin{aligned}I(a,b)&=\int_0^\infty\frac{\mathrm dt}{\sqrt{(a^2+t^2)(b^2+t^2)}}\eqbr=\frac12\int_{-\infty}^{\infty}\frac{\mathrm dt}{\sqrt{(a^2+t^2)(b^2+t^2)}}.\end{aligned}
\]

\emph{Step 2: Newman's substitution.} Let $a_1=\frac{a+b}2$, $b_1=\sqrt{ab}$, and substitute $t=\frac12\bigl(x-\frac{ab}x\bigr)$, an increasing bijection from $]0,\infty[$
onto $\R$, with $\mathrm dt=\frac12\bigl(1+\frac{ab}{x^2}\bigr)\mathrm dx=\frac{x^2+ab}{2x^2}\dd x$. Expanding the squares,
\[
\begin{gathered}t^2+b_1^2=\frac{(x^2+ab)^2}{4x^2},\qq t^2+a_1^2=\frac{(x^2+a^2)(x^2+b^2)}{4x^2}.\end{gathered}
\]
(For the second: $t^2+a_1^2=\frac14\bigl(x^2-2ab+\frac{a^2b^2}{x^2}+a^2+2ab+b^2\bigr)$.) Write $D=\sqrt{(x^2+a^2)(x^2+b^2)}$. Then
$\sqrt{(a_1^2+t^2)(b_1^2+t^2)}=\frac{(x^2+ab)D}{4x^2}$, so
\[
\begin{aligned}\frac{\mathrm dt}{\sqrt{(a_1^2+t^2)(b_1^2+t^2)}}&=\frac{x^2+ab}{2x^2}\cdot\frac{4x^2}{(x^2+ab)D}\dd x\eqbr=\frac{2\dd x}D.\end{aligned}
\]
By Step~1, $I(a_1,b_1)=\frac12\int_0^\infty\frac{2\,\mathrm dx}{\sqrt{(x^2+a^2)(x^2+b^2)}}=I(a,b)$.

\emph{Step 3: the limit.} By induction $I(a,b)=I(a_n,b_n)$ for every $n$, with the sequences of Proposition~\ref{prop:agm} (after swapping $a$ and $b$ if
needed). Since $b_n^2\le a_n^2\cos^2\theta+b_n^2\sin^2\theta\le a_n^2$, we have $\frac\pi{2a_n}\le I(a_n,b_n)\le\frac\pi{2b_n}$, and both bounds tend to $\frac\pi{2M(a,b)}$.
\end{proof}

The proof is short because all the work is in Step~2, which is just the right change of variable. Gauss's own proof used a more complicated trigonometric substitution; the one above appears in a short paper of D.~J.~Newman (1985).

In terms of the modulus $k\in[0,1[$ and the complementary modulus $k'=\sqrt{1-k^2}$, the \emph{complete elliptic integrals} of the first and second kind are
\[
\begin{gathered}K(k)=\int_0^{\pi/2}\frac{\mathrm d\theta}{\sqrt{1-k^2\sin^2\theta}},\qq E(k)=\int_0^{\pi/2}\sqrt{1-k^2\sin^2\theta}\dd\theta.\end{gathered}
\]
Since $1-k^2\sin^2\theta=\cos^2\theta+k'^2\sin^2\theta$, we have $K(k)=I(1,k')=\frac\pi{2M(1,k')}$.

\begin{corollary}[Landen's transformation]\label{cor:landen}
For $0\le k<1$, let $k_1=\frac{1-k'}{1+k'}$. Then $K(k)=(1+k_1)K(k_1)$.
\end{corollary}
\begin{proof}
By Step~2 above and homogeneity, $K(k)=I(1,k')=I\bigl(\frac{1+k'}2,\sqrt{k'}\bigr)=\frac2{1+k'}I\bigl(1,\frac{2\sqrt{k'}}{1+k'}\bigr)$. The last argument is the complementary
modulus of $k_1$, since $1-\frac{4k'}{(1+k')^2}=\frac{(1-k')^2}{(1+k')^2}$. And $1+k_1=\frac2{1+k'}$.
\end{proof}

Since $k_1\approx\frac{k^2}4$ for small $k$, iterating Landen's transformation drives the modulus to $0$ quadratically, where $K(0)=\frac\pi2$: this is the AGM again,
seen through the modulus.

\begin{corollary}[the lemniscate]\label{cor:lemniscate}
Let $\varpi=2\int_0^1\frac{\mathrm dr}{\sqrt{1-r^4}}$, so that the lemniscate $(x^2+y^2)^2=x^2-y^2$ has length $2\varpi$ (Chapter~7). Then
\[
\begin{gathered}\frac\varpi\pi=\frac1{M(1,\sqrt2)},\qq\varpi=2.62205755429211981\ldots\end{gathered}
\]
\end{corollary}
\begin{proof}
Substitute $r=\cos\theta$: $\mathrm dr=-\sin\theta\dd\theta$ and $1-\cos^4\theta=\sin^2\theta\,(1+\cos^2\theta)$, so
$\int_0^1\frac{\mathrm dr}{\sqrt{1-r^4}}=\int_0^{\pi/2}\frac{\mathrm d\theta}{\sqrt{1+\cos^2\theta}}=\int_0^{\pi/2}\frac{\mathrm d\theta}{\sqrt{2\cos^2\theta+\sin^2\theta}}=I(\sqrt2,1)$. By Theorem~\ref{thm:gaussagm} this is
$\frac\pi{2M(1,\sqrt2)}$.
\end{proof}

\begin{figure}[ht]
\centering
\fig{ch12-lemniscate}
\caption{The lemniscate $(x^2+y^2)^2=x^2-y^2$, in polar coordinates $r^2=\cos2\theta$. Along it, $\mathrm ds=\frac{\mathrm dr}{\sqrt{1-r^4}}$, so the rust arc from the centre to
$(1,0)$ has length $\int_0^1\frac{\mathrm dr}{\sqrt{1-r^4}}=\frac\varpi2$, and the whole curve has length $2\varpi$, as the circle of diameter $1$ has length $\pi$.}
\label{fig:lemniscate}
\end{figure}

\begin{beauty}
This is the identity of the epigraph. Gauss had computed both numbers to eleven decimals, and his diary entry records the moment he saw
they agreed, before he had a proof. He then developed the theory of the lemniscatic functions, the inverses of $\int\frac{\mathrm dr}{\sqrt{1-r^4}}$, which are
the first elliptic functions; the ``entirely new field of analysis'' he announced is the theory of elliptic functions and modular forms, the subject
of Volume~13. With Chapter~7, the corollary also computes $\Gamma(\frac14)$ by the AGM (Exercise~\ref{ex:gamma14}).
\end{beauty}

\section{The perimeter of an ellipse}

The ellipse with semi-axes $a\ge b>0$, parametrised by $(a\cos t,b\sin t)$, has speed $\sqrt{a^2\sin^2t+b^2\cos^2t}$, so its perimeter is $L(a,b)=4J(a,b)$, where
\[
J(a,b)=\int_0^{\pi/2}\sqrt{a^2\cos^2\theta+b^2\sin^2\theta}\dd\theta
\]
(substitute $\theta=\frac\pi2-t$). With $k$ the eccentricity, $b=ak'$ and $J(a,b)=aE(k)$. For a circle, $J(a,a)=\frac{\pi a}2$. In the rest of this section, $a,b>0$, we write
$R=a^2\cos^2\theta+b^2\sin^2\theta$, so that $I=\int_0^{\pi/2}R^{-1/2}$ and $J=\int_0^{\pi/2}R^{1/2}$, and the AGM sequences start from $a_0=a$, $b_0=b$.

\begin{lemma}\label{lem:jderiv}
\begin{enumerate}[label=(\alph*)]
\item $\displaystyle J(a,b)=a^2b^2\int_0^{\pi/2}R^{-3/2}\dd\theta$.
\item $I$ is of class $C^1$ on $]0,\infty[^2$, $a\partial_aI+b\partial_bI=-I$, and, for $a\ne b$,
\[
J(a,b)=a^2I(a,b)+a(a^2-b^2)\,\partial_aI(a,b).
\]
\end{enumerate}
\end{lemma}
\begin{proof}
(a) Write $c=\cos^2\theta$ and $s=\sin^2\theta$. The derivative of $\sin\theta\cos\theta\,R^{-1/2}$ is
$(c-s)R^{-1/2}-\sin\theta\cos\theta\cdot\frac12R^{-3/2}\cdot2(b^2-a^2)\sin\theta\cos\theta$, that is $\bigl((c-s)R-(b^2-a^2)cs\bigr)R^{-3/2}$. Using $c+s=1$, the
numerator is $a^2c^2-b^2s^2$, and so is $\frac{R^2-a^2b^2}{a^2-b^2}$ when $a\ne b$, because
$R^2-a^2b^2=a^4c^2+b^4s^2-a^2b^2(c^2+s^2)=(a^2-b^2)(a^2c^2-b^2s^2)$. So for $a\ne b$
\[
\frac{\mathrm d}{\mathrm d\theta}\Bigl(\frac{\sin\theta\cos\theta}{\sqrt R}\Bigr)=\frac{R^{1/2}-a^2b^2R^{-3/2}}{a^2-b^2}.
\]
The left side integrates to $0$ over $[0,\frac\pi2]$, which gives (a). For $a=b$, both sides of (a) equal $\frac{\pi a}2$.

(b) The integrand of $I$ is $C^1$ in $(a,b)$, with partial derivatives continuous on $[0,\frac\pi2]\times\left]0,\infty\right[^2$, so Theorem~\ref{thm:paramcompact} (applied in each variable)
gives $\partial_aI=-a\int\cos^2\theta\,R^{-3/2}$ and $\partial_bI=-b\int\sin^2\theta\,R^{-3/2}$, both continuous. Let $P=\int c\,R^{-3/2}$ and $Q=\int s\,R^{-3/2}$. Then
$a^2P+b^2Q=\int R\cdot R^{-3/2}=I$, which is the identity $a\partial_aI+b\partial_bI=-I$, and $P+Q=\int R^{-3/2}=\frac J{a^2b^2}$ by (a). Eliminating $Q$:
$(a^2-b^2)P=I-\frac J{a^2}$. Since $\partial_aI=-aP$, this is the formula.
\end{proof}

\begin{proposition}\label{prop:jstep}
For $a,b>0$, with $a_1=\frac{a+b}2$ and $b_1=\sqrt{ab}$,
\[
2J(a_1,b_1)-J(a,b)=ab\,I(a,b).
\]
\end{proposition}
\begin{proof}
Both sides are continuous in $(a,b)$, so it is enough to treat $a\ne b$. Write $I_1=\partial_1I(a_1,b_1)$ for the partial derivative of $I$ in its first variable at
$(a_1,b_1)$, and $I=I(a,b)=I(a_1,b_1)$. Differentiate the identity $I(a,b)=I\bigl(\frac{a+b}2,\sqrt{ab}\bigr)$ of Theorem~\ref{thm:gaussagm} in $a$, by the chain rule:
\[
\partial_aI(a,b)=\frac12I_1+\frac{b_1}{2a}\,\partial_2I(a_1,b_1).
\]
By Lemma~\ref{lem:jderiv}(b) at $(a_1,b_1)$, $b_1\partial_2I(a_1,b_1)=-I-a_1I_1$, so $\partial_aI(a,b)=\frac{a-a_1}{2a}I_1-\frac I{2a}=\frac{a-b}{4a}I_1-\frac I{2a}$. Put this
into Lemma~\ref{lem:jderiv}(b) at $(a,b)$:
\[
J(a,b)=a^2I-\frac{a^2-b^2}2I+\frac{(a^2-b^2)(a-b)}4I_1.
\]
Lemma~\ref{lem:jderiv}(b) at $(a_1,b_1)$, where $a_1^2-b_1^2=\frac{(a-b)^2}4$, gives $J(a_1,b_1)=a_1^2I+\frac{(a+b)(a-b)^2}8I_1$. So the last term above is
$2\bigl(J(a_1,b_1)-a_1^2I\bigr)$, and
\[
\begin{aligned}J(a,b)&=2J(a_1,b_1)+\Bigl(a^2-\frac{a^2-b^2}2-\frac{(a+b)^2}2\Bigr)I\\&=2J(a_1,b_1)-abI.\end{aligned}
\]
\end{proof}

\begin{theorem}[the perimeter by the AGM]\label{thm:perimeter}
For $a\ge b>0$, with $c_0^2=a^2-b^2$ and $c_n$ as in Proposition~\ref{prop:agm},
\[
J(a,b)=\Bigl(a^2-\sum_{n\ge0}2^{n-1}c_n^2\Bigr)I(a,b),
\]
so the ellipse with semi-axes $a\ge b$ has perimeter
\[
L(a,b)=\frac{2\pi}{M(a,b)}\Bigl(a^2-\sum_{n\ge0}2^{n-1}c_n^2\Bigr).
\]
In terms of the modulus, with the AGM started from $1$ and $k'$, so that $c_0=k$,
\[
E(k)=\Bigl(1-\sum_{n\ge0}2^{n-1}c_n^2\Bigr)K(k).
\]
\end{theorem}
\begin{proof}
Let $I=I(a,b)=I(a_n,b_n)$ for all $n$, and $d_n=J(a_n,b_n)-a_n^2I$. Proposition~\ref{prop:jstep} at $(a_n,b_n)$ reads $J(a_n,b_n)=2J(a_{n+1},b_{n+1})-a_nb_nI$, that is
\[
\begin{aligned}d_n&=2d_{n+1}+\bigl(2a_{n+1}^2-a_nb_n-a_n^2\bigr)I\eqbr=2d_{n+1}-\frac{c_n^2}2I,\end{aligned}
\]
because $2a_{n+1}^2=\frac{a_n^2+2a_nb_n+b_n^2}2$. By induction, $d_0=2^Nd_N-I\sum_{n<N}2^{n-1}c_n^2$. Now $\frac\pi2b_n\le J(a_n,b_n)\le\frac\pi2a_n$ and $\frac\pi2a_n\le a_n^2I\le\frac{\pi a_n^2}{2b_n}$, so $|d_n|\le\frac{\pi(a_n^2-b_n^2)}{2b_n}\le\frac{\pi c_n^2}{2b}$. By Proposition~\ref{prop:agm}(b), $c_{n+1}\le\frac{c_n^2}{4M}$ with $M=M(a,b)$; since $c_n\to0$, there is $N_0$ with $\frac{c_{N_0}}{4M}\le\frac12$, and then
$c_{N_0+j}\le4M\cdot2^{-2^j}$. So $2^Nc_N^2\to0$, $2^Nd_N\to0$, and the series converges: $d_0=-I\sum_{n\ge0}2^{n-1}c_n^2$, which is the formula. The perimeter is
$4J$ and $I=\frac\pi{2M}$. For the last statement take $a=1$, $b=k'$.
\end{proof}

For the ellipse with semi-axes $2$ and $1$, the AGM gives
\[
L=9.68844822054767619842\ldots
\]
after five steps, with every digit correct. For the Earth's orbit,
with eccentricity $0.0167$, the perimeter is $2\pi a\times0.99993027\ldots$: an ellipse that close to a circle is a circle for most purposes.

\begin{figure}[ht]
\centering
\fig{ch12-perimeter}
\caption{The perimeter of the ellipse with semi-axes $1$ and $\sqrt{1-e^2}$, divided by $2\pi$, against the eccentricity $e$ (charcoal): it falls from $1$ (the circle) to
$\frac2\pi$ (the segment $[-1,1]$, travelled twice). Dashed: Euler's bounds $\frac{\pi(a+b)}{2\pi}$ (ochre) and $\frac{\pi\sqrt{2(a^2+b^2)}}{2\pi}$ (coral) of Exercise~\ref{ex:eulerbounds}.}
\label{fig:perimeter}
\end{figure}

\begin{think}
The area of an ellipse is $\pi ab$. Is there a formula of the same kind for its perimeter, built from $a$, $b$, square roots, $\pi$, logarithms and the
like? What would such a formula give when $b\to0$, and when $b=a$?
\end{think}

\begin{remark}
There is none: $E(k)$ is not an elementary function of $k$. This follows from Liouville's theory of integration in finite terms, which decides when an
integral of an algebraic function is elementary; see the survey of Marchisotto and Zakeri in the references. So Theorem~\ref{thm:perimeter}, a limit of
algebraic operations converging quadratically, is as close to a closed form as one can get. Ramanujan found remarkable approximations, such as
$L\approx\pi\bigl(3(a+b)-\sqrt{(3a+b)(a+3b)}\bigr)$, whose relative error is $3\cdot10^{-6}$ for $a=2b$ and at most $0.5\%$ for all ellipses; a numerical calculation tests them
against the AGM.
\end{remark}

\section{Legendre's relation and the computation of $\pi$}

\begin{proposition}\label{prop:kederiv}
On $]0,1[$, $K$ and $E$ are differentiable, with
\[
\begin{gathered}\frac{\mathrm dE}{\mathrm dk}=\frac{E-K}k,\qq\frac{\mathrm dK}{\mathrm dk}=\frac{E-k'^2K}{k\,k'^2}.\end{gathered}
\]
\end{proposition}
\begin{proof}
Let $\Delta=\sqrt{1-k^2\sin^2\theta}$. Differentiating under the integral sign (Theorem~\ref{thm:paramcompact}, on $[0,\frac\pi2]\times[0,k_1]$ for any $k_1<1$),
$\frac{\mathrm dE}{\mathrm dk}=-\int\frac{k\sin^2\theta}\Delta$ and $\frac{\mathrm dK}{\mathrm dk}=\int\frac{k\sin^2\theta}{\Delta^3}$. Since $k^2\sin^2\theta=1-\Delta^2$, the first is
$\frac1k\int\frac{\Delta^2-1}\Delta=\frac{E-K}k$, and the second is $\frac1k\bigl(\int\Delta^{-3}-K\bigr)$. Lemma~\ref{lem:jderiv}(a) with $a=1$, $b=k'$ (so that $R=\Delta^2$ and
$J=E$) gives $\int\Delta^{-3}=\frac E{k'^2}$.
\end{proof}

\begin{theorem}[Legendre's relation, 1811]\label{thm:legendre}
For $0<k<1$, with $K'=K(k')$ and $E'=E(k')$,
\[
EK'+E'K-KK'=\frac\pi2.
\]
\end{theorem}
\begin{proof}
Let $\Lambda(k)=EK'+E'K-KK'$, a function of $k$ alone since $k'=\sqrt{1-k^2}$. By Proposition~\ref{prop:kederiv} and the chain rule with $\frac{\mathrm dk'}{\mathrm dk}=-\frac k{k'}$,
\[
\begin{gathered}\frac{\mathrm dE'}{\mathrm dk}=-\frac{k(E'-K')}{k'^2},\qq\frac{\mathrm dK'}{\mathrm dk}=-\frac{E'-k^2K'}{k\,k'^2}.\end{gathered}
\]
Multiply $\Lambda'=E_kK'+EK'_k+E'_kK+E'K_k-K_kK'-KK'_k$ by $kk'^2$ and substitute. The terms in $EE'$ are $-EE'+E'E=0$. The terms in $EK'$ are
$k'^2+k^2-1=0$ times $EK'$; those in $E'K$ are $(-k^2-k'^2+1)E'K=0$; and those in $KK'$ are $(-k'^2+k^2+k'^2-k^2)KK'=0$. So $\Lambda$ is constant.

To find the constant, let $k\to0^+$. First, $K-E=\int\frac{k^2\sin^2\theta}\Delta\le\frac{\pi k^2}{2k'}$. Second, $K'=\int_0^{\pi/2}\frac{\mathrm d\theta}{\sqrt{\cos^2\theta+k^2\sin^2\theta}}=O\bigl(\ln\frac1k\bigr)$:
with $u=\frac\pi2-\theta$, on $u\le\frac\pi4$ the integrand is at most $\bigl(\frac{4u^2}{\pi^2}+\frac{k^2}2\bigr)^{-1/2}$, whose integral is $\frac\pi2\operatorname{arsinh}\frac1{\sqrt2\,k}$, and on $u\ge\frac\pi4$
it is at most $\sqrt2$. Third, $E'=\int\sqrt{\cos^2\theta+k^2\sin^2\theta}\to\int_0^{\pi/2}\cos\theta\dd\theta=1$, since the integrands differ by at most $k$. Hence
$\Lambda=E'K-(K-E)K'\to1\cdot\frac\pi2-0$.
\end{proof}

\begin{theorem}[Brent--Salamin, 1976]\label{thm:brentsalamin}
Run the AGM from $a_0=1$, $b_0=\frac1{\sqrt2}$, and let $c_n=\frac{a_{n-1}-b_{n-1}}2$ for $n\ge1$. Then
\[
\pi=\frac{4M(1,\frac1{\sqrt2})^2}{1-\sum_{n\ge1}2^{n+1}c_n^2},
\]
and the approximations $\pi_n=\frac{4a_{n+1}^2}{1-\sum_{j=1}^{n}2^{j+1}c_j^2}$ are correct to $1$, $4$, $9$, $20$, $42$, $85$, $173$, $347$ digits for $n=1,\dots,8$.
\end{theorem}
\begin{proof}
At $k=k'=\frac1{\sqrt2}$, Legendre's relation reads $2EK-K^2=\frac\pi2$. By Theorems~\ref{thm:gaussagm} and~\ref{thm:perimeter}, with $M=M(1,\frac1{\sqrt2})$, $K=\frac\pi{2M}$ and
$E=\bigl(1-\sum_{n\ge0}2^{n-1}c_n^2\bigr)K$, where $c_0^2=\frac12$ (and $c_n=\frac{a_{n-1}-b_{n-1}}2$ for $n\ge1$ by Proposition~\ref{prop:agm}(b)). So
\[
\frac\pi2=K^2\Bigl(1-\sum_{n\ge0}2^nc_n^2\Bigr)=\frac{\pi^2}{4M^2}\Bigl(\frac12-\sum_{n\ge1}2^nc_n^2\Bigr),
\]
which rearranges to the formula. The digit counts are computed numerically.
\end{proof}

\begin{remark}
Each step of the algorithm costs a few multiplications and one square root of numbers with $N$ digits. With the fast multiplication of Chapter~9 and
Newton's method for square roots, both cost $O(N\log N\log\log N)$ operations, and about $\log_2N$ steps are needed, so $N$ digits of $\pi$ cost
$O(N(\log N)^2\log\log N)$: far less than the $O(N^2)$ of an arctangent series such as Machin's $\frac\pi4=4\arctan\frac15-\arctan\frac1{239}$ with schoolbook arithmetic. Salamin and Brent
found the algorithm independently in 1976; every ingredient had been known to Gauss and Legendre. Since then, most records have used instead the series
of the Chudnovsky brothers (1988), $\frac1\pi=12\sum_{n\ge0}\frac{(-1)^n(6n)!\,(13591409+545140134n)}{(3n)!\,(n!)^3\,640320^{3n+3/2}}$, which gives $14$ digits per term and is
evaluated by binary splitting with fast multiplication. Its proof uses modular forms, the subject of Volume~13. The three methods can be compared numerically.
\end{remark}

\section{Exercises}

\begin{exercise}[$\star$]\label{ex:agmsteps}
Let $a=1$, $b=\frac1{\sqrt2}$. Using Proposition~\ref{prop:agm}(b), show that $a_n-b_n\le8b\,\bigl(\frac{a-b}{8b}\bigr)^{2^n}$, and deduce that $10$ steps give $M(1,\frac1{\sqrt2})$ to
more than $1000$ digits.
\end{exercise}

\begin{exercise}[$\star$]\label{ex:eulerbounds}
Show that the perimeter of an ellipse with semi-axes $a,b$ satisfies $\pi(a+b)\le L(a,b)\le\pi\sqrt{2(a^2+b^2)}$, with equality only for a circle.
\end{exercise}

\begin{exercise}[$\star\star$]\label{ex:pendulum}
A pendulum of length $\ell$ released at rest from the angle $\theta_0\in\left]0,\pi\right[$ satisfies $\ddot\theta=-\frac g\ell\sin\theta$. Show that its period is
\[
T=4\sqrt{\frac\ell g}\,K\Bigl(\sin\frac{\theta_0}2\Bigr)=\frac{2\pi\sqrt{\ell/g}}{M\bigl(1,\cos\frac{\theta_0}2\bigr)}.
\]
By how much is the period at $\theta_0=\frac\pi2$ longer than the small-oscillation period $2\pi\sqrt{\ell/g}$?
\end{exercise}

\begin{exercise}[$\star\star$]\label{ex:gamma14}
Using Exercise~\ref{ex:lemniscate} of Chapter~7 and Corollary~\ref{cor:lemniscate}, show that
\[
\Gamma\Bigl(\frac14\Bigr)=\frac{(2\pi)^{3/4}}{\sqrt{M(1,\sqrt2)}},
\]
and compute $\Gamma(\frac14)$ to ten digits with five steps of the AGM.
\end{exercise}

\begin{exercise}[$\star\star\star$]\label{ex:agmlog}
Show that, as $x\to0^+$,
\[
\frac\pi{2M(1,x)}=\ln\frac4x+o(1).
\]
So the AGM computes logarithms, with the same quadratic speed.
\end{exercise}

\begin{exercise}[$\star$]\label{pb:cyclewalk}
A token moves on the vertices $0,1,\dots,k-1$ of a regular polygon ($k\ge3$); at each step it moves to one of the two
neighbouring vertices with probability $\frac12$ each, independently of the past. Let $q_n$ be the probability that it is at
$0$ after $n$ steps, having started at $0$.
\begin{enumerate}[label=(\alph*)]
\item Show that $q_n=\frac1k\sum_{j=0}^{k-1}\cos^n\frac{2\pi j}k$.
\item Give $q_n$ in closed form for $k=3$, $4$ and $6$.
\item Find the limit of $q_n$ as $n\to\infty$, according to the parities of $k$ and $n$.
\end{enumerate}
\end{exercise}

\begin{exercise}[$\star\star$]\label{pb:detsquare}
Let $A$ be an $n\times n$ matrix whose entries are independent and uniform on $\{-1,1\}$.
\begin{enumerate}[label=(\alph*)]
\item Show that $\mathbb E[\det A]=0$ and $\mathbb E[(\det A)^2]=n!$.
\item Deduce that some $n\times n$ matrix with entries $\pm1$ has $|\det A|\ge\sqrt{n!}$.
\item Show that every real $n\times n$ matrix with entries in $[-1,1]$ satisfies $|\det A|\le n^{n/2}$ (Hadamard's
inequality), and compare the two bounds as $n\to\infty$.
\end{enumerate}
\end{exercise}

\begin{exercise}[$\star\star\star$]\label{pb:hadamard}
A \emph{Hadamard matrix} of order $n$ is an $n\times n$ matrix $H$ with entries $\pm1$ and $HH^{\mathsf T}=nI_n$.
\begin{enumerate}[label=(\alph*)]
\item Show that if a Hadamard matrix of order $n>2$ exists, then $4\mid n$.
\item Show that if $H$ is a Hadamard matrix of order $n$, then $\begin{pmatrix}H&H\\H&-H\end{pmatrix}$ is one of order $2n$.
\item (Paley, 1933.) Let $p\equiv3\pmod4$ be a prime, $\chi$ the Legendre symbol modulo $p$ (with $\chi(0)=0$), and $Q$ the
$p\times p$ matrix with entries $Q_{ab}=\chi(b-a)$, $a,b\in\mathbb F_p$. Show that $Q^{\mathsf T}=-Q$, that every row of $Q$
sums to $0$, and that $QQ^{\mathsf T}=pI_p-J_p$, where $J_p$ is the all-ones matrix. Deduce a Hadamard matrix of order $p+1$.
\item List the orders $n\le64$ obtained from (b) and (c).
\end{enumerate}
\end{exercise}

\begin{exercise}[$\star\star$]\label{pb:rank}
Let $G$ act transitively on a set $X$ of $n\ge2$ points, let $g$ be uniform in $G$ and $F$ its number of fixed points. The
\emph{rank} $r$ of the action is the number of orbits of $G$ on $X\times X$.
\begin{enumerate}[label=(\alph*)]
\item Show that $\mathbb E[F]=1$ and $\mathbb E[F^2]=r\ge2$.
\item Show that the following are equivalent: $G$ is $2$-transitive (transitive on ordered pairs of distinct points);
$r=2$; $\operatorname{Var}F=1$; the character $\pi-1$ is irreducible, where $\pi(g)=F(g)$ is the permutation character.
\item Find the rank of the dihedral group $D_n$ acting on the $n$ vertices of a regular polygon.
\end{enumerate}
\end{exercise}

\begin{exercise}[$\star\star$]\label{pb:derangements}
Let $G$ act transitively and faithfully on a set $X$ of $n\ge2$ points.
\begin{enumerate}[label=(\alph*)]
\item (Cameron and Cohen.) Show that at least a proportion $\frac1n$ of the elements of $G$ have no fixed point.
\item Show that equality holds if and only if $G$ is \emph{sharply $2$-transitive}: for all pairs $(x,y)$ and $(x',y')$ of
distinct points there is exactly one $g$ with $gx=x'$ and $gy=y'$. Check it for the group of maps $t\mapsto at+b$ of
$\mathbb F_p$, $a\ne0$.
\item (Jordan.) Deduce that a finite group is never the union of the conjugates of a proper subgroup.
\end{enumerate}
\end{exercise}

\begin{exercise}[$\star\star$]\label{pb:noroots}
\emph{Roots of a random polynomial.} Let $f$ be a uniformly random monic polynomial of degree $n\ge1$ over $\mathbb F_q$ and $R$ its
number of roots in $\mathbb F_q$.
\begin{enumerate}[label=(\alph*)]
\item Show that if $n\ge q$, the values $f(a)$, $a\in\mathbb F_q$, are independent and uniform, so that $R$ is exactly binomial
with parameters $q$ and $\frac1q$. In particular $P(R=0)=(1-\frac1q)^q$.
\item For $n<q$, show that
\[
P(R=0)=\sum_{k=0}^n(-1)^k\binom qk q^{-k},
\]
and compare with the probability $\sum_{k=0}^n\frac{(-1)^k}{k!}$ that a random permutation of $n$ points has no fixed point.
\end{enumerate}
\end{exercise}

\begin{exercise}[$\star$]\label{pb:gaussmeasure}
\emph{The Gauss measure.} Let $T(x)=\{1/x\}$ on $]0,1[$ (fractional part), so that $a_{n+1}(x)=\lfloor1/T^nx\rfloor$ is the
$(n+1)$-st partial quotient of $x$, and let $\mu$ be the probability measure with density $\frac1{(1+x)\log2}$ on $[0,1]$.
\begin{enumerate}[label=(\alph*)]
\item Show that $\mu(T^{-1}[0,t[)=\mu([0,t[)$ for every $t\in[0,1]$: $\mu$ is invariant under $T$.
\item Deduce that, if $x$ has law $\mu$, then $P(a_n=k)=\log_2\frac{(k+1)^2}{k(k+2)}$ for every $n\ge1$ and $k\ge1$.
\item If $x$ is uniform instead, compute $P(a_1=k)$, and compare with (b) for $k=1$.
\end{enumerate}
\end{exercise}

\begin{exercise}[$\star\star$]\label{pb:digitsclt}
\emph{Binary digits.} Let $s(n)$ be the number of $1$'s in the binary expansion of $n$, and $n$ uniform in $\{0,\dots,N-1\}$.
\begin{enumerate}[label=(\alph*)]
\item If $N=2^k$, show that $s(n)$ is binomial with parameters $k$ and $\frac12$.
\item For general $N$, show that $\bigl(s(n)-\frac12\log_2N\bigr)/\sqrt{\frac14\log_2N}$ converges in law to a standard Gaussian.
\end{enumerate}
\end{exercise}

\begin{exercise}[$\star\star$]\label{pb:kolberg}
\emph{Both parities, forever (Kolberg, 1959).} Show that $p(n)$ is even for infinitely many $n$ and odd for infinitely many $n$.
\end{exercise}

\begin{exercise}[$\star\star$]\label{pb:kummer}
\emph{Kummer's theorem (1852).} Let $p$ be prime.
\begin{enumerate}[label=(\alph*)]
\item Show that $v_p\binom{m+n}m$ is the number of carries when $m$ and $n$ are added in base $p$.
\item Deduce that every prime power dividing $\binom nk$ is at most $n$.
\item Deduce that $\binom nk$ divides $\operatorname{lcm}(1,2,\dots,n)$, and that $\operatorname{lcm}(1,\dots,n)\ge\frac{2^n}{n+1}$.
\end{enumerate}
\end{exercise}

\begin{exercise}[$\star\star$]\label{pb:pascalfour}
\emph{Modulo $4$.} Let $n\ge1$, and let $B(n)$ be the number of occurrences of the block $10$ in the binary expansion of $n$
(for instance $B(10)=B(1010_2)=2$).
\begin{enumerate}[label=(\alph*)]
\item Show that exactly $B(n)\,2^{s(n)-1}$ entries of row $n$ are $\equiv2\pmod4$.
\item Which rows contain no entry $\equiv2\pmod4$?
\end{enumerate}
\end{exercise}

\begin{exercise}[$\star\star$]\label{pb:centralp}
\emph{Central binomial coefficients modulo $p$.} Let $p$ be an odd prime and $f(x)=\sum_{n\ge0}\binom{2n}nx^n=(1-4x)^{-1/2}$.
\begin{enumerate}[label=(\alph*)]
\item Show that $f(x)\equiv(1-4x)^{(p-1)/2}f(x^p)\pmod p$, coefficient by coefficient.
\item Deduce that $\binom{2n}n\equiv\prod_i\binom{2n_i}{n_i}\pmod p$, where $n=\sum_in_ip^i$ in base $p$.
\item Show that $p\nmid\binom{2n}n$ if and only if every base-$p$ digit of $n$ is at most $\frac{p-1}2$, and that the proportion of
such $n$ below $p^k$ tends to $0$.
\end{enumerate}
\end{exercise}

\begin{exercise}[$\star\star$]\label{pb:sternfib}
\emph{Stern meets Pascal and Fibonacci.} Keep the notation of Problem~\ref{pb:stern}, and let $F_n$ be the Fibonacci numbers.
\begin{enumerate}[label=(\alph*)]
\item Show that $s(n+1)$ is the number of $k$ such that $\binom{n-k}k$ is odd.
\item Show that $s(n)\le F_{k+2}$ for $2^k\le n\le2^{k+1}$.
\item Show that equality holds for $n=\frac13\bigl(2^{k+2}-(-1)^k\bigr)$.
\end{enumerate}
\end{exercise}

\begin{exercise}[$\star\star$]\label{pb:thuemorse}
\emph{The Thue--Morse series is algebraic.} Let $t_n=s(n)\bmod2\in\mathbb F_2$ (here $s(n)$ is again the binary digit sum) and
$T(x)=\sum_nt_nx^n\in\mathbb F_2[[x]]$.
\begin{enumerate}[label=(\alph*)]
\item Show that $(1+x)^3T^2+(1+x)^2T+x=0$ in $\mathbb F_2[[x]]$.
\item Show that $(t_n)$ is not eventually periodic.
\item Deduce that $T$ is algebraic of degree $2$ over $\mathbb F_2(x)$ but not rational.
\end{enumerate}
\end{exercise}

\begin{exercise}[$\star\star$]\label{pb:hermite}
\emph{Powers of $e$.}
\begin{enumerate}[label=(\alph*)]
\item Show that $e^k$ is irrational for every integer $k\ge1$.
\item Deduce that $e^r$ is irrational for every rational $r\ne0$, and that $\log r$ is irrational for every rational $r>0$, $r\ne1$.
\end{enumerate}
\end{exercise}

\begin{exercise}[$\star\star$]\label{pb:centralsum}
\emph{Reciprocals of central binomial coefficients.}
\begin{enumerate}[label=(\alph*)]
\item Show that $\displaystyle\frac1{n\binom{2n}n}=\int_0^1t^{n-1}(1-t)^n\dd t$ for $n\ge1$, and deduce that $\displaystyle\sum_{n\ge1}\frac1{n\binom{2n}n}=\frac\pi{3\sqrt3}$.
\item Show that $\displaystyle\sum_{n\ge0}\frac1{\binom{2n}n}=\frac43+\frac{2\pi\sqrt3}{27}$.
\end{enumerate}
\end{exercise}

\begin{exercise}[$\star$]\label{pb:lehmer}
\emph{Series with $\binom{2n}n/4^n$.} Let $c_n=\binom{2n}n4^{-n}$.
\begin{enumerate}[label=(\alph*)]
\item Show that $\displaystyle\sum_{n\ge0}\frac{c_n}{2n+1}=\frac\pi2$ and $\displaystyle\sum_{n\ge1}\frac{c_n}n=2\log2$.
\item Show that $\sum_nc_n^2$ diverges, while $\displaystyle\sum_{n\ge0}\frac{c_n^2}{n+1}=\frac4\pi$.
\end{enumerate}
\end{exercise}

\begin{exercise}[$\star\star$]\label{pb:catalanmoments}
\emph{Catalan numbers are moments.} Let $C_n=\frac1{n+1}\binom{2n}n$.
\begin{enumerate}[label=(\alph*)]
\item Show that $\displaystyle\frac1{2\pi}\int_{-2}^2x^{2n}\sqrt{4-x^2}\dd x=C_n$, while the odd moments vanish.
\item Show that $\displaystyle\frac1{2\pi}\int_0^4x^n\sqrt{\frac{4-x}x}\dd x=C_n$.
\item From (b), show that $C_n\sim\frac{4^n}{\sqrt\pi\,n^{3/2}}$.
\end{enumerate}
\end{exercise}

\begin{exercise}[$\star\star$]\label{pb:apcount}
\emph{Counting progressions with Fourier.} Let $p$ be an odd prime, $A\subseteq\mathbb F_p$, and let $T(A)$ be the number of pairs
$(x,d)\in\mathbb F_p^2$ with $x$, $x+d$, $x+2d$ all in $A$ (the pairs with $d=0$ are included).
\begin{enumerate}[label=(\alph*)]
\item Show that $T(A)=\frac1p\sum_{r\in\mathbb F_p}\hat A(r)^2\hat A(-2r)$.
\item Deduce that $\bigl|T(A)-\frac{|A|^3}p\bigr|\le|A|\max_{r\ne0}|\hat A(r)|$.
\item Let $Q$ be the set of non-zero squares modulo $p$. Show that $\bigl|T(Q)-\frac{(p-1)^3}{8p}\bigr|\le\frac{(\sqrt p+1)(p-1)}4$; so $T(Q)=\frac{p^2}8+O(p^{3/2})$.
\item Let $|A|=\delta p$. Show that if $|\hat A(r)|\le\frac12\delta^2p$ for all $r\ne0$ and $\delta^2p>2$, then $A$ contains a
progression $x,x+d,x+2d$ with $d\ne0$.
\end{enumerate}
\end{exercise}

\begin{exercise}[$\star\star\star$]\label{pb:dftspectrum}
\emph{The eigenvalues of the discrete Fourier transform.} Let $n\ge2$ and let $F$ be the unitary DFT on $\C^n$:
$(Fx)_j=\frac1{\sqrt n}\sum_{k=0}^{n-1}x_ke(-\frac{jk}n)$, indices modulo $n$.
\begin{enumerate}[label=(\alph*)]
\item Show that $F^2$ is the reversal $x_k\mapsto x_{-k}$, that $F^4=I$, and that $F$ is diagonalisable with eigenvalues among
$1,-1,-i,i$. Call their multiplicities $m_1$, $m_{-1}$, $m_{-i}$, $m_i$.
\item Express $\operatorname{tr}F$ and $\operatorname{tr}F^2$ in terms of the multiplicities, and compute both traces.
\item Deduce the multiplicities: writing $n=4m+s$ with $0\le s\le3$, the vector $(m_1,m_{-1},m_{-i},m_i)$ is
\[
\begin{array}{c|c}
s&(m_1,m_{-1},m_{-i},m_i)\\\hline
0&(m+1,m,m,m-1)\\
1&(m+1,m,m,m)\\
2&(m+1,m+1,m,m)\\
3&(m+1,m+1,m+1,m)
\end{array}
\]
for $s=0,1,2,3$.
\end{enumerate}
\end{exercise}

\begin{exercise}[$\star\star\star$]\label{pb:erdosturan}
\emph{An Erdős--Turán inequality.} Let $x_1,\dots,x_N\in[0,1[$, $S(h)=\sum_{n=1}^Ne(hx_n)$, and let $D$ be the \emph{discrepancy}, the
supremum over intervals $I\subseteq[0,1[$ of $\bigl|\frac1N\#\{n:x_n\in I\}-|I|\bigr|$.
\begin{enumerate}[label=(\alph*)]
\item Let $0<\delta<1$ and let $J$ be an interval of length $L\le1-\delta$. The $1$-periodic \emph{trapezoid} $t$ is the convolution
of the periodic indicator of $J$ with $\frac1\delta1_{[0,\delta]}$. Show that $0\le t\le1$, $\int_0^1t=L$, and that its Fourier
coefficients satisfy $|\hat t(h)|\le\min\bigl(\frac1{\pi|h|},\frac1{\pi^2h^2\delta}\bigr)$ for $h\ne0$.
\item Deduce that
\[
D\le2\delta+\frac2N\sum_{h\ge1}\min\Bigl(\frac1{\pi h},\frac1{\pi^2h^2\delta}\Bigr)|S(h)|.
\]
\item Let $p$ be an odd prime and $x_k=\{\frac{k^2}p\}$ for $0\le k<p$. Show that the discrepancy of these $p$ points is at most
$\frac{\log p+10}{\pi\sqrt p}$.
\end{enumerate}
\end{exercise}

\begin{exercise}[$\star\star$]\label{pb:circlewalk}
\emph{A random walk on the circle.} Let $\alpha$ be real, $\varepsilon_1,\varepsilon_2,\dots$ independent fair signs, and
$X_n=\{\alpha(\varepsilon_1+\cdots+\varepsilon_n)\}\in[0,1[$.
\begin{enumerate}[label=(\alph*)]
\item Show that if $\alpha$ is irrational, $X_n$ converges in law to the uniform distribution: $\mathbb E[g(X_n)]\to\int_0^1g$ for every
continuous $1$-periodic $g$.
\item Let $\alpha=\frac aq$ in lowest terms, $q\ge2$. Show that if $q$ is odd, $X_n$ converges in law to the uniform distribution on
$\{0,\frac1q,\dots,\frac{q-1}q\}$, and that if $q$ is even, $X_{2n}$ and $X_{2n+1}$ converge in law to the uniform distributions on two
disjoint sets of $\frac q2$ points, so that $X_n$ does not converge.
\end{enumerate}
\end{exercise}

\begin{exercise}[$\star\star\star$]\label{pb:fejer}
\emph{Non-negative trigonometric polynomials.} Let $T(\theta)=\sum_{|k|\le n}c_ke^{ik\theta}$ be real and $T\ge0$.
\begin{enumerate}[label=(\alph*)]
\item (Fejér--Riesz, 1916.) Show that $T(\theta)=|P(e^{i\theta})|^2$ for a polynomial $P\in\C[z]$ of degree at most $n$.
\item (Fejér.) Suppose $c_0=1$. Show that $T(0)\le n+1$, with equality only for the Fejér kernel $\frac1{n+1}\bigl|\sum_{k=0}^ne^{ik\theta}\bigr|^2$.
\item (Egerváry--Szász.) Suppose $T(\theta)=1+\sum_{k=1}^na_k\cos k\theta$ with real $a_k$. Show that $a_1\le2\cos\frac\pi{n+2}$, and that equality is
attained.
\end{enumerate}
\end{exercise}

\begin{exercise}[$\star\star\star$]\label{pb:bernstein}
\emph{Bernstein's inequality.} Let $T$ be a real trigonometric polynomial of degree at most $n$ and $\|T\|=\max|T|$.
\begin{enumerate}[label=(\alph*)]
\item Show that if $\|T\|\le1$, then $T'(\theta)^2+n^2T(\theta)^2\le n^2$ for all $\theta$. Deduce $\|T'\|\le n\|T\|$, and that the same holds for
complex-valued $T$.
\item Deduce that a real polynomial $P$ of degree $n$ with $|P|\le1$ on $[-1,1]$ satisfies $|P'(x)|\le\frac n{\sqrt{1-x^2}}$ for $|x|<1$.
\item Show that both inequalities are sharp.
\end{enumerate}
\end{exercise}

\begin{exercise}[$\star\star\star$]\label{pb:bernsteinabs}
\emph{Absolutely convergent Fourier series.} Let $f$ be $2\pi$-periodic with $|f(x)-f(y)|\le C|x-y|^\alpha$ for all $x,y$, where
$0<\alpha\le1$.
\begin{enumerate}[label=(\alph*)]
\item Show that for every $h$, $\frac1{2\pi}\int_{-\pi}^\pi|f(x+h)-f(x-h)|^2\dd x=\sum_k4\sin^2(kh)\,|c_k(f)|^2$.
\item (S.~Bernstein, 1914.) Suppose $\alpha>\frac12$. Show that $\sum_k|c_k(f)|<\infty$, and that the Fourier series of $f$ converges
uniformly to $f$.
\item Let $P_j$ be the Rudin--Shapiro polynomials of Problem~\ref{pb:rudin}, and $f(x)=\sum_{j\ge0}2^{-j}e^{i2^jx}P_j(e^{ix})$. Show that $f$
satisfies the hypothesis with $\alpha=\frac12$, but that $\sum_k|c_k(f)|=\infty$. So $\frac12$ is the threshold.
\end{enumerate}
\end{exercise}

\begin{exercise}[$\star\star$]\label{pb:fantaussky}
\emph{Polygons by the discrete Fourier transform.} Let $n\ge3$ and $\omega=e(\frac1n)$; indices are taken modulo $n$.
\begin{enumerate}[label=(\alph*)]
\item (Fan, Taussky and Todd, 1955.) Let $x_0,\dots,x_{n-1}\in\C$ with $\sum_kx_k=0$. Show that
\[
\sum_k|x_{k+1}-x_k|^2\ge4\sin^2\frac\pi n\sum_k|x_k|^2,
\]
with equality if and only if $x_k=a\omega^k+b\omega^{-k}$ for some $a,b\in\C$.
\item Let $z_0,\dots,z_{n-1}$ be the vertices of a closed polygon in $\C$, with signed area $A=\frac12\operatorname{Im}\sum_k\overline{z_k}z_{k+1}$. Show that
\[
\sum_k|z_{k+1}-z_k|^2\ge4\tan\frac\pi n\cdot A,
\]
with equality if and only if the polygon is a positively oriented regular $n$-gon (or a single point).
\item Deduce that among equilateral $n$-gons of given perimeter, the regular one has the largest area.
\end{enumerate}
\end{exercise}

\begin{exercise}[$\star\star\star$]\label{pb:involasym}
\emph{How many involutions?} Let $t_n$ be the number of involutions of $\{1,\dots,n\}$, so that $\sum_nt_n\frac{z^n}{n!}=T(z)=e^{z+z^2/2}$ (Volume~5,
Exercise~2.4). Let $r=r_n>0$ be the root of $r+r^2=n$, and $b=r+2r^2$.
\begin{enumerate}[label=(\alph*)]
\item Show that $\frac{t_n}{n!}=\frac1{2\pi}\int_{-\pi}^\pi T(re^{i\theta})\,r^{-n}e^{-in\theta}\dd\theta$.
\item Show that $|T(re^{i\theta})|\le T(r)\exp\bigl(-2r\sin^2\frac\theta2-r^2\sin^2\theta\bigr)$.
\item Show that
\[
\frac{t_n}{n!}\sim\frac{T(r)}{r^n\sqrt{2\pi b}}.
\]
(Cut the integral at $|\theta|=r^{-4/5}$.)
\item Deduce that $t_n\sim\frac1{\sqrt2}\,n^{n/2}e^{-n/2+\sqrt n-1/4}$.
\end{enumerate}
\end{exercise}

\begin{exercise}[$\star\star\star$]\label{pb:distinctcycles}
\emph{All cycles of different lengths.} Let $d_n$ be the probability that a uniform random permutation of $\{1,\dots,n\}$ has cycles of pairwise distinct lengths.
\begin{enumerate}[label=(\alph*)]
\item Show that $\sum_{n\ge0}d_nz^n=\prod_{k\ge1}\bigl(1+\frac{z^k}k\bigr)$ for $|z|<1$, and compute $d_1,\dots,d_5$.
\item Show that this product equals $\frac{h(z)}{1-z}$, where $h(z)=\prod_{k\ge1}\bigl(1+\frac{z^k}k\bigr)e^{-z^k/k}$, and that $h(1)=e^{-\gamma}$.
\item For a power series $f=\sum a_jz^j$ write $\|f\|=\sum_j|a_j|$. Show that $\|fg\|\le\|f\|\,\|g\|$ and that $\bigl\|(1+\frac{z^k}k)e^{-z^k/k}-1\bigr\|\le\frac1{k^2}$.
Deduce that the coefficients $h_j$ of $h$ satisfy $\sum_j|h_j|<\infty$ and $\sum_jh_j=h(1)$.
\item Deduce that $d_n\to e^{-\gamma}=0.5614\ldots$
\end{enumerate}
\end{exercise}

\begin{exercise}[$\star\star\star$]\label{pb:ramanujan}
\emph{Ramanujan's $\frac13$.} Let $Q(n)$ be as in Problem~\ref{pb:poissonise}, let $R(n)=1+\frac n{n+1}+\frac{n^2}{(n+1)(n+2)}+\cdots$, and define $\theta_n$ by
$\sum_{k=0}^{n-1}\frac{n^k}{k!}+\theta_n\frac{n^n}{n!}=\frac{e^n}2$.
\begin{enumerate}[label=(\alph*)]
\item Show that $\frac{n^n}{n!}Q(n)=\sum_{k<n}\frac{n^k}{k!}$ and $\frac{n^n}{n!}R(n)=\sum_{k\ge n}\frac{n^k}{k!}$. Deduce that $Q(n)+R(n)=\frac{n!\,e^n}{n^n}$ and $\theta_n=\frac{R(n)-Q(n)}2$.
\item Show that
\[
Q(n)=n\int_1^\infty s^{n-1}e^{n(1-s)}\dd s,\qq R(n)=n\int_0^1s^{n-1}e^{n(1-s)}\dd s.
\]
\item Let $\varphi(s)=s-1-\ln s$ and $v(s)=\operatorname{sign}(s-1)\sqrt{2\varphi(s)}$, an increasing bijection from $\left]0,\infty\right[$ onto $\R$; write $s(v)$ for its inverse.
Show that
\[
Q(n)=n\int_0^\infty e^{-nv^2/2}g(v)\dd v,\qq R(n)=n\int_{-\infty}^0e^{-nv^2/2}g(v)\dd v,\qq g(v)=\frac v{s(v)-1},
\]
and that $0<g(v)\le1+|v|$ and $g(v)=1-\frac v3+O(v^2)$, uniformly on $\R$.
\item Deduce that $Q(n)=\sqrt{\frac{\pi n}2}-\frac13+O\bigl(\frac1{\sqrt n}\bigr)$ and $R(n)=\sqrt{\frac{\pi n}2}+\frac13+O\bigl(\frac1{\sqrt n}\bigr)$, so that $\theta_n\to\frac13$.
\item Deduce that $e^{-n}\sum_{k=0}^n\frac{n^k}{k!}=\frac12+\frac2{3\sqrt{2\pi n}}+O\bigl(\frac1n\bigr)$.
\end{enumerate}
\end{exercise}

\begin{exercise}[$\star\star$]\label{pb:mappings}
\emph{Random maps.} Let $f$ be a uniform random map from $\{1,\dots,n\}$ to itself. Its \emph{functional graph} has an arrow $x\to f(x)$ for each $x$;
each connected component contains exactly one cycle, and $x$ is \emph{cyclic} if $f^k(x)=x$ for some $k\ge1$.
\begin{enumerate}[label=(\alph*)]
\item Show that $x$ lies on a cycle of length exactly $k$ with probability $\frac{n^{\underline k}}{n^{k+1}}$. Deduce that the expected number of cyclic points is
$Q(n)\sim\sqrt{\pi n/2}$.
\item Show that the expected number of components is $C_n=\sum_{k=1}^n\frac{n^{\underline k}}{k\,n^k}$.
\item Show that $C_n=\frac12\bigl(\ln(2n)+\gamma\bigr)+o(1)$. (For $n=100$, $C_n=2.979$ and the approximation is $2.938$.)
\end{enumerate}
\end{exercise}

\begin{exercise}[$\star\star$]\label{pb:xplusinv}
\emph{Adding the reciprocal.} Let $x_0=1$ and $x_{n+1}=x_n+\frac1{x_n}$.
\begin{enumerate}[label=(\alph*)]
\item Show that $x_n^2\ge2n+1$, and that $x_n^2\le2n+2+\ln n$ for $n\ge1$.
\item Show that $x_n^2=2n+\frac12\ln n+C+o(1)$ for a constant $C$. (Numerically, one finds $C=1.72314\ldots$)
\item Deduce that $x_n=\sqrt{2n}+\frac{\ln n+2C}{4\sqrt{2n}}+o\bigl(\frac1{\sqrt n}\bigr)$.
\end{enumerate}
\end{exercise}

\begin{exercise}[$\star\star\star$]\label{pb:hyperfactorial}
\emph{The hyperfactorial and Glaisher's constant.} Let $H(n)=1^1\,2^2\,3^3\cdots n^n$.
\begin{enumerate}[label=(\alph*)]
\item Differentiate the identity of Volume~5, Proposition~13.7 with respect to $s$ and set $s=-1$, to show that
\[
\ln H(n)=\Bigl(\frac{n^2}2+\frac n2+\frac1{12}\Bigr)\ln n-\frac{n^2}4+\ln A+\frac1{720n^2}+O\Bigl(\frac1{n^4}\Bigr),
\]
where $\ln A=\frac1{12}-\zeta'(-1)$.
\item Direct computation gives $A=1.2824271291\ldots$, Glaisher's constant. Check (a) numerically for $n=10$.
\end{enumerate}
\end{exercise}

\begin{exercise}[$\star\star$]\label{pb:banach}
\emph{Banach's matchboxes.} A smoker carries two boxes of $N$ matches, one in each pocket, and each time takes a match from a pocket chosen at random. At the moment he first
finds a box empty, let $R$ be the number of matches in the other.
\begin{enumerate}[label=(\alph*)]
\item Show that $P(R=r)=\binom{2N-r}N2^{-(2N-r)}$ for $0\le r\le N$.
\item Writing $a_r=P(R=r)$, show that $(2N-r)a_{r+1}=2(N-r)a_r$, and deduce that $\mathbb E[R]=(2N+1)\binom{2N}N4^{-N}-1$.
\item Deduce that $\mathbb E[R]=2\sqrt{N/\pi}-1+O\bigl(\frac1{\sqrt N}\bigr)$. With $50$ matches per box, how many are left on average?
\end{enumerate}
\end{exercise}

\begin{exercise}[$\star\star$]\label{pb:primezeta}
\emph{Computing $\sum_p\frac1{p^2}$ to twelve digits.} For real $s>1$ let $P(s)=\sum_pp^{-s}$.
\begin{enumerate}[label=(\alph*)]
\item Show that $\log\zeta(s)=\sum_{k\ge1}\frac{P(ks)}k$.
\item Show that $P(s)=\sum_{k\ge1}\frac{\mu(k)}k\log\zeta(ks)$, the double series converging absolutely.
\item Show that if the sum in (b) is stopped at $k=K$, the error is at most $\frac{4}{K}2^{-(K+1)s}$ for $s\ge2$.
\item Summing the primes up to $10^7$ gives $\sum_{p\le10^7}p^{-2}=0.4522474141\ldots$, wrong in the eighth decimal. With (b) and $K=20$, one finds numerically
$P(2)=0.452247420041\ldots$, correct to twelve digits. Estimate the tail $\sum_{p>10^7}p^{-2}$ and check the two numbers against each other.
\end{enumerate}
\end{exercise}

\begin{exercise}[$\star\star\star$]\label{pb:dickmanint}
\emph{Dickman's function integrates to $e^\gamma$.}
\begin{enumerate}[label=(\alph*)]
\item Show that $u\rho(u)=\int_{u-1}^u\rho(t)\dd t$ for $u\ge1$. Deduce that $\rho>0$, that $\rho$ decreases on $[1,\infty[$, and that $\rho(u)\le\frac12\rho(u-1)$ for $u\ge2$.
\item Let $\hat\rho(s)=\int_0^\infty e^{-su}\rho(u)\dd u$ for $s\ge0$. Extending $\rho$ by $0$ on $]-\infty,0[$, show that $(u\rho(u))'=\rho(u)-\rho(u-1)$ for $u>0$, $u\ne1$, and deduce
that $s\hat\rho'(s)=-(1-e^{-s})\hat\rho(s)$ for $s>0$.
\item Deduce that $\hat\rho(s)=\hat\rho(0)\exp\bigl(-\int_0^s\frac{1-e^{-t}}t\dd t\bigr)$.
\item Show that $s\hat\rho(s)\to1$ as $s\to\infty$, and that $\int_0^s\frac{1-e^{-t}}t\dd t=\log s+\gamma+o(1)$. Conclude that $\int_0^\infty\rho(u)\dd u=e^\gamma=1.78107\ldots$
\item By Mertens's third theorem, $\sum\frac1n$ over the $y$-smooth $n$ equals $\prod_{p\le y}(1-\frac1p)^{-1}\sim e^\gamma\log y$. Assuming that $\Psi(t,y)\approx t\rho(\frac{\log t}{\log y})$,
show that this is exactly what (d) predicts.
\end{enumerate}
\end{exercise}

\begin{exercise}[$\star\star$]\label{pb:polymertens}
\emph{Mertens's theorem for polynomials.} Let $N_k$ be the number of monic irreducible polynomials of degree $k$ over $\mathbb F_q$, so that
$\sum_{k\mid m}kN_k=q^m$ (Volume~3, Chapter~9). Let $E_n=\prod_{k\le n}(1-q^{-k})^{N_k}$.
\begin{enumerate}[label=(\alph*)]
\item Show that $-\log E_n=H_n+\varepsilon_n$, where $H_n=\sum_{m\le n}\frac1m$ and $0\le\varepsilon_n\le4(1+\log n)q^{-n/2}$. Deduce that $E_n\sim\frac{e^{-\gamma}}n$.
\item A monic polynomial of degree $n$ with no irreducible factor of degree at most $\frac n2$ is irreducible, and the probability of that is $\frac1n+O(\frac{q^{-n/2}}n)$
(Problem~\ref{pb:randpoly}). Independent sieving would predict $E_{\lfloor n/2\rfloor}$. Show that the prediction exceeds the truth by a factor tending to $2e^{-\gamma}=1.1229\ldots$, the same factor as for the integers.
\item Show that $ne^\gamma E_n=1-\frac1{2n}+O(\frac1{n^2})$ for fixed $q$. Numerically, one finds $ne^\gamma E_n=0.952015$ for $q=5$, $n=10$. Compare.
\end{enumerate}
\end{exercise}

\begin{figure}[ht]
\centering
\fig{ch07-pascal3}
\caption{Pascal's triangle modulo $3$, rows $0$ to $26$: charcoal dots are entries $\equiv1$, rust dots entries $\equiv2$, and the gaps
are multiples of $3$. Row $n$ has $2^{a_1}3^{a_2}$ dots, where $a_j$ counts the digits $j$ of $n$ in base $3$, and the charcoal dots
always win, by exactly $2^{a_1}$.}
\label{fig:pascal3}
\end{figure}

\section{Hints}

\begin{hint}{ex:agmsteps}
Put $e_n=\frac{a_n-b_n}{8b}$; then $e_{n+1}\le e_n^2$.
\end{hint}
\begin{hint}{ex:eulerbounds}
For the lower bound, compare
\[
\sqrt{a^2\cos^2\theta+b^2\sin^2\theta}
\quad\text{with}\quad a\cos^2\theta+b\sin^2\theta.
\]
For the upper bound use the Cauchy--Schwarz inequality.
\end{hint}
\begin{hint}{ex:pendulum}
Conservation of energy gives $\dot\theta^2$ as a function of $\theta$. Then substitute $\sin\frac\theta2=\sin\frac{\theta_0}2\sin\varphi$.
\end{hint}
\begin{hint}{ex:gamma14}
Both results compute the same integral $\int_0^1\frac{\mathrm dr}{\sqrt{1-r^4}}$.
\end{hint}
\begin{hint}{ex:agmlog}
Use the algebraic form $I(1,x)=\int_0^\infty\frac{\mathrm dt}{\sqrt{(1+t^2)(x^2+t^2)}}$ and split the integral at $t=1$.
\end{hint}

\begin{hint}{pb:cyclewalk}
The indicator of $S\equiv0\pmod k$ is $\frac1k\sum_j\omega^{jS}$ with $\omega=e^{2i\pi/k}$.
\end{hint}

\begin{hint}{pb:detsquare}
Expand $\det A$ as a sum over permutations, then expand the square. For (c), Gram--Schmidt.
\end{hint}

\begin{hint}{pb:hadamard}
(a) Normalise the first row to all $1$s and count the columns according to the signs in rows $2$ and $3$. (c) For $d\ne0$,
$\sum_c\chi(c)\chi(c+d)=-1$. Put $S=\begin{pmatrix}0&\mathbf1^{\mathsf T}\\-\mathbf1&Q\end{pmatrix}$ and $H=I+S$.
\end{hint}

\begin{hint}{pb:rank}
Fixed points of $g$ on $X\times X$ are pairs of fixed points. For the character, $\langle\pi,\pi\rangle=\mathbb E[F^2]$.
\end{hint}

\begin{hint}{pb:derangements}
Look at $\mathbb E[(F-1)(F-n)]$ and at the sign of $(F-1)(F-n)$ when $F\ge1$.
\end{hint}

\begin{hint}{pb:noroots}
(a) Evaluation at all points of $\mathbb F_q$ is a surjective linear map on polynomials of degree $<q$. (b) Inclusion--exclusion.
\end{hint}

\begin{hint}{pb:gaussmeasure}
(a) $T^{-1}[0,t[=\bigcup_k\,]\frac1{k+t},\frac1k]$, and the series telescopes.
\end{hint}

\begin{hint}{pb:digitsclt}
(b) Split $\{0,\dots,N-1\}$ according to the binary expansion of $N$ into dyadic blocks, on each of which $s(n)$ is a shifted
binomial.
\end{hint}

\begin{hint}{pb:kolberg}
Otherwise $\sum p(n)x^n$ would be rational modulo $2$, and so would $E(x)$. What does a rational series over $\mathbb F_2$ look
like?
\end{hint}

\begin{hint}{pb:kummer}
(a) Legendre: $v_p(n!)=\frac{n-s_p(n)}{p-1}$. (b) A carry out of position $i$ needs a digit of $n$ in position $i+1$.
\end{hint}

\begin{hint}{pb:pascalfour}
By Kummer, $\binom nk\equiv2\pmod4$ means exactly one carry in $k+(n-k)$. Where can it happen?
\end{hint}

\begin{hint}{pb:centralp}
(a) $f^2(1-4x)=1$ and $f(x)^p\equiv f(x^p)$. (b) $\binom{-1/2}j(-4)^j=\binom{2j}j$.
\end{hint}

\begin{hint}{pb:sternfib}
(a) $\binom{n-k}k$ is odd if and only if $k$ and $n-2k$ have no binary digit $1$ in common. (b) Of two consecutive integers, one is
even.
\end{hint}

\begin{hint}{pb:thuemorse}
(a) $t_{2n}=t_n$, $t_{2n+1}=1+t_n$, and $T(x)^2=T(x^2)$. (b) A period $P$ could be halved if even; if it is odd, compare
$t_{2n+P}$ and $t_{2n+1+P}$.
\end{hint}

\begin{hint}{pb:hermite}
Imitate Niven's proof for $\pi^2$ with $f(x)=\frac{x^n(1-x)^n}{n!}$ and $F=\sum_j(-1)^jk^{2n-j}f^{(j)}$: compute $\frac{\mathrm d}{\mathrm dx}\bigl(e^{kx}F(x)\bigr)$.
\end{hint}

\begin{hint}{pb:centralsum}
(a) The Beta integral. (b) $\frac1{\binom{2n}n}=(2n+1)\int_0^1(t-t^2)^n\dd t$ and $\sum_n(2n+1)u^n=\frac{1+u}{(1-u)^2}$.
\end{hint}

\begin{hint}{pb:lehmer}
(a) $\sum c_nx^{2n}=\frac1{\sqrt{1-x^2}}$. (b) $c_{n+1}=\frac{2n+1}{2n+2}c_n$; look for $b_n=\lambda nc_n^2$ with $b_n-b_{n+1}=\frac{c_n^2}{n+1}$.
\end{hint}

\begin{hint}{pb:catalanmoments}
(a) $x=2\sin\theta$ and the Beta integral. (b) $x=y^2$. (c) Laplace's method near $x=4$.
\end{hint}

\begin{hint}{pb:apcount}
(a) A triple $(x,y,z)$ is a progression iff $x+z=2y$. (b) Isolate $r=0$ and use Parseval for the rest. (c) Write $1_Q(x)=\frac12(1+\chi(x))$
for $x\ne0$, with $\chi$ the Legendre symbol, and use the Gauss sum. (d) Progressions with $d=0$ number $|A|$.
\end{hint}

\begin{hint}{pb:dftspectrum}
(a) $(F^2x)_j=\frac1n\sum_kx_k\sum_le(-\frac{l(j+k)}n)$. A matrix killed by a polynomial with simple roots is diagonalisable. (b) $\operatorname{tr}F$ is a
quadratic Gauss sum, whose value Gauss determined (Volume~5, Theorem~10.14). (c) Four linear equations in four unknowns.
\end{hint}

\begin{hint}{pb:erdosturan}
(a) The Fourier coefficients of a convolution multiply. (b) Choose $J$ so that $t\ge1_I$ for an upper bound, and $t\le1_I$ for a lower
bound; expand $t$ in its (absolutely convergent) Fourier series and sum over the points. (c) Gauss sums have modulus $\sqrt p$;
take $\delta=p^{-1/2}$.
\end{hint}

\begin{hint}{pb:circlewalk}
$\mathbb E[e(hX_n)]=\cos^n(2\pi h\alpha)$. When is $|\cos(2\pi h\alpha)|=1$?
\end{hint}

\begin{hint}{pb:fejer}
(a) Look at the roots of the polynomial $z^nT(z)$: they come in pairs $\alpha$, $1/\bar\alpha$, and those on the unit circle have even
multiplicity. (b) Cauchy--Schwarz on $P(1)$. (c) $a_1=2\operatorname{Re}\sum p_kp_{k+1}$ is a quadratic form; the adjacency matrix of a path.
\end{hint}

\begin{hint}{pb:bernstein}
(a) If the inequality fails at $\theta_0$, choose $\beta$ with $T(\theta_0)=\cos(n(\theta_0-\beta))$ and count the zeros of $\cos(n(\theta-\beta))-T(\theta)$. (b) $T(\theta)=P(\cos\theta)$.
(c) $\sin n\theta$ and the Chebyshev polynomials.
\end{hint}

\begin{hint}{pb:bernsteinabs}
(a) Parseval for $f(x+h)-f(x-h)$. (b) Dyadic blocks $2^{j-1}<|k|\le2^j$ with $h=\frac\pi{3\cdot2^j}$, then Cauchy--Schwarz in each block. (c) Split
$f$ at the scale of $h$: Bernstein's inequality for the low blocks, the sup norm for the high ones.
\end{hint}

\begin{hint}{pb:fantaussky}
(a) Expand $x_k=\sum_mc_m\omega^{mk}$ and use Parseval for the DFT. (b) Do the same for the area. (c) For an equilateral polygon,
$L^2=n\sum|z_{k+1}-z_k|^2$.
\end{hint}

\begin{hint}{pb:involasym}
(a) Cauchy's formula on the circle $|z|=r$. (b) Compute $\operatorname{Re}(re^{i\theta}+\frac12r^2e^{2i\theta})$. (c) Expand $re^{i\theta}+\frac12r^2e^{2i\theta}-in\theta$ to order $2$ in $\theta$: the
linear term vanishes because $r+r^2=n$. (d) $r=\sqrt n-\frac12+\frac1{8\sqrt n}+O(n^{-3/2})$, $r+\frac{r^2}2=\frac{n+r}2$, and Stirling for $n!$.
\end{hint}

\begin{hint}{pb:distinctcycles}
(a) The exponential formula (Volume~5, Theorem~2.9), each cycle length being used at most once. (b) $\prod_ke^{z^k/k}=\frac1{1-z}$, and $\prod_{k\le N}(1+\frac1k)=N+1$.
(c) $(1+w)e^{-w}=1-\sum_{j\ge2}(-1)^j\frac{j-1}{j!}w^j$. The partial products form a Cauchy sequence for $\|\cdot\|$. (d) $d_n=h_0+\dots+h_n$.
\end{hint}

\begin{hint}{pb:ramanujan}
(a) $\frac{n^{n-j}}{(n-j)!}=\frac{n^n}{n!}\cdot\frac{n^{\underline j}}{n^j}$. (b) $\sum_{k\ge n}e^{-x}\frac{x^k}{k!}=\int_0^x\frac{t^{n-1}e^{-t}}{(n-1)!}\dd t$ (differentiate both sides), then put $x=n$ and $t=ns$. (c) $s^{n-1}e^{n(1-s)}=\frac{e^{-n\varphi(s)}}s$
and $\mathrm d\varphi=\frac{s-1}s\mathrm ds=v\,\mathrm dv$. For the bounds, compare $\varphi$ with $\frac{(s-1)^2}2$. (d) $n\int_0^\infty v^je^{-nv^2/2}\dd v$ for $j=0,1,2$. (e) The sum is
$\frac{Q(n)+1}{Q(n)+R(n)}$.
\end{hint}

\begin{hint}{pb:mappings}
(a) $x,f(x),\dots,f^{k-1}(x)$ distinct and $f^k(x)=x$. (b) Each cycle of length $k$ has $k$ points. (c) Compare $\frac{n^{\underline k}}{n^k}$ with $e^{-k^2/(2n)}$ for $k\le n^{2/3}$, then split
$\sum_k\frac1ke^{-k^2/(2n)}$ into $\sum_{k\le\sqrt n}\frac1k$ and a Riemann sum. You will need $\int_0^\infty(e^{-u}-[u\le1])\frac{\mathrm du}u=-\gamma$.
\end{hint}

\begin{hint}{pb:xplusinv}
Work with $y_n=x_n^2$, which satisfies $y_{n+1}=y_n+2+\frac1{y_n}$. Then $\frac1{y_k}-\frac1{2k+1}=O(\frac{\ln k}{k^2})$, and $\sum_{k=0}^{n-1}\frac1{2k+1}=H_{2n}-\frac12H_n$.
\end{hint}

\begin{hint}{pb:hyperfactorial}
Differentiate term by term; the factor $(s)_{2k-1}$ vanishes at $s=-1$ for $k\ge2$, so only its derivative survives there. Take $M=3$.
\end{hint}

\begin{hint}{pb:banach}
(a) The box found empty is chosen for the $(N+1)$-th time while the other has been chosen $N-r$ times. (b) Sum the relation over $r$. (c) $\binom{2N}N4^{-N}=\frac1{\sqrt{\pi N}}(1+O(\frac1N))$
(Volume~5, Proposition~13.12).
\end{hint}

\begin{hint}{pb:primezeta}
(a) Take logarithms in the Euler product (Volume~5, Theorem~12.3). (b) $\sum_{k\mid m}\mu(k)=[m=1]$. (c) $\log\zeta(\sigma)\le\zeta(\sigma)-1\le2^{1-\sigma}$ for $\sigma\ge3$.
\end{hint}

\begin{hint}{pb:dickmanint}
(a) Both sides have the same derivative and agree at $u=1$. If $u_0$ were the first zero of $\rho$, the identity at $u_0$ would fail. (b) Integrate $e^{-su}(u\rho(u))'$ by parts.
(d) Split $\int_0^s=\int_0^1+\int_1^s$ and use $\gamma=\int_0^1\frac{1-e^{-t}}t\dd t-\int_1^\infty\frac{e^{-t}}t\dd t$, which follows from $\Gamma'(1)=-\gamma$ (Volume~5, Exercise~7.7).
(e) $\sum\frac1n=\int_{1^-}^\infty\frac{\mathrm d\Psi(t,y)}t=\int_1^\infty\frac{\Psi(t,y)}{t^2}\dd t$.
\end{hint}

\begin{hint}{pb:polymertens}
(a) Expand $-\log(1-q^{-k})=\sum_j\frac{q^{-jk}}j$ and group the terms with $jk=m\le n$. The others have $j\ge2$ and $(j-1)k\ge\frac n2$. (c) $H_n=\log n+\gamma+\frac1{2n}+O(\frac1{n^2})$.
\end{hint}

\section{Solutions}

\begin{solution}{ex:agmsteps}
\idea Quadratic convergence means $e_{n+1}\le e_n^2$ for a suitably normalised error, and then $e_n\le e_0^{2^n}$: the exponent doubles.
\writeup Proposition~\ref{prop:agm}(b) gives $a_{n+1}-b_{n+1}\le\frac{(a_n-b_n)^2}{8b}$, that is $e_{n+1}\le e_n^2$, so $e_n\le e_0^{2^n}$ by induction. Here
$e_0=\frac{1-0.7071}{8\cdot0.7071}\approx0.0518$, so $\log_{10}e_0\approx-1.286$ and $a_n-b_n\le5.66\cdot10^{-1.286\cdot2^n}$. Since $M$ lies between $b_n$ and $a_n$, the error of $a_n$
is below $a_n-b_n$. For $n=10$, $1.286\cdot1024>1300$. (The bound is pessimistic: the actual difference after $10$ steps is about $10^{-1396}$.)
\end{solution}

\begin{solution}{ex:eulerbounds}
\idea The perimeter is an average of $\sqrt{a^2\cos^2+b^2\sin^2}$; comparing a quadratic mean with an arithmetic mean gives one bound, and Cauchy--Schwarz the other.
\writeup With $c=\cos^2\theta$, $s=\sin^2\theta$, $c+s=1$: $(ac+bs)^2\le a^2c+b^2s$ by convexity of $x\mapsto x^2$ (with weights $c$, $s$), so
$\sqrt{a^2c+b^2s}\ge ac+bs$, and integrating over $[0,\frac\pi2]$, where $\int c=\int s=\frac\pi4$, gives $J\ge\frac{\pi(a+b)}4$. By Cauchy--Schwarz,
$J\le\sqrt{\frac\pi2}\sqrt{\int(a^2c+b^2s)}=\sqrt{\frac\pi2\cdot\frac{\pi(a^2+b^2)}4}$. Multiply by $4$. Equality in the first needs $ac+bs$ to be the quadratic mean for all
$\theta$, which forces $a=b$; in the second, the integrand must be constant, which again forces $a=b$. Figure~\ref{fig:perimeter} shows both bounds; the lower one is
much better for moderate eccentricity.
\end{solution}

\begin{solution}{ex:pendulum}
\idea Energy conservation turns the period into an integral $\int\frac{\mathrm d\theta}{\dot\theta}$, and a trigonometric substitution puts it in the form $K(k)$; Gauss's theorem
then gives the AGM.
\writeup Multiplying the equation by $\dot\theta$ and integrating, $\dot\theta^2=\frac{2g}\ell(\cos\theta-\cos\theta_0)=\frac{4g}\ell\bigl(k^2-\sin^2\frac\theta2\bigr)$ with $k=\sin\frac{\theta_0}2$.
By symmetry, the period is four times the time from $\theta=0$ to $\theta_0$, so $T=4\int_0^{\theta_0}\frac{\mathrm d\theta}{\dot\theta}=2\sqrt{\frac\ell g}\int_0^{\theta_0}\frac{\mathrm d\theta}{\sqrt{k^2-\sin^2(\theta/2)}}$.
Substitute $\sin\frac\theta2=k\sin\varphi$, $\varphi\in[0,\frac\pi2]$: then $\frac12\cos\frac\theta2\dd\theta=k\cos\varphi\dd\varphi$, $\sqrt{k^2-\sin^2\frac\theta2}=k\cos\varphi$ and
$\cos\frac\theta2=\sqrt{1-k^2\sin^2\varphi}$, so
\[
T=4\sqrt{\frac\ell g}\int_0^{\pi/2}\frac{\mathrm d\varphi}{\sqrt{1-k^2\sin^2\varphi}}=4\sqrt{\frac\ell g}K(k).
\]
With $k'=\cos\frac{\theta_0}2$ and $K(k)=\frac\pi{2M(1,k')}$ this is the formula. At $\theta_0=\frac\pi2$, $M(1,\cos\frac\pi4)=M(1,0.7071)=0.8472\ldots$, so the period is $18.0\%$ longer.
(For $10^\circ$ it is $0.19\%$ longer; for $179^\circ$, almost four times as long.)
\end{solution}

\begin{solution}{ex:gamma14}
\idea Two different evaluations of the same integral: one by the Beta function (Chapter~7), one by the AGM (this chapter).
\writeup Chapter~7 gives $\int_0^1\frac{\mathrm dr}{\sqrt{1-r^4}}=\frac{\Gamma(1/4)^2}{4\sqrt{2\pi}}$, and Corollary~\ref{cor:lemniscate} gives $\frac\pi{2M(1,\sqrt2)}$. Equating,
$\Gamma(\frac14)^2=\frac{2\pi\sqrt{2\pi}}{M(1,\sqrt2)}=\frac{(2\pi)^{3/2}}{M(1,\sqrt2)}$. After five steps, $a_5=1.19814023473559220744$, already equal to $M(1,\sqrt2)$ to all twenty-one digits shown, and the formula gives
$\Gamma(\frac14)=3.62560990822\ldots$, which agrees with the value computed from Euler's integral.
\end{solution}

\begin{solution}{ex:agmlog}
\idea For small $x$, $\frac1{\sqrt{x^2+t^2}}$ is close to $\frac1t$ except near $0$, where it is cut off at $t\approx x$: this produces $\ln\frac1x$, and the constants come from
exact antiderivatives.
\writeup By Theorem~\ref{thm:gaussagm} and Step~1 of its proof, $\frac\pi{2M(1,x)}=I(1,x)=\int_0^\infty\frac{\mathrm dt}{\sqrt{(1+t^2)(x^2+t^2)}}$. Split at $t=1$.

On $[1,\infty[$, the integrand converges to $\frac1{t\sqrt{1+t^2}}$ as $x\to0$, uniformly, and is bounded by it, so by dominated convergence (Corollary~\ref{cor:dct}) that part tends to
$\int_1^\infty\frac{\mathrm dt}{t\sqrt{1+t^2}}=\operatorname{arsinh}1=\ln(1+\sqrt2)$, using the antiderivative $-\operatorname{arsinh}\frac1t$.

On $[0,1]$, write the integrand as $\frac1{\sqrt{x^2+t^2}}+\frac{g(t)}{\sqrt{x^2+t^2}}$ with $g(t)=\frac1{\sqrt{1+t^2}}-1$. First, $\int_0^1\frac{\mathrm dt}{\sqrt{x^2+t^2}}=\operatorname{arsinh}\frac1x
=\ln\frac{1+\sqrt{1+x^2}}x=\ln\frac2x+o(1)$. Second, $|g(t)|\le\frac{t^2}2$, so $\bigl|\frac{g(t)}{\sqrt{x^2+t^2}}\bigr|\le\frac t2$, and by Corollary~\ref{cor:dct} (the convergence is uniform, the difference being at most $\frac x2$)
$\int_0^1\frac{g(t)}{\sqrt{x^2+t^2}}\dd t\to\int_0^1\frac{g(t)}t\dd t$. This integral equals $\bigl[-\operatorname{arsinh}\frac1t-\ln t\bigr]_0^1=\bigl[-\ln\bigl(1+\sqrt{1+t^2}\bigr)\bigr]_0^1=\ln2-\ln(1+\sqrt2)$.

Adding the three limits, $I(1,x)=\ln\frac2x+\ln2-\ln(1+\sqrt2)+\ln(1+\sqrt2)+o(1)=\ln\frac4x+o(1)$. Numerically, one finds the error $o(1)$ to be $6.8\cdot10^{-3}$ at $x=0.1$, $1.25\cdot10^{-4}$
at $x=0.01$ and $3.6\cdot10^{-12}$ at $x=10^{-6}$, consistent with the known bound $O(x^2\ln\frac1x)$ (Borwein and Borwein, Chapter~7). Taking $x=2^{-m}$ gives $(m+2)\ln2\approx\frac\pi{2M(1,2^{-m})}$, with an error of order $m4^{-m}$: with $\pi$ from Brent--Salamin, this computes $\ln2$ to about $0.6m$ digits.
\end{solution}

\begin{solution}{pb:cyclewalk}
\idea The position after $n$ steps is a sum $S_n=X_1+\dots+X_n$ of independent $\pm1$, read modulo $k$. To pick out the event
$S_n\equiv0$ one uses the roots of unity filter, and the expectation of $\omega^{jS_n}$ factorises by independence.

\writeup (a) Let $\omega=e^{2i\pi/k}$. For an integer $s$, $\frac1k\sum_{j=0}^{k-1}\omega^{js}$ is $1$ if $k\mid s$ and $0$
otherwise (Volume~1, Chapter~5). So
\[
q_n=\mathbb E\Bigl[\frac1k\sum_j\omega^{jS_n}\Bigr]=\frac1k\sum_j\prod_{i=1}^n\mathbb E\bigl[\omega^{jX_i}\bigr]
=\frac1k\sum_j\Bigl(\frac{\omega^j+\omega^{-j}}2\Bigr)^n=\frac1k\sum_j\cos^n\frac{2\pi j}k,
\]
using independence (Volume~4, Chapter~2).

(b) For $k=3$ the cosines are $1,-\frac12,-\frac12$, so $q_n=\frac13\bigl(1+2(-\frac12)^n\bigr)$. For $k=4$ they are
$1,0,-1,0$, so $q_n=\frac14\bigl(1+(-1)^n\bigr)$ for $n\ge1$ (and $q_0=1$). For $k=6$ they are $1,\frac12,-\frac12,-1,-\frac12,
\frac12$, so $q_n=\frac16\bigl(1+2\cdot2^{-n}+2(-\frac12)^n+(-1)^n\bigr)$.

(c) $\cos\frac{2\pi j}k=\pm1$ only for $j=0$ and, if $k$ is even, $j=\frac k2$. The other terms tend to $0$. So $q_n\to\frac1k$
if $k$ is odd; if $k$ is even, $q_n\to\frac2k$ along even $n$ and $q_n=0$ for odd $n$ (an even polygon is bipartite).
\tools Volume~1, Chapter~5 (roots of unity filter); Volume~4, Chapter~2 (independence) and Chapter~8 (periodic chains).
\end{solution}

\begin{solution}{pb:detsquare}
\idea The determinant is a signed sum of $n!$ products. Squaring it creates pairs of products, and the expectation kills every
pair in which some entry appears an odd number of times. Only the $n!$ diagonal pairs survive, each contributing $1$. A random
variable that has mean square $n!$ must sometimes be at least $\sqrt{n!}$ in absolute value.

\writeup (a) $\det A=\sum_\sigma\operatorname{sgn}(\sigma)\prod_ia_{i\sigma(i)}$. Each product has mean $0$, its factors being
independent with mean $0$. Next,
\[
(\det A)^2=\sum_{\sigma,\tau}\operatorname{sgn}(\sigma\tau)\prod_ia_{i\sigma(i)}a_{i\tau(i)}.
\]
If $\sigma=\tau$ the product is $\prod_ia_{i\sigma(i)}^2=1$. If $\sigma\ne\tau$, pick $i$ with $\sigma(i)\ne\tau(i)$: the entry
$a_{i\sigma(i)}$ appears exactly once in the product (row $i$ contributes $a_{i\sigma(i)}a_{i\tau(i)}$, and the other rows
contribute other rows' entries), so the product is $a_{i\sigma(i)}$ times an independent variable, and has mean $0$. Hence
$\mathbb E[(\det A)^2]=n!$ (Volume~4, Chapter~2).

(b) If $|\det A|<\sqrt{n!}$ for every $\pm1$ matrix, then $\mathbb E[(\det A)^2]<n!$.

(c) Let $r_1,\dots,r_n$ be the rows. If they are dependent, $\det A=0$. Otherwise, Gram--Schmidt gives
$r_i^*=r_i-\sum_{j<i}c_{ij}r_j^*$ with $r_1^*,\dots,r_n^*$ orthogonal (Volume~1, Chapter~7). Subtracting from a row a combination
of the previous rows does not change the determinant, so $\det A=\det A^*$ with $A^*$ the matrix of rows $r_i^*$, and
$(\det A^*)^2=\det(A^*A^{*\mathsf T})=\prod_i\|r_i^*\|^2$. As $r_i^*$ is the orthogonal projection of $r_i$ on the orthogonal
of the previous rows, $\|r_i^*\|\le\|r_i\|\le\sqrt n$. So $|\det A|\le n^{n/2}$. By Stirling, $\frac{\sqrt{n!}}{n^{n/2}}\sim
(2\pi n)^{1/4}e^{-n/2}$: the typical size of a random determinant is exponentially smaller than the maximum.
\tools Volume~4, Chapter~2 (expectation, independence); Volume~1, Chapter~7 (determinants, Gram--Schmidt) and Chapters~5--6 (Stirling's formula).
\leads Tao and Vu (2006) proved that $|\det A|=\sqrt{n!}\,e^{O(\sqrt{n\log n})}$ with probability tending to $1$, and Nguyen and
Vu (2014) that $\log|\det A|$ is asymptotically normal. Equality in Hadamard's bound is Problem~\ref{pb:hadamard}.
\end{solution}

\begin{solution}{pb:hadamard}
\idea A Hadamard matrix is a set of $n$ pairwise orthogonal $\pm1$ vectors. Orthogonality of $\pm1$ vectors means they agree
in exactly half of the coordinates, and three rows force a split into quarters. For a construction, one needs $\pm1$
vectors that are ``uncorrelated'', and the quadratic residues modulo $p$ are the classical source: shifting the Legendre
symbol by $d$ leaves it almost orthogonal to itself (Jacobsthal's sum). A border of ones and a diagonal fix the ``almost''.

\writeup (a) Multiplying a column by $-1$ keeps the Hadamard property ($HH^{\mathsf T}=\sum_c(\text{column }c)(\text{column
}c)^{\mathsf T}$), so we may assume row $1$ is all ones. Let $n\ge3$, and let $\alpha,\beta,\gamma,\delta$ be the numbers of
columns where rows $2$ and $3$ carry $(+,+)$, $(+,-)$, $(-,+)$, $(-,-)$. Orthogonality of rows $1,2,3$ gives
\[
\alpha+\beta-\gamma-\delta=0,\qquad\alpha-\beta+\gamma-\delta=0,\qquad\alpha-\beta-\gamma+\delta=0,
\]
and $\alpha+\beta+\gamma+\delta=n$. Adding all four, $4\alpha=n$; similarly $\beta=\gamma=\delta=\frac n4$. So $4\mid n$.

(b) By block multiplication,
\[
\begin{pmatrix}H&H\\H&-H\end{pmatrix}\begin{pmatrix}H^{\mathsf T}&H^{\mathsf T}\\H^{\mathsf T}&-H^{\mathsf T}\end{pmatrix}
=\begin{pmatrix}2HH^{\mathsf T}&0\\0&2HH^{\mathsf T}\end{pmatrix}=2nI_{2n}.
\]

(c) Since $p\equiv3\pmod4$, $\chi(-1)=-1$ (Volume~2, Chapter~6), so $Q_{ba}=\chi(a-b)=-\chi(b-a)=-Q_{ab}$. Each row sums to
$\sum_b\chi(b-a)=\sum_c\chi(c)=0$, since there are as many residues as non-residues. For the diagonal of $QQ^{\mathsf T}$:
$\sum_c\chi(c-a)^2=p-1$. Off the diagonal, with $d=a-b\ne0$ and $c$ replaced by $c+b$:
\[
(QQ^{\mathsf T})_{ab}=\sum_c\chi(c-a)\chi(c-b)=\sum_c\chi(c)\chi(c+d).
\]
For $c\ne0$, $\chi(c)\chi(c+d)=\chi(c^2)\chi(1+dc^{-1})=\chi(1+dc^{-1})$, and as $c$ runs over $\mathbb F_p^*$, $1+dc^{-1}$
runs over $\mathbb F_p\setminus\{1\}$. So the sum is $\sum_{u\ne1}\chi(u)=-\chi(1)=-1$. Hence $QQ^{\mathsf T}=pI-J$.

Now let $S=\begin{pmatrix}0&\mathbf1^{\mathsf T}\\-\mathbf1&Q\end{pmatrix}$, of order $p+1$, and $H=I+S$. The diagonal of
$Q$ is $\chi(0)=0$, so $H$ has entries $\pm1$. Next $S^{\mathsf T}=-S$, and
\[
SS^{\mathsf T}=\begin{pmatrix}0&\mathbf1^{\mathsf T}\\-\mathbf1&Q\end{pmatrix}\begin{pmatrix}0&-\mathbf1^{\mathsf T}\\
\mathbf1&Q^{\mathsf T}\end{pmatrix}=\begin{pmatrix}p&(Q\mathbf1)^{\mathsf T}\\Q\mathbf1&J+QQ^{\mathsf T}\end{pmatrix}=pI_{p+1},
\]
since $Q\mathbf1=0$. Therefore $HH^{\mathsf T}=(I+S)(I-S)=I-S^2=I+SS^{\mathsf T}=(p+1)I_{p+1}$.

(d) Paley's construction gives the orders $4$, $8$, $12$, $20$, $24$, $32$, $44$, $48$, $60$ (from $p=3,7,11,19,23,31,43,47,59$), Sylvester's doubling
of $1$ gives $1$, $2$, $4$, $8$, $16$, $32$, $64$, and doubling the Paley orders gives $24,40,48,64$ and then $48$ again. So up to $64$ we obtain
every multiple of $4$ except $28$, $36$, $52$ and $56$. (Those exist too, by other constructions of Paley and Williamson.)
\tools Volume~2, Chapter~6 (Legendre symbol, $\chi(-1)$, residue counts); Volume~1, Chapter~7 (block matrices); Volume~3,
Chapter~9 (the field $\mathbb F_p$).
\leads The \emph{Hadamard conjecture} says that a Hadamard matrix exists for every multiple of $4$; it is open, and the
smallest unknown order is $668$ (the order $428$ was found by Kharaghani and Tayfeh-Rezaie in 2005). In computer science the
Sylvester matrices are the Walsh--Hadamard transform, and their rows form the Hadamard code.
\end{solution}

\begin{solution}{pb:rank}
\idea A pair $(x,y)$ is fixed by $g$ exactly when both points are, so $F^2$ counts the fixed points of $g$ on $X\times X$, and
Burnside turns its mean into a number of orbits. The diagonal is one orbit; the rest is at least one more.

\writeup (a) By Burnside, $\mathbb E[F]$ is the number of orbits of $G$ on $X$, which is $1$. The fixed points of $g$ on $X\times X$
are the pairs of fixed points, $F(g)^2$ of them, so $\mathbb E[F^2]=r$. The diagonal $\{(x,x)\}$ is an orbit, because $G$ is
transitive, and its complement is non-empty since $n\ge2$; so $r\ge2$.

(b) $r=2$ if and only if the off-diagonal pairs form a single orbit, which is $2$-transitivity. $\operatorname{Var}F=r-1$, so
this is also $\operatorname{Var}F=1$. Finally $\pi$ is real-valued, so $\langle\pi,\pi\rangle=\frac1{|G|}\sum\pi(g)^2=r$ and
$\langle\pi,1\rangle=\mathbb E[F]=1$. With $\psi=\pi-1$, $\langle\psi,1\rangle=0$ and $\langle\psi,\psi\rangle=r-2\langle\psi,1\rangle-1=r-1$.
A character $\psi=\sum m_i\chi_i$ has $\langle\psi,\psi\rangle=\sum m_i^2$, so it is irreducible if and only if
$\langle\psi,\psi\rangle=1$ (Volume~3, Chapter~11), that is, $r=2$.

(c) A rotation or reflection of the $n$-gon preserves the distance between two vertices along the polygon, which takes the
values $0,1,\dots,\lfloor\frac n2\rfloor$. Conversely, if $(x,y)$ and $(x',y')$ are at the same distance, a rotation takes $x$ to $x'$
and then $y$ to $y'$ or to the mirror image of $y'$ in the axis through $x'$, and the reflection in that axis repairs it. So
the orbits on pairs are the distance classes, and $r=\lfloor\frac n2\rfloor+1$. For $n=3$ this is $2$: $D_3=\mathfrak S_3$ is
$2$-transitive on the triangle.
\tools Volume~3, Chapters~3 and~11 (Burnside, inner products of characters); Volume~4, Chapter~2 (variance).
\leads Problem~\ref{pb:fixmoments} is the case $G=\mathfrak S_n$, where all moments are Bell numbers. The higher moments
$\mathbb E[F^j]$ count orbits on $X^j$, and a group is $j$-transitive exactly when they agree with those of $\mathfrak S_n$ up to
order $j$.
\end{solution}

\begin{solution}{pb:derangements}
\idea We know $\mathbb E[F]=1$ and $\mathbb E[F^2]=r\ge2$. Find a quadratic function of $F$ that is $\le0$ whenever $F\ge1$ and
positive at $F=0$: $(F-1)(F-n)$ is the natural choice, since $1\le F\le n$ when $F\ge1$. Its mean is $r-1\ge1$, and only the
derangements can produce it.

\writeup (a) Let $D$ be the set of $g$ with $F(g)=0$. For $g\notin D$, $1\le F(g)\le n$, so $(F-1)(F-n)\le0$. For $g\in D$ it is $n$.
So
\[
n\,P(D)\ge\mathbb E[(F-1)(F-n)]=\mathbb E[F^2]-(n+1)\mathbb E[F]+n=r-1\ge1,
\]
by Problem~\ref{pb:rank}. Hence $P(D)\ge\frac1n$.

(b) Equality needs $r=2$ and $(F-1)(F-n)=0$ for every $g\notin D$, that is $F(g)\in\{0,1,n\}$ for all $g$. Since the action is
faithful, $F(g)=n$ only for $g=e$. So equality holds if and only if $G$ is $2$-transitive and every $g\ne e$ fixes at most one
point. If so, the stabiliser of two distinct points is trivial, so the $g$ with $gx=x'$, $gy=y'$ is unique: $G$ is sharply
$2$-transitive. Conversely, sharply $2$-transitive means $2$-transitive with trivial two-point stabilisers, so every $g\ne e$
fixes at most one point.

For $t\mapsto at+b$ on $\mathbb F_p$: given $x\ne y$ and $x'\ne y'$, the system $ax+b=x'$, $ay+b=y'$ has the unique solution
$a=\frac{x'-y'}{x-y}\ne0$, $b=x'-ax$. So the group, of order $p(p-1)$, is sharply $2$-transitive, and its derangements are the
$p-1$ translations $t\mapsto t+b$, $b\ne0$: a proportion $\frac{p-1}{p(p-1)}=\frac1p$.

(c) Let $H<G$ be proper and suppose $G=\bigcup_{g\in G}gHg^{-1}$. Then $G$ acts transitively on the
$n=[G:H]\ge2$ cosets $gH$, and $x$ fixes $gH$ if and only if $g^{-1}xg\in H$, that is $x\in gHg^{-1}$. So every element would fix a
coset. The action may not be faithful, but (a) applies to the faithful action of $G/K$, $K$ the kernel, and gives an element
with no fixed point. Contradiction.
\tools Volume~3, Chapter~3 (actions on cosets); Volume~4, Chapter~2 (expectations).
\leads Jordan's theorem (1872) has a famous arithmetic consequence: an irreducible polynomial of degree at least $2$ in
$\Z[X]$ has no root modulo $p$ for infinitely many primes $p$, a set of positive density by Chebotarev's theorem (see Serre's
article below). Fulman and Guralnick proved a conjecture of Boston and Shalev: for simple groups the proportion of
derangements is bounded below by an absolute constant.
\end{solution}

\begin{solution}{pb:noroots}
\idea Evaluation at the $q$ points of $\mathbb F_q$ is a linear map, onto as soon as the degree allows interpolation. The roots then
behave like $q$ independent coins of bias $\frac1q$, exactly when $n\ge q$, and like the first terms of inclusion--exclusion
otherwise.

\writeup Write $f=x^n+g$ with $g$ uniform among polynomials of degree $<n$. For a set $E\subset\mathbb F_q$ with $|E|\le n$, the linear map
$g\mapsto(g(a))_{a\in E}$ is onto (Lagrange interpolation with polynomials of degree $<|E|$), so it takes every value equally often,
and so does $f\mapsto(f(a))_{a\in E}$, which is a translate of it.

(a) If $n\ge q$, take $E=\mathbb F_q$: $(f(a))_a$ is uniform on $\mathbb F_q^q$, the coordinates are independent and uniform, the events
$f(a)=0$ are independent with probability $\frac1q$, and $R\sim\mathrm{Bin}(q,\frac1q)$. Hence $P(R=0)=(1-\frac1q)^q$.

(b) By inclusion--exclusion,
\[
P(R=0)=\sum_{E\subset\mathbb F_q}(-1)^{|E|}P(f=0\text{ on }E).
\]
For $|E|\le n$ this probability is $q^{-|E|}$ by the
remark above; for $|E|>n$ it is $0$, a monic polynomial of degree $n$ having at most $n$ roots. So
$P(R=0)=\sum_{k\le n}(-1)^k\binom qkq^{-k}$, and since
\[
\binom qkq^{-k}=\frac1{k!}\prod_{i<k}\Bigl(1-\frac iq\Bigr)\longrightarrow\frac1{k!},
\]
as $q\to\infty$ this tends to $\sum_{k\le n}\frac{(-1)^k}{k!}$, the derangement probability: a random polynomial has no root as often as a random
permutation has no fixed point, in agreement with Problem~\ref{pb:randpoly}(c), roots being factors of degree $1$.

\tools Volume~1, Chapter~6 (Lagrange interpolation); Volume~4, Chapter~1 (inclusion--exclusion, binomial law); Chapter~4 of this book
(derangements).
\leads For $n\ge q$ the answer does not depend on $n$ at all. Over $\Z$ the analogue is the number of roots modulo $p$ of a random
polynomial, and for a fixed polynomial the Chebotarev density theorem gives the law of the number of roots modulo a random prime.
\end{solution}

\begin{solution}{pb:gaussmeasure}
\idea The preimage of $[0,t[$ is a union of intervals, one for each value of the first partial quotient, and their Gauss measures
telescope back to $\log_2(1+t)$.

\writeup (a) $Tx<t$ with $\lfloor1/x\rfloor=k$ means $k\le\frac1x<k+t$, that is $x\in\,]\frac1{k+t},\frac1k]$. So
\[
\begin{aligned}
\mu(T^{-1}[0,t[)&=\frac1{\log2}\sum_{k\ge1}\Bigl(\log\Bigl(1+\frac1k\Bigr)-\log\Bigl(1+\frac1{k+t}\Bigr)\Bigr)\\
&=\frac1{\log2}\sum_{k\ge1}\Bigl(\log\frac{k+1}k-\log\frac{k+1+t}{k+t}\Bigr).
\end{aligned}
\]
The partial sum up to $K$ is $\log(K+1)-\log\frac{K+1+t}{1+t}\to\log(1+t)$, so the value is $\mu([0,t[)$. Two probability measures that
agree on all intervals $[0,t[$ agree on all Borel sets, so $\mu\circ T^{-1}=\mu$.

(b) $a_n(x)=a_1(T^{n-1}x)$, and $T^{n-1}x$ has law $\mu$ by (a) and induction. So
$P(a_n=k)=\mu(]\frac1{k+1},\frac1k])=\log_2\frac{1+1/k}{1+1/(k+1)}=\log_2\frac{(k+1)^2}{k(k+2)}$.

(c) $P(a_1=k)=\frac1k-\frac1{k+1}=\frac1{k(k+1)}$. For $k=1$: $0.5$ for uniform $x$, $\log_2\frac43=0.415$ for the Gauss measure.

\tools Volume~2, Chapter~8; Volume~4, Chapter~5 (laws determined by distribution functions).
\leads $\mathbb E[a_1]=\infty$ under both measures, since the probabilities decay like $\frac1{k^2}$. Khinchin proved that nevertheless
$(a_1\cdots a_n)^{1/n}\to2.6854\ldots$ for almost every $x$, a consequence of the ergodic theorem (Part~II).
\end{solution}

\begin{solution}{pb:digitsclt}
\idea For $N=2^k$ the bits of a uniform $n$ are independent fair coins. A general $N$ is a union of a few dyadic blocks, each a shifted
copy of the first case, and almost all the mass sits in the big blocks.

\writeup (a) $n\mapsto$ its $k$ bits is a bijection from $\{0,\dots,2^k-1\}$ to $\{0,1\}^k$, so the bits are independent and fair, and
$s(n)$ is their sum.

(b) Write $N=2^{k_1}+\dots+2^{k_r}$ with $k_1>\dots>k_r$, and $M_i=2^{k_1}+\dots+2^{k_{i-1}}$. The interval $[0,N[$ is the disjoint union of
the blocks $B_i=[M_i,M_i+2^{k_i}[$, and for $n\in B_i$, $s(n)=(i-1)+s(n-M_i)$, where $n-M_i$ is uniform on $[0,2^{k_i}[$ when $n$ is uniform
on $B_i$. So the law of $s(n)$ is the mixture, with weights $w_i=2^{k_i}/N$, of the laws of $(i-1)+\mathrm{Bin}(k_i,\frac12)$.

Let $k=k_1$, so $\log_2N\in[k,k+1[$, and fix $h\ge1$. The blocks with $i>h$ have total weight at most $2^{k-h+1}/N\le2^{1-h}$. For
$i\le h$, $k_i\in[k-h,k]$ and the shift is at most $h$; by the de Moivre--Laplace theorem (Volume~4, Chapter~6),
$\bigl((i-1)+\mathrm{Bin}(k_i,\frac12)-\frac12\log_2N\bigr)/\sqrt{\frac14\log_2N}\to N(0,1)$ in law, uniformly in $i\le h$, because the shift
and the change of variance are $O(h)$ and $O(h/k)$ against a standard deviation $\sqrt k/2$. Hence for every $x$, the distribution
function of the normalised $s(n)$ at $x$ is within $2^{1-h}+o(1)$ of $\Phi(x)$. Let $N\to\infty$, then $h\to\infty$.

\tools Volume~4, Chapter~6 (central limit theorem); Volume~4, Chapter~2 (binomial law).
\leads The mean $\frac1N\sum_{n<N}s(n)=\frac12\log_2N+F(\log_2N)$ with $F$ a continuous, nowhere differentiable periodic function
(Trollope 1968, Delange 1975): the blocks make the mean oscillate. Gelfond (1968) proved that $s(n)\bmod2$ is equidistributed in
arithmetic progressions, and Mauduit and Rivat (2010) that it is equidistributed along the primes.
\end{solution}

\begin{solution}{pb:kolberg}
\idea If $p(n)$ had eventually constant parity, the generating function of partitions would be a rational function modulo $2$, and
so would its inverse, the pentagonal series. But the pentagonal series has longer and longer gaps, which a rational series over
a finite field cannot have.

\writeup Work in $\mathbb F_2[[x]]$ and let $\bar P=\sum_n(p(n)\bmod2)x^n$. If $p(n)$ is even for all $n\ge N$, then $\bar P$ is a polynomial;
if $p(n)$ is odd for all $n\ge N$, then $\bar P$ is a polynomial plus $\frac{x^N}{1-x}$. In both cases $\bar P=\frac UV$ with polynomials
$U,V$ and $U(0)=V(0)=1$ (as $p(0)=1$). Since $P(x)E(x)=1$ in $\Z[[x]]$, the reduction $\bar E$ of $E$ is $\frac VU$, rational too.

A rational series $\frac VU=\sum e_nx^n$ over $\mathbb F_2$ with $U=1+u_1x+\dots+u_dx^d$ satisfies $e_n=u_1e_{n-1}+\dots+u_de_{n-d}$ for
$n>\max(d,\deg V)$. The vector $(e_{n-1},\dots,e_{n-d})$ takes at most $2^d$ values, so it repeats, and from then on the sequence
is periodic. An eventually periodic sequence with infinitely many non-zero terms has bounded gaps between them.

But by the pentagonal number theorem, $e_n\ne0$ exactly when $n=\frac{j(3j\mp1)}2$ for some $j\ge0$, and the gap between the
consecutive values $\frac{j(3j+1)}2$ and $\frac{(j+1)(3j+2)}2$ is $2j+1$, unbounded. This contradiction shows that neither parity
can hold from some point on.

\tools Volume~5, Theorem~2.12 (pentagonal number theorem); Volume~5, Theorem~2.5 (rational series and recurrences).
\leads The argument gives no quantitative information. Nicolas, Ruzsa and Sárközy (1998) and Ahlgren (1999) showed that each
parity occurs, up to logarithmic factors, at least $\sqrt N$ times below $N$, still far from the expected $\frac N2$. Every arithmetic
progression contains infinitely many $n$ with $p(n)$ even (Ono, 1996) and infinitely many with $p(n)$ odd (Radu, 2012).
\end{solution}

\begin{solution}{pb:kummer}
\idea Legendre's formula counts the factors $p$ in $n!$ through the digit sum, and each carry in an addition costs exactly $p-1$
in digit sum.

\writeup (a) By Volume~13, Exercise~2.5, $v_p(n!)=\frac{n-s_p(n)}{p-1}$. Hence
\[
v_p\binom{m+n}m=\frac{s_p(m)+s_p(n)-s_p(m+n)}{p-1}.
\]
Add $m$ and $n$ digit by digit. In a position with incoming carry $c_{i-1}\in\{0,1\}$ the digit of the sum is
$m_i+n_i+c_{i-1}-pc_i$, where $c_i$ is the outgoing carry. Summing over $i$ (with $c_{-1}=0$ and no carry beyond the last digit),
$s_p(m+n)=s_p(m)+s_p(n)+\sum_ic_{i-1}-p\sum_ic_i=s_p(m)+s_p(n)-(p-1)\sum_ic_i$. So $v_p\binom{m+n}m=\sum_ic_i$.

(b) Add $k$ and $n-k$. If $n$ has $L+1$ digits ($p^L\le n<p^{L+1}$), a carry out of position $L$ would make the sum at least
$p^{L+1}$, so carries can only occur out of positions $0,\dots,L-1$: at most $L$ of them. Hence $p^{v_p\binom nk}\le p^L\le n$.

(c) By (b), for each prime $p$ the power $p^{v_p\binom nk}$ is at most $n$, hence divides the highest power of $p$ that is $\le n$,
which is the power of $p$ in $\operatorname{lcm}(1,\dots,n)$. So $\binom nk$ divides the lcm. The largest of the $n+1$ numbers
$\binom nk$ is at least their average $\frac{2^n}{n+1}$.

\tools Volume~13, Exercise~2.5 (Legendre's formula via digit sums); Volume~2, Chapter~2 (lcm and prime powers).
\leads The bound (c) is the starting point of Chebyshev-type estimates: $\log\operatorname{lcm}(1,\dots,n)=\psi(n)$, so (c) gives
$\psi(n)\ge n\log2-\log(n+1)$ (compare Volume~6, Chapter~1). Kummer found (a) in 1852, in a long paper on reciprocity laws.
\end{solution}

\begin{solution}{pb:pascalfour}
\idea An entry is $2\bmod4$ when it is divisible by $2$ exactly once, that is, when there is exactly one carry. One carry can only
happen at a place where $n$ has the pattern $10$.

\writeup (a) By Problem~\ref{pb:kummer}, $\binom nk\equiv2\pmod4$ if and only if adding $k$ and $m=n-k$ in binary produces exactly one carry,
say out of position $i$. In position $i$ there is no incoming carry, and $k_i+m_i=n_i+2$, which forces $k_i=m_i=1$ and $n_i=0$. In
position $i+1$ there is an incoming carry and no outgoing one, so $k_{i+1}+m_{i+1}+1=n_{i+1}\le1$, which forces $k_{i+1}=m_{i+1}=0$ and
$n_{i+1}=1$. In every other position there is no carry at all, so $k_j+m_j=n_j$: where $n_j=0$ both digits are $0$, and where $n_j=1$
we may choose which of $k_j,m_j$ is $1$. So the $k$ with exactly one carry correspond to a choice of a block $10$ in $n$ (positions
$i+1,i$) and a free choice in each of the other $s(n)-1$ positions where $n$ has a digit $1$: $B(n)2^{s(n)-1}$ of them.

(b) Rows with $B(n)=0$: binary expansions with no $1$ immediately followed by a $0$, that is, $n=11\cdots1_2=2^a-1$ (and $n=0$).

\tools Problem~\ref{pb:kummer}.
\leads There are similar rules for $\binom nk$ modulo any prime power, in terms of blocks of digits; Granville's survey listed
below explains them, and how such counts become ``fractal'' pictures as in Figure~\ref{fig:pascal3}.
\end{solution}

\begin{solution}{pb:centralp}
\idea The series $f$ satisfies $f^2(1-4x)=1$, and modulo $p$ raising to the power $p$ is the substitution $x\mapsto x^p$. This turns
the algebraic equation into a self-similarity that reads the coefficients digit by digit.

\writeup (a) For a power series $g$ with integer coefficients, $g(x)^p\equiv g(x^p)\pmod p$ (the multinomial coefficients of the
$p$-th power are divisible by $p$ except for the pure powers, and $a^p\equiv a$). Since $f^2=\frac1{1-4x}$,
\[
(1-4x)^{\frac{p-1}2}f(x^p)\equiv(1-4x)^{\frac{p-1}2}f(x)^p=(1-4x)^{\frac{p-1}2}f(x)\bigl(f(x)^2\bigr)^{\frac{p-1}2}=f(x).
\]

(b) Let
\[
Q(x)=(1-4x)^{(p-1)/2}=\sum_{j\le(p-1)/2}\binom{(p-1)/2}j(-4)^jx^j.
\]
Since $\frac{p-1}2\equiv-\frac12\pmod p$ and
$\binom{-1/2}j(-4)^j=\binom{2j}j$ (Volume~4, Exercise~3.7), the coefficient of $x^j$ in $Q$ is $\equiv\binom{2j}j$ for $j\le\frac{p-1}2$; and for
$\frac{p-1}2<j<p$, $\binom{2j}j=\frac{(2j)!}{j!^2}$ is divisible by $p$ because $p\le2j<2p$. So $Q(x)\equiv\sum_{j<p}\binom{2j}jx^j$. Writing
$n=n_0+pn'$ with $0\le n_0<p$, the coefficient of $x^n$ in $Q(x)f(x^p)$ is $\binom{2n_0}{n_0}\binom{2n'}{n'}$ modulo $p$ (only the term
$x^{n_0}$ of $Q$ can be combined with a power of $x^p$, as $\deg Q<p$). Induction on the number of digits gives the product
formula.

(c) A digit $n_i\le\frac{p-1}2$ gives $\binom{2n_i}{n_i}=\frac{(2n_i)!}{n_i!^2}$ with $2n_i<p$, not divisible by $p$; a digit $n_i>\frac{p-1}2$ gives a
multiple of $p$, as in (b). Below $p^k$ the good $n$ are $\bigl(\frac{p+1}2\bigr)^k$ out of $p^k$, a proportion
$\bigl(\frac{p+1}{2p}\bigr)^k\to0$.

\tools Volume~4, Exercise~3.7 (the series of $(1-4x)^{-1/2}$); Volume~2, Chapter~3 (Fermat's little theorem).
\leads The same trick works for any power series $f$ with $f^d\in\Q(x)$, and more generally for diagonals of rational functions.
It gives ``Lucas congruences'' for many combinatorial sequences, such as the Apéry numbers of the irrationality proof of $\zeta(3)$
(Chapter~11; Gessel, 1982). In base $p$
these series are \emph{$p$-automatic}, the theme of Problem~\ref{pb:thuemorse}.
\end{solution}

\begin{solution}{pb:sternfib}
\idea Hyperbinary representations are pairs of binary numbers without common digits, and such pairs are what Lucas's theorem
detects. The bound is an inequality version of the Fibonacci recursion, because of two consecutive integers one is even.

\writeup (a) Given a hyperbinary representation of $n$, let $k=\sum_{e_i=2}2^i$ and $j=\sum_{e_i=1}2^i$. Then $n=2k+j$ and $k$, $j$
have no binary digit $1$ in common. Conversely such a pair $(k,j)$ gives back the $e_i$. So $s(n+1)$ is the number of $k\ge0$
with $2k\le n$ such that $k$ and $j=n-2k$ have disjoint binary supports. By Lucas (Volume~2, Exercise~3.9), $\binom{k+j}k$ is odd if
and only if every binary digit of $k$ is at most the digit of $k+j$ in the same place, that is, if and only if there is no carry in
$k+j$, that is, if and only if $k$ and $j$ have disjoint supports. As $k+j=n-k$, this is the oddness of $\binom{n-k}k$.

(b) Let $M_k=\max\{s(n):2^k\le n\le2^{k+1}\}$, so $M_0=1=F_2$ and $M_1=\max(s(2),s(3),s(4))=2=F_3$. For $n$ in $[2^{k+1},2^{k+2}]$: if
$n=2m$ is even, $s(n)=s(m)\le M_k$. If $n=2m+1$, then $m$ and $m+1$ lie in $[2^k,2^{k+1}]$, and one of them is even, equal to $2r$
with $r\in[2^{k-1},2^k]$, so $s(n)=s(m)+s(m+1)\le M_k+M_{k-1}$. Hence $M_{k+1}\le M_k+M_{k-1}$ and $M_k\le F_{k+2}$ by induction.

(c) Let $a_k=\frac{2^{k+2}-(-1)^k}3$, so $a_1=3$, $a_2=5$, $a_3=11$, and $a_{k+1}=2a_k+\varepsilon_k$ with $\varepsilon_k=(-1)^k$ (check:
$2a_k+(-1)^k=\frac{2^{k+3}-2(-1)^k+3(-1)^k}3=a_{k+1}$). Then $s(a_{k+1})=s(a_k)+s(a_k+\varepsilon_k)$, and
$a_k+\varepsilon_k=2a_{k-1}+\varepsilon_{k-1}+\varepsilon_k=2a_{k-1}$, so $s(a_k+\varepsilon_k)=s(a_{k-1})$. Thus $s(a_{k+1})=s(a_k)+s(a_{k-1})$ with
$s(a_1)=2$, $s(a_2)=3$, and $s(a_k)=F_{k+2}$. Finally $2^k\le a_k\le2^{k+1}$.

\tools Volume~2, Exercise~3.9 (Lucas); Volume~1, Chapter~1 (Fibonacci numbers).
\leads By (a) and the identity $\sum_k\binom{n-k}k=F_{n+1}$, $s(n+1)$ counts the odd numbers on the $n$-th ``shallow diagonal''
of Pascal's triangle, whose sum is $F_{n+1}$. The second maximum point of each row, $\frac{5\cdot2^k+(-1)^k}3$ in the figure, is the
mirror image: $s(2^{k+1}+2^k-n)=s(n)$ for $2^k\le n\le2^{k+1}$, which follows from the reflection symmetry of the Stern--Brocot tree.
See Lehmer's and Northshield's papers below.
\end{solution}

\begin{solution}{pb:thuemorse}
\idea The rules $t_{2n}=t_n$ and $t_{2n+1}=1+t_n$ say that $T(x)$ is expressed through $T(x^2)$, which modulo $2$ is $T(x)^2$. Periodicity
is incompatible with these rules, because a period could be halved until it becomes odd, and an odd period contradicts the
alternation $t_{2n+1}\ne t_{2n}$.

\writeup (a) Since $s(2n)=s(n)$ and $s(2n+1)=s(n)+1$,
\[
T(x)=\sum_nt_nx^{2n}+\sum_n(1+t_n)x^{2n+1}=(1+x)T(x^2)+\frac x{1-x^2}.
\]
In $\mathbb F_2[[x]]$, $T(x^2)=T(x)^2$ and $1-x^2=(1+x)^2$. So $T=(1+x)T^2+\frac x{(1+x)^2}$; multiply by $(1+x)^2$.

(b) Suppose $t_{n+P}=t_n$ for all $n\ge N$, with $P\ge1$ minimal. If $P$ is even, then for $n\ge N$, $t_{n+P/2}=t_{2n+P}=t_{2n}=t_n$, so $P/2$ is a
period from $N$ on, contradicting minimality. So $P$ is odd. For $n\ge N$,
\[
t_{2n+P}=t_{2n}=t_n\quad\text{and}\quad t_{2n+P}=t_{2(n+\frac{P-1}2)+1}=1+t_{n+\frac{P-1}2},
\]
so $t_{n+\frac{P-1}2}=1+t_n$. Similarly $t_{2n+1+P}=t_{2n+1}=1+t_n$ and $t_{2n+1+P}=t_{2(n+\frac{P+1}2)}=t_{n+\frac{P+1}2}$, so $t_{n+\frac{P+1}2}=1+t_n$.
Hence $t_{m+1}=t_m$ for all large $m$: the sequence is eventually constant. But $t_{2n+1}=1+t_{2n}$ for every $n$. Contradiction.

(c) By (a), $T$ is a root of $(1+x)^3Y^2+(1+x)^2Y+x$, so it is algebraic of degree at most $2$. If $T$ were rational, its coefficients
would satisfy a linear recurrence from some point on, and be eventually periodic as in Problem~\ref{pb:kolberg}; this contradicts
(b). So the degree is exactly $2$.

\tools Volume~5, Theorem~2.5 (rational series and recurrences); Volume~3, Chapter~9 (the Frobenius map in characteristic $p$).
\leads \emph{Christol's theorem} (1979): a power series over $\mathbb F_p$ is algebraic over $\mathbb F_p(x)$ if and only if its sequence
of coefficients is \emph{$p$-automatic}, that is, $a_n$ can be computed by a finite automaton reading the base-$p$ digits of $n$. Thue--Morse,
the Catalan numbers modulo $2$ and the central binomial coefficients modulo $p$ (Problem~\ref{pb:centralp}) are examples. See
Allouche and Shallit's book.
\end{solution}

\begin{solution}{pb:hermite}
\idea A polynomial $f$ whose derivatives at $0$ and $1$ are all integers, and an antiderivative of $e^{kx}f(x)$ of the shape $e^{kx}F(x)$,
turn $\int_0^1e^{kx}f$ into an integer combination of $1$ and $e^k$; and the factor $\frac1{n!}$ in $f$ makes the integral tiny.

\writeup (a) Let $f(x)=\frac{x^n(1-x)^n}{n!}=\frac1{n!}\sum_{j=n}^{2n}c_jx^j$ with $c_j\in\Z$. Then $f^{(j)}(0)=\frac{j!}{n!}c_j\in\Z$ for $n\le j\le2n$, and $0$ for
other $j$; since $f(1-x)=f(x)$, also $f^{(j)}(1)\in\Z$. Let $F=\sum_{j=0}^{2n}(-1)^jk^{2n-j}f^{(j)}$. Then
\[
kF+F'=\sum_{j=0}^{2n}(-1)^jk^{2n+1-j}f^{(j)}+\sum_{j=0}^{2n}(-1)^jk^{2n-j}f^{(j+1)}=k^{2n+1}f,
\]
because the sums telescope and $f^{(2n+1)}=0$. So $\frac{\mathrm d}{\mathrm dx}\bigl(e^{kx}F(x)\bigr)=k^{2n+1}e^{kx}f(x)$, and
\[
k^{2n+1}\int_0^1e^{kx}f(x)\dd x=e^kF(1)-F(0).
\]
If $e^k=\frac ab$ with $a,b\ge1$ integers, then
\[
M_n=bk^{2n+1}\int_0^1e^{kx}f(x)\dd x=aF(1)-bF(0)
\]
is an integer. It is positive, since
$f>0$ on $]0,1[$. But $0<f\le\frac{4^{-n}}{n!}$, so $M_n\le\frac{bk e^k(k^2/4)^n}{n!}\to0$. For large $n$, $0<M_n<1$: a contradiction.

(b) Let $r=\frac pq$ with $p\ne0$, $q\ge1$. If $e^r$ were rational, so would be $(e^r)^q=e^p$ and $e^{-p}=\frac1{e^p}$, one of which is $e^k$
with $k=|p|\ge1$; this contradicts (a). If $\log r=\frac pq$ for a rational $r>0$, $r\ne1$, then $p\ne0$ and $e^{p/q}=r$ would be
rational.

\tools{Volume~5, Theorem~3.8 (the same method for $\pi^2$); Volume~1, Exercise~8.4 ($e$ is irrational).}
\leads Hermite (1873) pushed this idea, with several functions $e^{kx}$ at once, to prove that $e$ is transcendental: it is the root of
no non-zero polynomial with integer coefficients. Lindemann (1882) proved that $\pi$ is transcendental, which shows that the circle
cannot be squared with ruler and compass.
\end{solution}

\begin{solution}{pb:centralsum}
\idea $\frac{(n-1)!\,n!}{(2n)!}$ is a Beta integral, and summing the integrands is a geometric series in $t(1-t)$, which is at most
$\frac14$. What remains is a rational function whose denominator $1-t+t^2$ has roots at primitive sixth roots of unity: hence $\pi\sqrt3$.

\writeup (a) By Volume~5, Theorem~7.11, $\int_0^1t^{n-1}(1-t)^n\dd t=B(n,n+1)=\frac{(n-1)!\,n!}{(2n)!}=\frac1{n\binom{2n}n}$. Summing the non-negative
integrands, $\sum_{n\ge1}t^{n-1}(1-t)^n=\frac{1-t}{1-t+t^2}$. By the symmetry $t\mapsto1-t$, which fixes $1-t+t^2$, its integral is half of
\[
\int_0^1\frac{\mathrm dt}{1-t+t^2}=\frac2{\sqrt3}\Bigl[\arctan\frac{2t-1}{\sqrt3}\Bigr]_0^1=\frac2{\sqrt3}\cdot\frac\pi3=\frac{2\pi}{3\sqrt3}.
\]
So the sum is $\frac\pi{3\sqrt3}=0.6046\ldots$

(b) $\frac1{\binom{2n}n}=(2n+1)B(n+1,n+1)=(2n+1)\int_0^1u^n\dd t$ with $u=t-t^2$, and $\sum_{n\ge0}(2n+1)u^n=\frac{1+u}{(1-u)^2}$ (differentiate the geometric
series). Let $q=1-t+t^2=1-u$. Then $\frac{1+u}{(1-u)^2}=\frac{2-q}{q^2}=\frac2{q^2}-\frac1q$. With $t-\frac12=\frac{\sqrt3}2\tan\theta$, $q=\frac34\sec^2\theta$ and
$\theta\in[-\frac\pi6,\frac\pi6]$:
\[
\int_0^1\frac{\mathrm dt}{q^2}=\frac{8\sqrt3}9\int_{-\pi/6}^{\pi/6}\cos^2\theta\dd\theta=\frac{8\sqrt3}9\Bigl(\frac\pi6+\frac{\sqrt3}4\Bigr)=\frac{4\pi\sqrt3}{27}+\frac23.
\]
So $\sum_{n\ge0}\frac1{\binom{2n}n}=\frac{8\pi\sqrt3}{27}+\frac43-\frac{2\pi}{3\sqrt3}=\frac43+\frac{2\pi\sqrt3}{27}=1.7364\ldots$

\tools{Volume~5, Theorem~7.11 (Beta and Gamma); Volume~1, Chapter~8 (integrals of rational functions).}
\leads D.~H.~Lehmer's 1985 article collects dozens of such sums. The appearance of $\sqrt3$ is no accident: these series are values
of $L$-functions of characters modulo $3$, and the same technique with $\frac1{n^3\binom{2n}n}$ leads to Apéry's series for $\zeta(3)$.
\end{solution}

\begin{solution}{pb:lehmer}
\idea Everything follows from $\sum_nc_nx^n=\frac1{\sqrt{1-x}}$, the binomial series for the exponent $-\frac12$. Part (b) hides a telescoping
sum, found by trying a multiple of $nc_n^2$.

\writeup (a) For $|x|<1$, $\sum_nc_nx^{2n}=(1-x^2)^{-1/2}$; integrating from $0$ to $x$, $\arcsin x=\sum_n\frac{c_n}{2n+1}x^{2n+1}$. The terms are
non-negative and $c_n\le1$, so the series converges at $x=1$ and, by monotone convergence, its sum is $\arcsin1=\frac\pi2$. Next,
$\sum_{n\ge1}c_nx^{n-1}=\frac{(1-x)^{-1/2}-1}x$. Integrating from $0$ to $1$ and substituting $v=\sqrt{1-x}$ ($x=1-v^2$, $\mathrm dx=-2v\dd v$):
\[
\sum_{n\ge1}\frac{c_n}n=\int_0^1\frac{1/v-1}{1-v^2}\,2v\dd v=\int_0^1\frac{2\dd v}{1+v}=2\log2,
\]
again with monotone convergence at the end point.

(b) $c_n\sim\frac1{\sqrt{\pi n}}$ (Volume~5, Proposition~13.12), so $c_n^2\sim\frac1{\pi n}$ and $\sum c_n^2$ diverges like the harmonic series. Let $b_n=-4nc_n^2$. Using
$c_{n+1}=\frac{2n+1}{2n+2}c_n$,
\[
b_n-b_{n+1}=4c_n^2\Bigl(\frac{(2n+1)^2}{4(n+1)}-n\Bigr)=4c_n^2\cdot\frac1{4(n+1)}=\frac{c_n^2}{n+1}.
\]
So $\sum_{n=0}^{N-1}\frac{c_n^2}{n+1}=b_0-b_N=4Nc_N^2\to\frac4\pi$.

\tools{Volume~1, Chapter~8 (Taylor series); Volume~5, Proposition~13.12 (central binomial coefficient).}
\leads $\sum c_n^2x^n=\frac2\pi K(\sqrt x)$, a complete elliptic integral, whose behaviour near $x=1$ explains the divergence (Volume~5, Chapter~14, on the
AGM, meets the same function). Ramanujan's series for $\frac1\pi$, such as $\sum_n\binom{2n}n^3\frac{42n+5}{2^{12n+4}}=\frac1\pi$, belong to this family.
\end{solution}

\begin{solution}{pb:catalanmoments}
\idea Under $x=2\sin\theta$ the moment becomes a Wallis-type integral, which the Beta function evaluates.

\writeup (a) Odd moments vanish by symmetry. With $x=2\sin\theta$, $\sqrt{4-x^2}=2\cos\theta$ and $\mathrm dx=2\cos\theta\dd\theta$:
\[
\frac1{2\pi}\int_{-2}^2x^{2n}\sqrt{4-x^2}\dd x=\frac{4^{n+1}}{\pi}\int_0^{\pi/2}\sin^{2n}\theta\cos^2\theta\dd\theta=\frac{4^{n+1}}{2\pi}B\Bigl(n+\frac12,\frac32\Bigr)
\]
by Volume~5, Exercise~7.3. Now $B(n+\frac12,\frac32)=\frac{\Gamma(n+\frac12)\Gamma(\frac32)}{\Gamma(n+2)}$, with $\Gamma(n+\frac12)=\frac{(2n)!}{4^nn!}\sqrt\pi$ and $\Gamma(\frac32)=\frac{\sqrt\pi}2$, so
the moment is $\frac{4^{n+1}}{2\pi}\cdot\frac{(2n)!\,\pi}{2\cdot4^nn!\,(n+1)!}=\frac{(2n)!}{n!\,(n+1)!}=C_n$.

(b) With $x=y^2$, $\mathrm dx=2y\dd y$ and $\sqrt{\frac{4-x}x}\,\mathrm dx=2\sqrt{4-y^2}\dd y$; so the integral is $\frac1{2\pi}\int_0^22y^{2n}\sqrt{4-y^2}\dd y$, which is (a).

(c) Put $x=4(1-u)$: $C_n=\frac{2\cdot4^n}\pi\int_0^1(1-u)^{n-1/2}\sqrt u\dd u$. As $(1-u)^{n}=e^{n\log(1-u)}$ is concentrated near $u=0$, Laplace's method
(Volume~5, Theorem~13.9; here it is even exact: the integral is $B(\frac32,n+\frac12)$) gives $\int_0^1(1-u)^{n-1/2}\sqrt u\dd u\sim\int_0^\infty e^{-nu}\sqrt u\dd u=\Gamma(\frac32)n^{-3/2}=\frac{\sqrt\pi}2n^{-3/2}$.
So $C_n\sim\frac{4^n}{\sqrt\pi\,n^{3/2}}$.

\tools{Volume~5, Exercise~7.3, Proposition~7.12 ($\Gamma(\frac12)$), Theorem~11.9 (Laplace); Volume~4, Exercise~1.7 (Catalan numbers count Dyck paths).}
\leads (a) is the moment form of Wigner's semicircle law (1955): the eigenvalues of a large random symmetric matrix, scaled by
$\sqrt n$, are distributed with density $\frac1{2\pi}\sqrt{4-x^2}$, and the proof computes $\frac1n\operatorname{tr}(A^{2k})$ by counting closed walks that
turn out to be Dyck paths. (b) is the Marchenko--Pastur law for $AA^{\mathsf T}$. Random matrices return in Part~III.
\end{solution}

\begin{solution}{pb:apcount}
\idea A progression is a solution of the linear equation $x+z=2y$, so its count is a Fourier sum. The trivial frequency gives the
number of progressions in a random set of the same density, and the others are controlled by the largest non-trivial
coefficient: small coefficients mean random-looking, and random-looking sets have the expected number of progressions.

\writeup (a) The map $(x,d)\mapsto(x,x+d,x+2d)$ is a bijection from $\mathbb F_p^2$ onto the triples $(x,y,z)$ with $x+z=2y$ (given such a triple,
$d=y-x$ and then $z=2y-x=x+2d$). So
\[
T(A)=\sum_{x,y,z\in A}\frac1p\sum_re\Bigl(\frac{r(x+z-2y)}p\Bigr)=\frac1p\sum_r\hat A(r)^2\hat A(-2r).
\]

(b) The term $r=0$ is $\frac{|A|^3}p$. For $r\ne0$, $-2r\ne0$ since $p$ is odd, so the rest is at most
$\frac1p\max_{s\ne0}|\hat A(s)|\sum_r|\hat A(r)|^2=|A|\max_{s\ne0}|\hat A(s)|$ by Parseval.

(c) Let $\chi$ be the Legendre symbol, $\chi(0)=0$. For $x\ne0$, $1_Q(x)=\frac12(1+\chi(x))$, so for $r\ne0$,
\[
\hat Q(r)=\frac12\sum_{x\ne0}e\Bigl(\frac{rx}p\Bigr)+\frac12\sum_x\chi(x)e\Bigl(\frac{rx}p\Bigr)=\frac{-1+\chi(r)g}2,
\]
where $g=\sum_x\chi(x)e(\frac xp)$ is the Gauss sum (substitute $x\mapsto\bar rx$), and $|g|=\sqrt p$. Hence $|\hat Q(r)|\le\frac{\sqrt p+1}2$, and (b) with
$|Q|=\frac{p-1}2$ gives the claim.

(d) By (b), $T(A)\ge\delta^3p^2-\delta p\cdot\frac12\delta^2p=\frac12\delta^3p^2$. The pairs with $d=0$ number $|A|=\delta p$, and $\frac12\delta^3p^2>\delta p$ because
$\delta^2p>2$. So some pair has $d\ne0$.

\tools{Volume~5, Theorem~9.3 and Proposition~9.13 (Parseval, the Gauss sum); Volume~2, Chapter~6 (the Legendre symbol).}
\leads Part~(d) is the first half of Roth's proof (1953) that a subset of $\{1,\dots,N\}$ of density $\delta$ contains a three-term
progression once $N$ is large: either all non-trivial coefficients are small, and (d) finds a progression, or one is large, and
the set is denser on a long subprogression, where one starts again. Meshulam's version over $\mathbb F_3^n$ is in Volume~8. Behrend (1946)
built progression-free sets of size $Ne^{-c\sqrt{\log N}}$, and Kelley and Meka (2023) proved a nearly matching upper bound of the
shape $Ne^{-c(\log N)^\beta}$.
\end{solution}

\begin{solution}{pb:dftspectrum}
\idea $F^4=I$ pins the spectrum to four points, and two traces, one of which is Gauss's famous sum, pin the multiplicities.

\writeup (a) By orthogonality,
\[
(F^2x)_j=\frac1n\sum_le\Bigl(-\frac{jl}n\Bigr)\sum_kx_ke\Bigl(-\frac{lk}n\Bigr)=\sum_kx_k\cdot\frac1n\sum_le\Bigl(-\frac{l(j+k)}n\Bigr)=x_{-j}.
\]
So $F^2=R$, the reversal, and $F^4=R^2=I$. The polynomial $X^4-1$ kills $F$ and has four simple roots, so $F$ is diagonalisable (its
minimal polynomial divides $X^4-1$, hence has simple roots), with eigenvalues among $1,-1,-i,i$.

(b) In an eigenbasis,
\[
\operatorname{tr}F=m_1-m_{-1}-im_{-i}+im_i,\qquad\operatorname{tr}F^2=m_1+m_{-1}-m_{-i}-m_i.
\]
Now $\operatorname{tr}F^2=\operatorname{tr}R$ is the number of $k\in\Z/n\Z$ with $k=-k$: $1$ if $n$ is odd and $2$ if $n$ is even. And
$\operatorname{tr}F=\frac1{\sqrt n}\sum_ke(-\frac{k^2}n)=\frac{\overline{G(n)}}{\sqrt n}$, where $G(n)=\sum_{k=0}^{n-1}e(\frac{k^2}n)$; by Gauss's theorem $\operatorname{tr}F$ is $1-i$, $1$, $0$, $-i$ for
$n\equiv0,1,2,3\pmod4$.

(c) Let $\sigma=m_1+m_{-1}$ and $\tau=m_{-i}+m_i$. Then $\sigma+\tau=n$ and $\sigma-\tau=\operatorname{tr}F^2$; and $m_1-m_{-1}=\operatorname{Re}\operatorname{tr}F$, $m_i-m_{-i}=\operatorname{Im}\operatorname{tr}F$. For
$n=4m$: $\sigma=2m+1$, $\tau=2m-1$, $m_1-m_{-1}=1$, $m_i-m_{-i}=-1$, so $(m_1,m_{-1},m_{-i},m_i)=(m+1,m,m,m-1)$. For $n=4m+1$: $\sigma=2m+1$, $\tau=2m$,
differences $1$ and $0$: $(m+1,m,m,m)$. For $n=4m+2$: $\sigma=2m+2$, $\tau=2m$, differences $0$, $0$: $(m+1,m+1,m,m)$. For $n=4m+3$: $\sigma=2m+2$,
$\tau=2m+1$, differences $0$ and $-1$: $(m+1,m+1,m+1,m)$.

\tools{Volume~5, Theorem~9.3 (DFT inversion) and Theorem~8.14 (Gauss's evaluation); Volume~1, Chapter~7 (diagonalisation).}
\leads The argument can be run backwards: Schur (1921) found the multiplicities independently, from the determinant of the DFT
matrix (a Vandermonde determinant), and deduced the value of $G(n)$, sign included, which had cost Gauss years of effort. On $L^2(\R)$ the
Fourier transform also satisfies $\mathcal F^4=I$; its eigenfunctions are the Hermite functions $H_k(\sqrt{2\pi}\,x)e^{-\pi x^2}$, with eigenvalue $(-i)^k$, and
McClellan and Parks (1972) studied discrete analogues of them for the DFT.
\end{solution}

\begin{solution}{pb:erdosturan}
\idea Indicators of intervals have badly decaying Fourier coefficients, so we sandwich them between continuous trapezoids, whose
coefficients decay like $h^{-2}$. The price is the width $\delta$ of the slopes; the gain is an absolutely convergent expansion that
turns counting points into exponential sums.

\writeup (a) Write $k=\frac1\delta1_{[0,\delta]}$, extended $1$-periodically like $1_J$. Then $t(x)=\frac1\delta\int_0^\delta1_J(x-u)\dd u$ is an average of values in
$\{0,1\}$, so $0\le t\le1$, and by Fubini $\int_0^1t=\int_0^11_J=L$. The Fourier coefficients of a periodic convolution multiply (Fubini
again): $\hat t(h)=\hat1_J(h)\hat k(h)$. For $h\ne0$,
\[
|\hat1_J(h)|=\Bigl|\int_Je(-hx)\dd x\Bigr|=\frac{|\sin\pi hL|}{\pi|h|}\le\frac1{\pi|h|},\qquad|\hat k(h)|=\frac{|\sin\pi h\delta|}{\pi|h|\delta}\le\min\Bigl(1,\frac1{\pi|h|\delta}\Bigr),
\]
and the bound follows.

(b) By (a), $\sum_h|\hat t(h)|<\infty$, so $\sum_h\hat t(h)e(hx)$ converges uniformly to a continuous function with the same coefficients as $t$, which is
$t$ itself (Volume~5, Corollary~8.10(b), after the change of variable $x\mapsto\frac x{2\pi}$). Hence
\[
\frac1N\sum_nt(x_n)=L+\frac1N\sum_{h\ne0}\hat t(h)S(h),
\]
and, as $S(-h)=\overline{S(h)}$, the error is at most
\[
E=\frac2N\sum_{h\ge1}\min\Bigl(\frac1{\pi h},\frac1{\pi^2h^2\delta}\Bigr)|S(h)|.
\]
Let $I$ have ends $a<b$ and length $\ell$,
and let $\nu=\frac1N\#\{n:x_n\in I\}$.

Upper bound. If $\ell\le1-2\delta$, take $J=[a-\delta,b]$, of length $L=\ell+\delta\le1-\delta$. For $x\in I$, $[x-\delta,x]\subseteq J$, so $t(x)=1$: thus $t\ge1_I$
and $\nu\le\frac1N\sum t(x_n)\le\ell+\delta+E$. If $\ell>1-2\delta$, simply $\nu\le1<\ell+2\delta$.

Lower bound. If $\ell\ge\delta$, take $J=[a,b-\delta]$, of length $\ell-\delta$. If $t(x)>0$ then $x-u\in J$ for a set of $u\in[0,\delta]$ of positive
length, which forces $a<x<b$; so $t\le1_I$ and $\nu\ge\ell-\delta-E$. If $\ell<\delta$, simply $\nu\ge0>\ell-\delta$.

In all cases $|\nu-\ell|\le2\delta+E$, which is the claim.

(c) For $p\nmid h$, $S(h)=\sum_{k=0}^{p-1}e(\frac{hk^2}p)=\sum_y(1+\chi(y))e(\frac{hy}p)=\chi(h)g$, where $\chi$ is the Legendre symbol and $g$ the Gauss sum; so
$|S(h)|=\sqrt p$. For $p\mid h$, $S(h)=p$. We use two elementary bounds, for $X\ge1$: $\sum_{h\le X}\frac1h\le1+\log X$, and
$\sum_{h>X}\frac1{h^2}\le\sum_{h>\lfloor X\rfloor}\frac1{h(h-1)}=\frac1{\lfloor X\rfloor}\le\frac2X$. Take $\delta=p^{-1/2}$. The terms with $p\nmid h$ contribute at most
\[
\frac{2\sqrt p}p\Bigl(\sum_{h\le\sqrt p}\frac1{\pi h}+\sum_{h>\sqrt p}\frac1{\pi^2h^2\delta}\Bigr)\le\frac2{\sqrt p}\Bigl(\frac{1+\frac12\log p}\pi+\frac2{\pi^2}\Bigr),
\]
and those with $h=mp$ at most $\frac2p\sum_{m\ge1}\frac p{\pi^2m^2p^2\delta}=\frac1{3p^{3/2}}$. Altogether
\[
D\le\frac{\log p+2+2\pi+\frac4\pi}{\pi\sqrt p}+\frac1{3p^{3/2}}\le\frac{\log p+10}{\pi\sqrt p},
\]
since $2+2\pi+\frac4\pi=9.55\ldots$ and $\frac1{3p^{3/2}}\le\frac{0.44}{\pi\sqrt p}$ for $p\ge3$.

\tools{Volume~5, Corollary~8.10 and Theorem~8.13 (Fejér, Weyl's criterion), Proposition~7.13 (Gauss sums).}
\leads Erdős and Turán (1948) proved the general form $D\le\frac C H+C\sum_{h\le H}\frac{|S(h)|}{hN}$ for every $H$; Vaaler (1985) found the optimal
smoothing functions, which give it with small explicit constants. For the squares modulo $p$, one finds numerically $D\sqrt p/\log p\approx0.14$
at $p=10007$, comfortably inside the bound of (c).
\end{solution}

\begin{solution}{pb:circlewalk}
\idea The Fourier coefficients of the law of $X_n$ are $n$-th powers $\cos^n(2\pi h\alpha)$. They tend to $0$ unless the cosine is $\pm1$, and
which frequencies survive is a question about $2h\alpha$ being an integer.

\writeup Since $e(hX_n)=e(h\alpha(\varepsilon_1+\cdots+\varepsilon_n))$ and the signs are independent,
\[
\mathbb E[e(hX_n)]=\prod_{k=1}^n\mathbb E[e(h\alpha\varepsilon_k)]=\cos^n(2\pi h\alpha),
\]
and $|\cos(2\pi h\alpha)|=1$ exactly when $2h\alpha\in\Z$.

(a) If $\alpha$ is irrational and $h\ne0$, then $2h\alpha\notin\Z$, so $\mathbb E[e(hX_n)]\to0=\int_0^1e(hx)\dd x$. By linearity $\mathbb E[P(X_n)]\to\int_0^1P$ for every
trigonometric polynomial $P$. Given a continuous $1$-periodic $g$ and $\varepsilon>0$, Fejér's theorem (Volume~5, Corollary~8.10(a)) gives a
trigonometric polynomial $P$ with $|g-P|\le\varepsilon$ everywhere, so $\limsup_n|\mathbb E[g(X_n)]-\int_0^1g|\le2\varepsilon$.

(b) Now $X_n$ lives in the finite set $\frac1q\Z/\Z$, and the law of a variable on this set is determined by the $q$ numbers $\mathbb E[e(hX)]$,
$0\le h<q$ (inversion of the DFT, Volume~5, Theorem~9.3); so it suffices to find the limits of these numbers. As $\gcd(a,q)=1$, $2h\alpha\in\Z$ iff
$q\mid2h$.

If $q$ is odd, this means $q\mid h$, and then $\cos(2\pi h\alpha)=1$; so $\mathbb E[e(hX_n)]\to1_{q\mid h}=\frac1q\sum_{j<q}e(\frac{hj}q)$, the coefficient of the uniform law on
$\frac1q\Z/\Z$.

If $q$ is even, $a$ is odd, and $q\mid2h$ means $h=l\frac q2$; then $\cos(2\pi h\alpha)=\cos(\pi la)=(-1)^l$. So $\mathbb E[e(hX_n)]$ tends to $0$ unless $h=l\frac q2$,
where it equals $(-1)^{ln}$. For $n$ even these are the coefficients of the uniform law on $\{0,\frac2q,\frac4q,\dots\}$, and for $n$ odd those of the
uniform law on $\{\frac1q,\frac3q,\dots\}$: indeed $\frac2q\sum_je(h(\frac1q+\frac{2j}q))=e(\frac hq)1_{\frac q2\mid h}=(-1)^l$ for $h=l\frac q2$. The two sets are disjoint, so
$X_n$ has no limit; this is the parity obstruction of Problem~\ref{pb:cyclewalk}, where $\alpha=\frac1k$.

\tools{Volume~5, Corollary~8.10 and Theorem~9.3; Volume~4, Chapter~5 (independence).}
\leads Compare with Weyl's theorem that $(\{n\alpha\})$ is equidistributed (Volume~5, Corollary~8.14): the deterministic walk and the random walk
are both governed by $\alpha$ being irrational. The \emph{speed} of convergence in (a), $\cos^n(2\pi h\alpha)\approx e^{-2\pi^2n\|2h\alpha\|^2}$ where $\|\cdot\|$ is the
distance to the nearest integer, depends on how well $2\alpha$ is approximated by rationals, that is, on its continued fraction (Volume~2,
Chapter~8).
\end{solution}

\begin{solution}{pb:fejer}
\idea The polynomial $z^nT(z)$ is self-reciprocal, so its roots come in pairs $\alpha$, $1/\bar\alpha$; on the circle the two members of a pair
have proportional distances to $z$, so keeping one member of each pair produces a square root of $T$. Once $T=|P|^2$, extremal problems
for $T$ become quadratic-form problems for the coefficients of $P$.

\writeup (a) If $T$ is constant, take $P=\sqrt{c_0}$. Otherwise reduce $n$ until $c_n\ne0$, and let $Q(z)=\sum_{k=-n}^nc_kz^{k+n}$, of degree $2n$ with $Q(0)=c_{-n}=\overline{c_n}\ne0$.
Since $c_{-k}=\overline{c_k}$, $z^{2n}\overline{Q(1/\bar z)}=Q(z)$ for $z\ne0$, so $\alpha\mapsto1/\bar\alpha$ permutes the roots of $Q$, preserving multiplicities. A root
$e^{i\theta_0}$ on the circle has even multiplicity: $T(\theta)=e^{-in\theta}Q(e^{i\theta})$ vanishes at $\theta_0$ to the order of that multiplicity and does not change
sign. So the roots can be listed as $\alpha_1,\dots,\alpha_n,1/\bar\alpha_1,\dots,1/\bar\alpha_n$ with $0<|\alpha_j|\le1$, and $Q(z)=c_n\prod_j(z-\alpha_j)(z-1/\bar\alpha_j)$. For $|z|=1$,
$|z-1/\bar\alpha|=\frac{|\bar\alpha z-1|}{|\alpha|}=\frac{|z|\,|\bar\alpha-\bar z|}{|\alpha|}=\frac{|z-\alpha|}{|\alpha|}$. Hence, as $T\ge0$,
\[
T(\theta)=|Q(e^{i\theta})|=\frac{|c_n|}{\prod_j|\alpha_j|}\prod_j|e^{i\theta}-\alpha_j|^2=|P(e^{i\theta})|^2,\qquad P(z)=\Bigl(\frac{|c_n|}{\prod_j|\alpha_j|}\Bigr)^{1/2}\prod_j(z-\alpha_j).
\]

(b) Write $P=\sum_{k=0}^np_kz^k$. By orthogonality, $c_0=\frac1{2\pi}\int|P(e^{i\theta})|^2\dd\theta=\sum|p_k|^2=1$, and by Cauchy--Schwarz
$T(0)=|\sum p_k|^2\le(n+1)\sum|p_k|^2=n+1$, with equality iff all $p_k$ are equal, that is $T=\frac1{n+1}|\sum_{k=0}^ne^{ik\theta}|^2$, the Fejér kernel $F_{n+1}$ of
Volume~5, Lemma~8.8.

(c) Now $c_k=\frac{a_k}2$ for $k\ge1$ and $c_0=1$. With $T=|P|^2$ as in (b), $c_1=\frac1{2\pi}\int|P|^2e^{-i\theta}\dd\theta=\sum_kp_{k+1}\overline{p_k}$, so with $u_k=|p_k|$,
\[
a_1=2\operatorname{Re}\sum_{k=0}^{n-1}p_{k+1}\overline{p_k}\le2\sum_{k=0}^{n-1}u_ku_{k+1}=u^{\mathsf T}Au,\qquad|u|=1,
\]
where $A$ is the adjacency matrix of the path with $n+1$ vertices. So $a_1\le\lambda_{\max}(A)$. The vectors $v^{(j)}_k=\sin\frac{(k+1)j\pi}{n+2}$, $0\le k\le n$, satisfy
$Av^{(j)}=2\cos\frac{j\pi}{n+2}v^{(j)}$ (by $\sin(x-y)+\sin(x+y)=2\sin x\cos y$, with $v_{-1}=v_{n+1}=0$), for $j=1,\dots,n+1$: these are $n+1$ distinct eigenvalues,
hence all of them, and the largest is $2\cos\frac\pi{n+2}$. For equality take $p_k=v^{(1)}_k/|v^{(1)}|>0$: then $T=|P|^2$ has real coefficients $c_k=c_{-k}$, so it is
a cosine polynomial $1+\sum a_k\cos k\theta$, and $a_1=2\sum p_kp_{k+1}=2\cos\frac\pi{n+2}$.

\tools{Volume~1, Chapter~6 (roots of polynomials) and Chapter~7 (symmetric matrices); Volume~5, Lemma~8.8 (the Fejér kernel).}
\leads Non-negative trigonometric polynomials are everywhere. Mertens's $3+4\cos\theta+\cos2\theta=2(1+\cos\theta)^2\ge0$ is the heart of the zero-free
region for $\zeta$ on the line $\operatorname{Re}s=1$ (Volume~6); Delsarte's linear programming bound for codes and the Cohn--Elkies bounds for sphere packings,
which Viazovska used in dimension $8$, optimise over non-negative or positive-definite functions; and engineers use the Fejér--Riesz
factorisation, under the name \emph{spectral factorisation}, to design filters.
\end{solution}

\begin{solution}{pb:bernstein}
\idea A trigonometric polynomial of degree $n$ has at most $2n$ zeros in a period. If $T$ climbed faster than the pure oscillation
$\cos(n(\theta-\beta))$ that has the same value at some point, their difference would have too many zeros.

\writeup (a) First suppose $\|T\|<1$, and suppose that $T'(\theta_0)^2+n^2T(\theta_0)^2>n^2$ for some $\theta_0$. Then $T'(\theta_0)\ne0$, as $|T(\theta_0)|<1$. Let $\gamma\in\left]0,\pi\right[$ with
$\cos\gamma=T(\theta_0)$, and choose $\beta$ with $n(\theta_0-\beta)=\pm\gamma$, the sign chosen so that $C(\theta)=\cos(n(\theta-\beta))$ has $C(\theta_0)=T(\theta_0)$ and $C'(\theta_0)$ of the
same sign as $T'(\theta_0)$; then $|C'(\theta_0)|=n\sin\gamma=n\sqrt{1-T(\theta_0)^2}<|T'(\theta_0)|$. The difference $D=C-T$ is a trigonometric polynomial of degree at most
$n$, and not $0$ since $C$ reaches $1$. At the $2n$ points $\theta_j=\beta+\frac{j\pi}n$ of a period, $C=(-1)^j$, so $D(\theta_j)$ has the sign of $(-1)^j$, and $D$ vanishes in
each of the $2n$ arcs $]\theta_j,\theta_{j+1}[$. Now look at the arc containing $\theta_0$, where $C$ is monotone; say $C$ decreases there, the other case
being symmetric. Then $D(\theta_j)>0>D(\theta_{j+1})$, and $T'(\theta_0)<C'(\theta_0)<0$, so $D'(\theta_0)>0$: $D$ is negative just before $\theta_0$ and positive just after. So $D$
has a zero in $]\theta_j,\theta_0[$, the zero $\theta_0$, and a zero in $]\theta_0,\theta_{j+1}[$: at least $2n+2$ zeros in a period, which is impossible. Hence
$T'^2+n^2T^2\le n^2$ when $\|T\|<1$, and by applying this to $(1-\varepsilon)T$ and letting $\varepsilon\to0$, also when $\|T\|\le1$. In particular $|T'|\le n$, and by
homogeneity $\|T'\|\le n\|T\|$. For complex $T$ and a given $\theta_0$, choose $\beta$ with $e^{i\beta}T'(\theta_0)=|T'(\theta_0)|$; the real polynomial $U=\operatorname{Re}(e^{i\beta}T)$ has
$\|U\|\le\|T\|$ and $U'(\theta_0)=|T'(\theta_0)|$, so $|T'(\theta_0)|\le n\|T\|$.

(b) $T(\theta)=P(\cos\theta)$ is a real trigonometric polynomial of degree $n$ with $\|T\|\le1$, and $T'(\theta)=-\sin\theta\,P'(\cos\theta)$. For $x=\cos\theta$ with
$0<\theta<\pi$, (a) gives $|P'(x)|\sin\theta\le n$, that is $|P'(x)|\le\frac n{\sqrt{1-x^2}}$.

(c) $T=\sin n\theta$ has $\|T\|=1$ and $T'(0)=n$. For the Chebyshev polynomial, $T_n(\cos\theta)=\cos n\theta$, so $T_n'(\cos\theta)=\frac{n\sin n\theta}{\sin\theta}$, and at $\theta=\frac\pi{2n}$,
$|T_n'(x)|=\frac n{\sqrt{1-x^2}}$.

\tools{Volume~1, Chapter~5 (trigonometric polynomials) and Chapter~3 (counting roots); Volume~5, Corollary~8.10 (Chebyshev polynomials).}
\leads Bernstein proved the inequality in 1912 to show that a function which trigonometric polynomials
of degree $n$ approximate within $n^{-\alpha}$ must be Hölder of order $\alpha$, for $0<\alpha<1$: the converse of Jackson's theorem. The sharper form
in (a) is the Bernstein--Szegő inequality. Near the ends of $[-1,1]$ the bound of (b) blows up, and the truth there is Markov's
inequality (1889): $|P'|\le n^2$ on $[-1,1]$, with equality for $T_n$ at $\pm1$. Problem~\ref{pb:bernsteinabs} uses (a) to build a function.
\end{solution}

\begin{solution}{pb:bernsteinabs}
\idea Hölder continuity controls the $L^2$ norm of differences, and Parseval converts this into $\ell^2$ control of each dyadic block of
coefficients. Cauchy--Schwarz turns $\ell^2$ into $\ell^1$ at the price of the square root of the block length; the books balance exactly
when $\alpha=\frac12$, and Rudin--Shapiro polynomials show that at $\alpha=\frac12$ they do not balance in our favour.

\writeup (a) $f$ is continuous, and $g(x)=f(x+h)-f(x-h)$ has coefficients $(e^{ikh}-e^{-ikh})c_k=2i\sin(kh)c_k$; apply Parseval (Volume~5, Theorem~8.11).

(b) The left side of (a) is at most $C^2(2h)^{2\alpha}$. For $j\ge0$ let $B_j=\{k:2^{j-1}<|k|\le2^j\}$, which has $2^j$ elements, and $h=\frac\pi{3\cdot2^j}$; for
$k\in B_j$, $\frac\pi6<|k|h\le\frac\pi3$, so $\sin^2(kh)\ge\frac14$ and
\[
\sum_{k\in B_j}|c_k|^2\le\sum_k4\sin^2(kh)|c_k|^2\le C^2\Bigl(\frac{2\pi}{3\cdot2^j}\Bigr)^{2\alpha}.
\]
By Cauchy--Schwarz,
\[
\sum_{k\in B_j}|c_k|\le2^{j/2}C\Bigl(\frac{2\pi}3\Bigr)^\alpha2^{-j\alpha}=C\Bigl(\frac{2\pi}3\Bigr)^\alpha2^{j(1/2-\alpha)},
\]
a geometric series in $j$ since $\alpha>\frac12$. So
$\sum_k|c_k|=|c_0|+\sum_j\sum_{B_j}|c_k|<\infty$. Then $\sum_kc_ke^{ikx}$ converges uniformly to a continuous function with the coefficients of $f$, which is $f$
(Volume~5, Corollary~8.10(b)).

(c) Let $g_j(x)=e^{i2^jx}P_j(e^{ix})$. By Problem~\ref{pb:rudin}, its frequencies are $2^j,\dots,2^{j+1}-1$, its coefficients are $\pm1$, and $\|g_j\|\le\sqrt{2^{j+1}}$. So
$f=\sum_j2^{-j}g_j$ converges uniformly, and integrating term by term, $c_k(f)=\pm2^{-j}$ for $2^j\le k<2^{j+1}$: $\sum_k|c_k(f)|=\sum_j2^{-j}2^j=\infty$. For the Hölder
bound, let $0<|h|\le1$ and $m\ge0$ with $2^{-m-1}<|h|\le2^{-m}$. As $g_j$ has degree less than $2^{j+1}$, Bernstein's inequality (Problem~\ref{pb:bernstein})
gives $\|g_j'\|\le2^{j+1}\sqrt{2^{j+1}}$. Hence
\[
\begin{aligned}|f(x+h)-f(x)|&\le\sum_{j\le m}2^{-j}|h|\,2^{3(j+1)/2}+\sum_{j>m}2^{-j}\cdot2\sqrt{2^{j+1}}\\&\le\frac{2^{3/2}|h|2^{m/2}+2^{3/2}2^{-(m+1)/2}}{1-2^{-1/2}}\le20|h|^{1/2},\end{aligned}
\]
since $2^{m/2}\le|h|^{-1/2}$ and $2^{-(m+1)/2}<|h|^{1/2}$. For $|h|>1$, $|f(x+h)-f(x)|\le2\|f\|\le2\|f\||h|^{1/2}$. So $f$ is Hölder of order $\frac12$.

\tools{Volume~5, Theorem~8.11 and Corollary~8.10; Problems~\ref{pb:rudin} and~\ref{pb:bernstein}.}
\leads The space of absolutely convergent Fourier series is an algebra, the Wiener algebra of Volume~8, whose invertible elements Wiener
characterised in 1932. Bernstein's theorem says that it contains every Hölder function of order $>\frac12$; (c), which is essentially
Rudin's motivation for his 1959 note, shows that order $\frac12$ is not enough. Hardy and Littlewood had shown this earlier with the series
$\sum_k\frac{e^{ick\log k}}ke^{ikx}$.
\end{solution}

\begin{solution}{pb:fantaussky}
\idea Write the polygon as a sum of regular polygons, one for each frequency $m$ (its discrete Fourier modes). Differences and areas
are both diagonal in this decomposition, and each mode is compared with its own constant; the extremal mode is $m=1$, the regular
polygon traversed once.

\writeup (a) Write $x_k=\sum_{m=0}^{n-1}c_m\omega^{mk}$ (DFT inversion, Volume~5, Theorem~9.3); $c_0=\frac1n\sum x_k=0$. Then $x_{k+1}-x_k=\sum_mc_m(\omega^m-1)\omega^{mk}$, and by
Parseval,
\[
\begin{aligned}\sum_k|x_{k+1}-x_k|^2&=n\sum_m|c_m|^2|\omega^m-1|^2=n\sum_{m=1}^{n-1}4\sin^2\frac{\pi m}n|c_m|^2\\&\ge4\sin^2\frac\pi n\cdot n\sum_m|c_m|^2=4\sin^2\frac\pi n\sum_k|x_k|^2,\end{aligned}
\]
since $\sin^2\frac{\pi m}n\ge\sin^2\frac\pi n$ for $1\le m\le n-1$, with equality only for $m=1$ and $m=n-1$. Equality holds iff $c_m=0$ for the other $m$, that is
$x_k=a\omega^k+b\omega^{-k}$.

(b) With $z_k=\sum_mc_m\omega^{mk}$, orthogonality gives $\sum_k\overline{z_k}z_{k+1}=n\sum_m|c_m|^2\omega^m$, so $2A=n\sum_m|c_m|^2\sin\frac{2\pi m}n$, while
$\sum_k|z_{k+1}-z_k|^2=n\sum_m|c_m|^2\,4\sin^2\frac{\pi m}n$. With $\theta_m=\frac{\pi m}n$,
\[
\sum_k|z_{k+1}-z_k|^2-4\tan\frac\pi n\cdot A=4n\sum_m|c_m|^2\sin\theta_m\Bigl(\sin\theta_m-\tan\frac\pi n\cos\theta_m\Bigr).
\]
The term $m=0$ vanishes. For $1\le m\le n-1$, $\sin\theta_m>0$; the bracket is positive if $\cos\theta_m\le0$, and equals $\cos\theta_m(\tan\theta_m-\tan\frac\pi n)\ge0$ if
$\cos\theta_m>0$, with equality only for $m=1$. So the difference is $\ge0$, with equality iff $z_k=c_0+c_1\omega^k$: a positively oriented regular $n$-gon
with centre $c_0$, or a point if $c_1=0$.

(c) If all sides have length $s$, then $L=ns$ and $\sum|z_{k+1}-z_k|^2=ns^2=\frac{L^2}n$, so (b) gives $A\le\frac{L^2}{4n\tan(\pi/n)}$, with equality for the regular
$n$-gon (orient the polygon positively, so that $A$ is its area).

\tools{Volume~5, Theorem~9.3 (DFT inversion and Parseval); Volume~1, Chapter~5 (roots of unity).}
\leads Ky Fan, Olga Taussky and John Todd (1955) proved (a) as a discrete analogue of Wirtinger's inequality, and (b) is the discrete
version of Hurwitz's Fourier proof of the isoperimetric inequality (Volume~5). The full polygonal isoperimetric inequality,
$L^2\ge4n\tan\frac\pi n\cdot A$ for every $n$-gon, needs a further argument (equalising the sides); the decomposition of a polygon into regular
modes is also the key to Problems~\ref{pb:midpoint} and~\ref{pb:napoleon}.
\end{solution}

\begin{solution}{pb:involasym}
\idea On the circle $|z|=r$ where $|T(z)|r^{-n}$ is smallest along the positive axis, the integrand is a sharp Gaussian bump around $\theta=0$. Laplace's method
does the rest, and Stirling turns the result into a closed form.

\writeup (a) $T$ is entire, so Cauchy's formula on $|z|=r$ (Volume~5, Theorem~11.10) gives $\frac{t_n}{n!}=\frac1{2\pi i}\oint\frac{T(z)}{z^{n+1}}\dd z$; parametrise $z=re^{i\theta}$.

(b) $|T(re^{i\theta})|=\exp\operatorname{Re}\bigl(re^{i\theta}+\frac12r^2e^{2i\theta}\bigr)=\exp\bigl(r\cos\theta+\frac12r^2\cos2\theta\bigr)$, and
$r(\cos\theta-1)+\frac12r^2(\cos2\theta-1)=-2r\sin^2\frac\theta2-r^2\sin^2\theta$.

(c) Let $\delta=r^{-4/5}$. \emph{Tails.} For $\delta\le|\theta|\le\frac\pi2$, $\sin^2\theta\ge\frac{4\theta^2}{\pi^2}$, so by (b) the integrand has modulus at most
$T(r)r^{-n}e^{-4r^2\delta^2/\pi^2}=T(r)r^{-n}e^{-4r^{2/5}/\pi^2}$. For $\frac\pi2\le|\theta|\le\pi$, $\sin^2\frac\theta2\ge\frac12$, so the modulus is at most $T(r)r^{-n}e^{-r}$. Both are
$o\bigl(T(r)r^{-n}/\sqrt b\bigr)$, since $b\le3r^2$.

\emph{Centre.} Let $G(\theta)=re^{i\theta}+\frac12r^2e^{2i\theta}-in\theta$. Then $G(0)=\ln T(r)$, $G'(0)=i(r+r^2)-in=0$, $G''(0)=-(r+2r^2)=-b$, and $|G'''|\le r+4r^2\le5r^2$. By Taylor's formula,
$G(\theta)=\ln T(r)-\frac b2\theta^2+E(\theta)$ with $|E(\theta)|\le r^2|\theta|^3\le r^{-2/5}$ for $|\theta|\le\delta$. Hence
\[
\int_{-\delta}^\delta e^{G(\theta)}\dd\theta=T(r)\bigl(1+O(r^{-2/5})\bigr)\int_{-\delta}^\delta e^{-b\theta^2/2}\dd\theta=T(r)\sqrt{\frac{2\pi}b}\,\bigl(1+o(1)\bigr),
\]
because $b\delta^2\ge2r^{2/5}\to\infty$, so the Gaussian integral over $|\theta|>\delta$ is negligible. Dividing by $2\pi r^n$ gives the claim.

(d) $r=\frac{\sqrt{4n+1}-1}2=\sqrt n-\frac12+\frac1{8\sqrt n}+O(n^{-3/2})$. Since $r^2=n-r$, $\ln T(r)=r+\frac{r^2}2=\frac{n+r}2$. Next $\frac r{\sqrt n}=1+\varepsilon$ with
$\varepsilon=-\frac1{2\sqrt n}+\frac1{8n}+O(n^{-3/2})$, and $\ln(1+\varepsilon)=\varepsilon-\frac{\varepsilon^2}2+O(\varepsilon^3)=-\frac1{2\sqrt n}+O(n^{-3/2})$, so $n\ln r=\frac n2\ln n-\frac{\sqrt n}2+O(n^{-1/2})$. Hence
\[
\ln\frac{T(r)}{r^n}=\frac n2+\frac{\sqrt n}2-\frac14-\frac n2\ln n+\frac{\sqrt n}2+o(1)=\frac n2+\sqrt n-\frac14-\frac n2\ln n+o(1).
\]
With $b\sim2n$ and $n!\sim\sqrt{2\pi n}\,n^ne^{-n}$,
\[
t_n\sim\sqrt{2\pi n}\,n^ne^{-n}\cdot\frac{e^{n/2+\sqrt n-1/4}n^{-n/2}}{\sqrt{4\pi n}}=\frac1{\sqrt2}n^{n/2}e^{-n/2+\sqrt n-1/4}.
\]

\tools{Volume~5, Theorem~11.10 (Cauchy's formula), Theorem~11.9 (Laplace's method) and the example of the central binomial coefficient in Chapter~11; Volume~1, Theorem~3.17 (Stirling).}
\leads This is the saddle-point method of Volume~5 applied to a function that is not a power, the setting that Hayman (1956) axiomatised; the result for
involutions is due to Chowla, Herstein and Moore (1951) and to Moser and Wyman (1955). Numerical evaluation shows that the ratio of $t_n$ to the approximation
in (d) is $1.0091$ for $n=1000$: the error is of relative order $\frac1{\sqrt n}$, which is why the approximation in (c), with ratio $0.99983$ at $n=1000$, is much better.
\end{solution}

\begin{solution}{pb:distinctcycles}
\idea The generating function is $\frac1{1-z}$ times a function that is tame at $z=1$; the partial sums of its coefficients then converge to its value at $1$, which
telescopes into $e^{-\gamma}$.

\writeup (a) A permutation is a set of cycles, and there are $(k-1)!$ cycles on a given $k$-set, so the exponential generating function of one cycle of length $k$
is $\frac{z^k}k$. Allowing each length at most once, the exponential formula (Volume~5, Theorem~2.9) gives $\sum_n(\text{number}/n!)z^n=\prod_k(1+\frac{z^k}k)$, and the number divided
by $n!$ is $d_n$. Expanding: $d_1=1$, $d_2=\frac12$ (only $1+\frac{z^2}2$ contributes, not $z\cdot z$), $d_3=\frac13+\frac12=\frac56$, $d_4=\frac14+\frac13=\frac7{12}$, $d_5=\frac15+\frac14+\frac16=\frac{37}{60}$.

(b) For $|z|<1$, $\prod_ke^{z^k/k}=e^{-\ln(1-z)}=\frac1{1-z}$, and all products converge absolutely, so $\prod(1+\frac{z^k}k)=\frac{h(z)}{1-z}$. At $z=1$,
$\prod_{k\le N}(1+\frac1k)e^{-1/k}=(N+1)e^{-H_N}\to e^{-\gamma}$.

(c) $\|fg\|\le\|f\|\,\|g\|$ is the triangle inequality applied to the Cauchy product. Next, $(1+w)e^{-w}=\sum_j\frac{(-w)^j}{j!}+\sum_j\frac{(-1)^jw^{j+1}}{j!}=1-\sum_{j\ge2}(-1)^j\frac{j-1}{j!}w^j$,
so with $w=\frac{z^k}k$,
\[
\Bigl\|\Bigl(1+\frac{z^k}k\Bigr)e^{-z^k/k}-1\Bigr\|\le\sum_{j\ge2}\frac{j-1}{j!}k^{-j}\le\frac1{k^2}\sum_{j\ge2}\frac{j-1}{j!}=\frac1{k^2}.
\]
Let $P_N$ be the product of the first $N$ factors. Then $\|P_N\|\le\prod_k(1+k^{-2})=:C<\infty$, and $\|P_{N+1}-P_N\|\le\|P_N\|\cdot(N+1)^{-2}\le C(N+1)^{-2}$, so $(P_N)$ is a Cauchy
sequence for $\|\cdot\|$ and converges to a series $\tilde h$ with $\|\tilde h\|<\infty$. Convergence in this norm implies convergence of each coefficient and uniform convergence on
$|z|\le1$, so $\tilde h=h$ on $|z|<1$ and $\tilde h(1)=\lim P_N(1)=h(1)$. Hence $\sum_j|h_j|<\infty$ and $\sum_jh_j=h(1)$.

(d) From $\sum d_nz^n=h(z)\sum z^n$, $d_n=h_0+\dots+h_n\to\sum_jh_j=e^{-\gamma}$.

\tools{Volume~5, Theorem~2.9 (exponential formula) and Corollary~11.6 ($H_N-\ln N\to\gamma$); Volume~1, Chapter~3 (absolute convergence, products).}
\leads The convergence is slow: $d_{300}=0.5633$, with an error of order $\frac1n$. Compare with P1, Chapter~4: a permutation with no cycle shorter than $m$ has probability
about $\frac{e^{-\gamma}}m$, and an integer without prime factors below $m$ has density about $\frac{e^{-\gamma}}{\ln m}$ (Mertens). The constant $e^{-\gamma}$ keeps appearing where
cycles or primes are forbidden to repeat or to be small. The norm $\|\cdot\|$ is that of the Wiener algebra, which reappears with Wiener's theorem in Volume~8.
\end{solution}

\begin{solution}{pb:ramanujan}
\idea Both $Q$ and $R$ are Laplace integrals of the same function, on the two sides of its peak. The peak is not symmetric: the skewness of $s^{n-1}e^{n(1-s)}$, seen through
$g(v)=1-\frac v3+\cdots$, gives $\mp\frac13$, and Ramanujan's $\frac13$ is half of $\frac13+\frac13$.

\writeup (a) $\frac{n^{n-j}}{(n-j)!}=\frac{n^n}{n!}\cdot\frac{n(n-1)\cdots(n-j+1)}{n^j}$, and summing over $j=1,\dots,n$ gives the first identity. Likewise $\frac{n^{n+j}}{(n+j)!}=\frac{n^n}{n!}\cdot\frac{n^j}{(n+1)\cdots(n+j)}$, summed
over $j\ge0$. Adding gives $e^n$. Then $\theta_n\frac{n^n}{n!}=\frac{e^n}2-\frac{n^n}{n!}Q=\frac{n^n}{n!}\bigl(\frac{Q+R}2-Q\bigr)$.

(b) The function $F(x)=\sum_{k\ge n}e^{-x}\frac{x^k}{k!}$ satisfies $F(0)=0$ and $F'(x)=\sum_{k\ge n}e^{-x}\bigl(\frac{x^{k-1}}{(k-1)!}-\frac{x^k}{k!}\bigr)=e^{-x}\frac{x^{n-1}}{(n-1)!}$ (telescoping). So
$\sum_{k\ge n}\frac{n^k}{k!}=\frac{e^n}{(n-1)!}\int_0^nt^{n-1}e^{-t}\dd t$, and with $t=ns$ and (a),
\[
R(n)=\frac{n!}{n^n}\cdot\frac{e^n}{(n-1)!}\,n^n\int_0^1s^{n-1}e^{-ns}\dd s=n\int_0^1s^{n-1}e^{n(1-s)}\dd s.
\]
The same computation with $\int_0^\infty t^{n-1}e^{-t}\dd t=(n-1)!$ gives
\[
Q+R=n\int_0^\infty s^{n-1}e^{n(1-s)}\dd s,
\]
hence the formula for $Q$.

(c) $\varphi$ is convex with minimum $\varphi(1)=0$, so $v$ is an increasing bijection, smooth with $v'(1)=1$ (since $\varphi(s)=\frac{(s-1)^2}2+O((s-1)^3)$), and its inverse is smooth. Since
$s^{n-1}e^{n(1-s)}=\frac{e^{-n\varphi(s)}}s$ and $v\,\mathrm dv=\mathrm d\varphi=\frac{s-1}s\mathrm ds$, we get $\frac{\mathrm ds}s=g(v)\,\mathrm dv$, which gives the two integrals. $g>0$ since $v$ and $s-1$ have the same sign.
For $s>1$, $\ln s\ge(s-1)-\frac{(s-1)^2}2$, so $\varphi\le\frac{(s-1)^2}2$, $v\le s-1$ and $g\le1$. For $s<1$ we need $|v|\le(1-s)(1+|v|)$, that is $|v|\le\frac{1-s}s=:w$; with $s=\frac1{1+w}$,
$2\varphi=2\ln(1+w)-\frac{2w}{1+w}$, and $k(w)=w^2-2\ln(1+w)+\frac{2w}{1+w}$ has $k(0)=0$ and $k'(w)=2w-\frac{2w}{(1+w)^2}\ge0$, so $2\varphi\le w^2$. Next, with $x=s-1$, $\varphi=\frac{x^2}2-\frac{x^3}3+O(x^4)$,
so $v=x(1-\frac x3+O(x^2))$, $x=v+\frac{v^2}3+O(v^3)$ and $g=\frac vx=1-\frac v3+O(v^2)$ near $0$. As $g$ is smooth, $|g(v)-1+\frac v3|\le Cv^2$ on $[-1,1]$; for $|v|\ge1$,
$|g(v)-1+\frac v3|\le2+2|v|\le4v^2$.

(d) With $I_j=n\int_0^\infty v^je^{-nv^2/2}\dd v$: $I_0=\sqrt{\frac{\pi n}2}$, $I_1=1$, $I_2=\frac{\sqrt{2\pi}}{2\sqrt n}$. So $Q=I_0-\frac13I_1+O(I_2)=\sqrt{\frac{\pi n}2}-\frac13+O(n^{-1/2})$, and the same on
$\left]-\infty,0\right]$ gives $R=\sqrt{\frac{\pi n}2}+\frac13+O(n^{-1/2})$. Hence $\theta_n=\frac{R-Q}2=\frac13+O(n^{-1/2})$.

(e) $e^{-n}\sum_{k\le n}\frac{n^k}{k!}=e^{-n}\frac{n^n}{n!}(Q+1)=\frac{Q+1}{Q+R}$. Now $Q+R=\frac{n!e^n}{n^n}=\sqrt{2\pi n}\,(1+O(\frac1n))$ (Volume~5, Theorem~13.8) and $Q+1=\frac12\sqrt{2\pi n}+\frac23+O(n^{-1/2})$, so the ratio is
$\frac12+\frac2{3\sqrt{2\pi n}}+O(\frac1n)$.

\tools{Volume~4, Exercise~6.5 (the limit $\frac12$); Volume~5, Theorem~13.8 (Stirling's series) and Theorem~11.9 (Laplace's method); Volume~1, Chapter~8 (inverse functions).}
\leads Numerically, one finds $n(\theta_n-\frac13)\to0.029630$, which is $\frac4{135}$: the next term of the expansion, which comes from the $v^2$ term of $g$. The same Laplace
integral, cut at other places, gives Knuth's analysis of linear probing: with a table of $m$ cells and $n$ keys, the cost of a search involves $Q$-like
sums at the ratio $\frac nm$, and the full table is exactly Ramanujan's case.
\end{solution}

\begin{solution}{pb:mappings}
\idea Following $x,f(x),f^2(x),\dots$ is drawing with replacement until a repeat, so the birthday sums reappear; the number of cycles weights the same terms
by $\frac1k$, which turns $\sqrt n$ into a logarithm.

\writeup (a) $x$ lies on a cycle of length exactly $k$ if and only if $x,f(x),\dots,f^{k-1}(x)$ are distinct and $f^k(x)=x$. The values $f(x),f(f(x)),\dots$ are fresh uniform draws as long as
no repeat occurs, so this has probability $\frac{n-1}n\cdots\frac{n-k+1}n\cdot\frac1n=\frac{n^{\underline k}}{n^{k+1}}$. Summing over $k$ and over the $n$ points gives $\mathbb E[\#\text{cyclic}]=\sum_k\frac{n^{\underline k}}{n^k}=Q(n)$,
and $Q(n)\sim\sqrt{\pi n/2}$ by Problem~\ref{pb:ramanujan} (or Volume~4, Exercise~2.7).

(b) The number of cycles is $\sum_x\frac{[x\text{ cyclic}]}{\text{length of its cycle}}$, so $C_n=n\sum_k\frac1k\cdot\frac{n^{\underline k}}{n^{k+1}}$.

(c) Let $w_k=\frac{n^{\underline k}}{n^k}$. \emph{Comparison.} $\ln w_k=\sum_{j<k}\ln(1-\frac jn)$. For $k\le n^{2/3}$ and $n\ge8$, $\frac jn\le\frac12$, and $-x-x^2\le\ln(1-x)\le-x$ gives
$\bigl|\ln w_k+\frac{k^2}{2n}\bigr|\le\frac k{2n}+\frac{k^3}{n^2}\le2$, hence $|w_k-e^{-k^2/(2n)}|\le Ce^{-k^2/(2n)}(\frac kn+\frac{k^3}{n^2})$. Then
$\sum_k\frac1ke^{-k^2/(2n)}(\frac kn+\frac{k^3}{n^2})=O(\frac{\sqrt n}n)+O(\frac{n^{3/2}}{n^2})\to0$. For $k>n^{2/3}$, $w_k\le w_{\lceil n^{2/3}\rceil}\le e^{-n^{1/3}/3}$ and $e^{-k^2/(2n)}\le e^{-n^{1/3}/2}$, so these terms total
at most $2ne^{-n^{1/3}/3}\to0$. So $C_n=S_n+o(1)$ with $S_n=\sum_{k\ge1}\frac1ke^{-k^2/(2n)}$.

\emph{The sum.} With $a=n^{-1/2}$ and $\chi(x)=\frac{e^{-x^2/2}-[x\le1]}x$, $S_n=\sum_{k\le\sqrt n}\frac1k+\sum_{k\ge1}a\chi(ka)$. The first sum is $\frac12\ln n+\gamma+o(1)$ (Volume~5, Corollary~13.6). The function
$\chi$ is bounded on $\left]0,1\right]$, monotone on $\left]0,1\right]$ and on $\left]1,\infty\right[$, and integrable, so its Riemann sums converge: $\sum_ka\chi(ka)\to\int_0^\infty\chi$. With $u=\frac{x^2}2$,
\[
\int_0^\infty\chi=\frac12\int_0^\infty\bigl(e^{-u}-[u\le\tfrac12]\bigr)\frac{\mathrm du}u=\frac12\Bigl(\int_0^\infty\bigl(e^{-u}-[u\le1]\bigr)\frac{\mathrm du}u+\int_{1/2}^1\frac{\mathrm du}u\Bigr)=\frac12(-\gamma+\ln2).
\]
The value $-\gamma$ comes from integrating by parts with $\ln u$, which turns the integral into $\int_0^\infty e^{-u}\ln u\dd u=\Gamma'(1)=-\gamma$ (Volume~5, Exercise~7.7). So
$C_n=\frac12\ln n+\gamma+\frac12(\ln2-\gamma)+o(1)=\frac12(\ln2n+\gamma)+o(1)$.

\tools{Volume~4, Chapter~1 (birthdays); Volume~5, Corollary~13.6 (harmonic numbers) and Exercise~5.7 ($\Gamma'(1)=-\gamma$).}
\leads Kruskal (1954) found (c), and Harris (1960) the distributions. The rho shape of the path $x,f(x),f^2(x),\dots$ (a tail then a cycle, both of length about
$\sqrt n$) is what Pollard's factoring method exploits. Numerically, one finds $C_n-\frac12(\ln2n+\gamma)\approx\frac{0.42}{\sqrt n}$.
\end{solution}

\begin{solution}{pb:xplusinv}
\idea Squaring linearises the recursion: $x_n^2$ grows by $2$ plus a small term $\frac1{x_n^2}$, which is about $\frac1{2n}$ and so adds up to a logarithm. Feeding the crude estimate back
gives the constant.

\writeup Let $y_n=x_n^2$; then $y_{n+1}=y_n+2+\frac1{y_n}$ and $y_0=1$.

(a) $y_{n+1}\ge y_n+2$ gives $y_n\ge2n+1$. Then $y_n=2n+1+\sum_{k=0}^{n-1}\frac1{y_k}\le2n+1+\sum_{k=0}^{n-1}\frac1{2k+1}\le2n+2+\int_0^{n-1}\frac{\mathrm dx}{2x+1}=2n+2+\frac12\ln(2n-1)$, which is at most
$2n+2+\ln n$ for $n\ge1$.

(b) Let $\delta_k=\frac1{y_k}-\frac1{2k+1}=\frac{2k+1-y_k}{y_k(2k+1)}$. By (a), $-\frac{1+\ln k}{(2k+1)^2}\le\delta_k\le0$ for $k\ge1$, and $\delta_0=0$, so $\sum_k\delta_k$ converges absolutely to some $\Delta$, and
$\sum_{k<n}\delta_k=\Delta+o(1)$. Also $\sum_{k=0}^{n-1}\frac1{2k+1}=H_{2n}-\frac12H_n=\ln2n+\gamma-\frac12(\ln n+\gamma)+o(1)=\frac12\ln n+\ln2+\frac\gamma2+o(1)$ (Volume~5, Corollary~13.6). So
\[
y_n=2n+\frac12\ln n+C+o(1),\qq C=1+\ln2+\frac\gamma2+\Delta.
\]

(c) With $a_n=\frac12\ln n+C+o(1)=O(\ln n)$, $\sqrt{2n+a_n}=\sqrt{2n}\bigl(1+\frac{a_n}{4n}+O(\frac{a_n^2}{n^2})\bigr)=\sqrt{2n}+\frac{a_n}{2\sqrt{2n}}+O\bigl(\frac{\ln^2n}{n^{3/2}}\bigr)$, and $\frac{a_n}{2\sqrt{2n}}=\frac{\ln n+2C}{4\sqrt{2n}}+o(\frac1{\sqrt n})$.

\tools{Volume~1, Chapter~3 (sequences, comparison with integrals); Volume~5, Corollary~13.6 (harmonic numbers).}
\leads Here $\Delta=-0.2586\ldots$, so $C=1.72314\ldots$; the constant depends on $x_0$, and we know no closed form for it. The same squaring trick handles $x_{n+1}=x_n+\frac1{x_n^p}$ (raise to the
power $p+1$). De Bruijn's \emph{Asymptotic Methods in Analysis} (1958) develops this bootstrapping for iterations in general.
\end{solution}

\begin{solution}{pb:hyperfactorial}
\idea $\ln H(N)=\sum k\ln k$ is minus the derivative of $\sum k^{-s}$ at $s=-1$. Differentiating the Euler--Maclaurin identity for $\zeta$ therefore gives the expansion of $\ln H$, with $\zeta'(-1)$ as
the constant.

\writeup Take $M=3$ in the identity above (valid for $\operatorname{Re}s>-5$) and differentiate in $s$; both sides are holomorphic there, and the integral may be differentiated under the
integral sign. At $s=-1$:
\begin{itemize}
\item $\frac{\mathrm d}{\mathrm ds}\sum_{k\le N}k^{-s}=-\sum k^{-s}\ln k$, which is $-\ln H(N)$;
\item $\frac{\mathrm d}{\mathrm ds}\frac{N^{1-s}}{s-1}=-\frac{N^{1-s}\ln N}{s-1}-\frac{N^{1-s}}{(s-1)^2}$, which is $\frac{N^2}2\ln N-\frac{N^2}4$;
\item $\frac{\mathrm d}{\mathrm ds}\bigl(-\frac{N^{-s}}2\bigr)=\frac{N^{-s}\ln N}2$, which is $\frac N2\ln N$;
\item $\frac{\mathrm d}{\mathrm ds}\bigl(\frac{B_2}2sN^{-1-s}\bigr)=\frac1{12}N^{-1-s}(1-s\ln N)$, which is $\frac1{12}(1+\ln N)$;
\item for $j=2,3$ the factor $(s)_{2j-1}$ vanishes at $s=-1$, so only its derivative survives: $(s)_3'(-1)=(-1)\cdot1=-1$, giving $\frac{B_4}{24}(-1)N^{-2}=\frac1{720N^2}$, and the $j=3$ term gives
$O(N^{-4})$;
\item likewise $(s)_6$ vanishes at $-1$, so the derivative of the remainder is
\[
-\frac{(s)_6'(-1)}{6!}\int_N^\infty\bar B_6(x)x^{-5}\dd x=O(N^{-4}).
\]
\end{itemize}
So
\[
\begin{aligned}\zeta'(-1)&=-\ln H(N)+\Bigl(\frac{N^2}2+\frac N2+\frac1{12}\Bigr)\ln N\\&\quad-\frac{N^2}4+\frac1{12}+\frac1{720N^2}+O(N^{-4}),\end{aligned}
\]
which is the claim with $\ln A=\frac1{12}-\zeta'(-1)$.

(b) $\zeta'(-1)=-0.1654211\ldots$, so $\ln A=0.2487545$ and $A=1.2824271$. For $n=10$, $\ln H(10)=\sum_{k\le10}k\ln k=102.0828\ldots$ and the formula without the $\frac1{720n^2}$ term is off by $1.39\cdot10^{-5}$, which is $\frac1{72000}$ to
three digits.

\tools{Volume~5, Proposition~13.7 and Chapter~11 (differentiating holomorphic functions).}
\leads Kinkelin (1860) and Glaisher (1878) found this constant; Glaisher's name stuck. With the functional equation one can show $\zeta'(2)=\frac{\pi^2}6(\gamma+\ln2\pi-12\ln A)$, so $A$ also
measures how $\zeta$ changes at $2$. $A$ is to $\prod k^k$ what $\sqrt{2\pi}$ is to $\prod k$: indeed $\ln\sqrt{2\pi}=-\zeta'(0)$, and the same differentiation at $s=0$ gives Stirling's formula.
\end{solution}

\begin{solution}{pb:banach}
\idea Finding a box empty means asking for it one time too many. The recursion between consecutive probabilities sums to an equation for the mean, which leaves a central
binomial coefficient.

\writeup (a) Suppose the left box is found empty. This happens when the left pocket is chosen for the $(N+1)$-th time, the right one having been chosen $N-r$ times before: $2N-r+1$
choices, the last one left. This has probability $\binom{2N-r}N2^{-(2N-r+1)}$, and doubling for the right box gives the formula.

(b) $\frac{a_{r+1}}{a_r}=2\binom{2N-r-1}N\binom{2N-r}N^{-1}=\frac{2(N-r)}{2N-r}$, which is the relation (also true at $r=N$ with $a_{N+1}=0$). Write it $2Na_{r+1}-ra_{r+1}=2Na_r-2ra_r$ and sum over
$0\le r\le N$. With $E=\mathbb E[R]=\sum ra_r$ and $\sum a_r=1$: the left side is $2N(1-a_0)-\sum_{r'\ge1}(r'-1)a_{r'}=2N(1-a_0)-E+(1-a_0)$, the right side is $2N-2E$. So $E=(2N+1)a_0-1$, with
$a_0=\binom{2N}N4^{-N}$.

(c) $a_0=\frac1{\sqrt{\pi N}}(1+O(\frac1N))$, so $E=\frac{2N+1}{\sqrt{\pi N}}-1+O(N^{-1/2})=2\sqrt{\frac N\pi}-1+O(N^{-1/2})$. For $N=50$: $E=7.04$, against $2\sqrt{50/\pi}-1=6.98$.

\tools{Volume~4, Chapter~2 (expectations); Volume~5, Proposition~13.12 (central binomial coefficients).}
\leads Feller's textbook made this ``Banach match box problem'' famous. The answer $2\sqrt{N/\pi}$ is the expected distance of a simple random walk from its start after about $2N$ steps,
which is what $R$ really is: the difference between the two boxes when one runs out.
\end{solution}

\begin{solution}{pb:primezeta}
\idea The logarithm of the Euler product is a sum of prime zeta values at multiples of $s$; Möbius inversion over the multiples solves for $P(s)$. The values $\zeta(ks)$ are
extremely close to $1$, so the series converges like $4^{-k}$ for $s=2$.

\writeup (a) For $s>1$, $\zeta(s)=\prod_p(1-p^{-s})^{-1}$ (Volume~5, Theorem~12.3), so $\log\zeta(s)=\sum_p\sum_{j\ge1}\frac{p^{-js}}j=\sum_{j\ge1}\frac{P(js)}j$, all terms being positive.

(b) Using (a) at $ks$,
\[
\sum_{k\ge1}\frac{\mu(k)}k\log\zeta(ks)=\sum_{k\ge1}\sum_{j\ge1}\frac{\mu(k)}{kj}P(jks)=\sum_{m\ge1}\frac{P(ms)}m\sum_{k\mid m}\mu(k)=P(s).
\]
The rearrangement is legitimate because $\sum_{k,j}\frac{P(jks)}{kj}=\sum_m\frac{\tau(m)P(ms)}m$, and $P(ms)\le\zeta(ms)-1\le2^{1-ms}$ for $ms\ge3$, so the series converges.

(c) For $\sigma\ge3$, $\log\zeta(\sigma)\le\zeta(\sigma)-1\le2^{-\sigma}+\int_2^\infty t^{-\sigma}\dd t=2^{-\sigma}\bigl(1+\frac2{\sigma-1}\bigr)\le2^{1-\sigma}$. For $k\ge K+1$ and $s\ge2$,
$ks\ge4$, so the error is at most
\[
\sum_{k>K}\frac{2^{1-ks}}k\le\frac{2^{1-(K+1)s}}{(K+1)(1-2^{-s})}\le\frac83\cdot\frac{2^{-(K+1)s}}{K+1}\le\frac4K2^{-(K+1)s}.
\]
For $s=2$ and $K=20$ this is $0.2\cdot2^{-42}<5\cdot10^{-14}$.

(d) By the prime number theorem, $\sum_{p>x}p^{-2}=\int_x^\infty t^{-2}\dd\pi(t)\sim\int_x^\infty\frac{\mathrm dt}{t^2\log t}\sim\frac1{x\log x}$, which is $6.2\cdot10^{-9}$ at $x=10^7$. And indeed
$0.452247420041-0.452247414181=5.86\cdot10^{-9}$. (The slightly smaller true value comes from the next term of $\frac1{x\log x}(1-\frac1{\log x}+\dots)$.)

\tools{Volume~5, Theorem~12.3 (Euler product) and Corollary~10.4 (Möbius inversion); Volume~6, Theorem~5.5.}
\leads The same inversion computes any sum $\sum_pf(p)$ for which $\sum_nf(n)$ has an Euler product: Mertens's constant and the twin prime constant are computed this way to hundreds of digits.
Summing the primes directly is hopeless beyond a few digits: the tail decays only like $\frac1{x\log x}$.
\end{solution}

\begin{solution}{pb:dickmanint}
\idea The delay equation $u\rho'(u)=-\rho(u-1)$ becomes an ordinary differential equation for the Laplace transform, because multiplying by $u$ is differentiating in $s$ and the delay is
multiplication by $e^{-s}$. The constant of integration is fixed at $s=\infty$ by $\rho(0)=1$, and Euler's constant appears from the asymptotics of the exponential integral.

\writeup (a) For $u>1$, $\frac{\mathrm d}{\mathrm du}(u\rho(u))=\rho(u)+u\rho'(u)=\rho(u)-\rho(u-1)=\frac{\mathrm d}{\mathrm du}\int_{u-1}^u\rho$, and both sides equal $1$ at $u=1$. If $\rho$ had a zero,
let $u_0>1$ be the first; then $\rho>0$ on $[u_0-1,u_0[$ and $u_0\rho(u_0)=\int_{u_0-1}^{u_0}\rho>0$, a contradiction. So $\rho>0$, and $\rho'(u)=-\frac{\rho(u-1)}u<0$ for $u>1$. For
$u\ge2$, $\rho$ decreases on $[u-1,u]$, so $u\rho(u)\le\rho(u-1)$ and $\rho(u)\le\frac12\rho(u-1)$. In particular $\rho(u)\le2^{2-u}$, so $\hat\rho(s)$ is finite and smooth for
$s>-\log2$, with derivatives obtained under the integral sign.

(b) On $]0,1[$, $u\rho(u)=u$ has derivative $1=\rho(u)-\rho(u-1)$, as $\rho(u-1)=0$; on $]1,\infty[$ this is (a). Since $u\rho(u)$ is continuous and decays exponentially,
\[
\int_0^\infty e^{-su}(u\rho(u))'\dd u=\Bigl[e^{-su}u\rho(u)\Bigr]_0^\infty+s\int_0^\infty ue^{-su}\rho(u)\dd u=-s\hat\rho'(s),
\]
while the left side is $\hat\rho(s)-\int_1^\infty e^{-su}\rho(u-1)\dd u=(1-e^{-s})\hat\rho(s)$.

(c) $\hat\rho>0$, and $(\log\hat\rho)'(s)=-\frac{1-e^{-s}}s$ for $s>0$; the right side extends continuously to $s=0$, and $\hat\rho$ is continuous on $[0,\infty[$, so integrating from $0$ gives the formula.

(d) $s\hat\rho(s)=s\int_0^1e^{-su}\dd u+s\int_1^\infty e^{-su}\rho(u)\dd u=1-e^{-s}+\theta$ with $0\le\theta\le s\int_1^\infty e^{-su}\dd u=e^{-s}$, so $s\hat\rho(s)\to1$. Next,
\[
\int_0^s\frac{1-e^{-t}}t\dd t=\int_0^1\frac{1-e^{-t}}t\dd t+\log s-\int_1^s\frac{e^{-t}}t\dd t\to\log s+\gamma+o(1),
\]
because $\int_0^1\frac{1-e^{-t}}t\dd t-\int_1^\infty\frac{e^{-t}}t\dd t=\gamma$. To see this, integrate by parts in $\Gamma'(1)=\int_0^\infty e^{-t}\log t\dd t=-\gamma$ (Volume~5, Exercise~7.7):
$\int_0^1e^{-t}\log t\dd t=\bigl[(1-e^{-t})\log t\bigr]_0^1-\int_0^1\frac{1-e^{-t}}t\dd t=-\int_0^1\frac{1-e^{-t}}t\dd t$ and $\int_1^\infty e^{-t}\log t\dd t=\int_1^\infty\frac{e^{-t}}t\dd t$. Finally, by (c),
\[
\int_0^\infty\rho=\hat\rho(0)=s\hat\rho(s)\cdot\exp\Bigl(\int_0^s\frac{1-e^{-t}}t\dd t-\log s\Bigr)\to1\cdot e^\gamma.
\]

(e) By Abel summation, $\sum_{n\ y\text{-smooth}}\frac1n=\int_{1^-}^\infty\frac{\mathrm d\Psi(t,y)}t=\int_1^\infty\frac{\Psi(t,y)}{t^2}\dd t$, since $\Psi(t,y)/t\to0$. With $\Psi(t,y)\approx t\rho(\frac{\log t}{\log y})$ and
$t=y^u$, $\frac{\mathrm dt}t=\log y\dd u$, this is $\approx\log y\int_0^\infty\rho(u)\dd u=e^\gamma\log y$, which is Mertens's third theorem. So $\int\rho=e^\gamma$ is the
consistency condition between Dickman's function and Mertens's theorem.

\tools{Volume~5, Exercise~7.7 ($\Gamma'(1)=-\gamma$); Laplace transforms and differentiation under the integral sign (Volume~5); Volume~6, Theorem~1.6.}
\leads The formula of (c) is the starting point of de Bruijn's asymptotic analysis of $\rho$, which gives $\rho(u)=e^{-u\log u-u\log\log u+u+o(u)}$ by a saddle point applied to the inverse Laplace
transform. Its counterpart for permutations is the probability that a random permutation has no cycle longer than $\frac nu$, which also tends to $\rho(u)$.
\end{solution}

\begin{solution}{pb:polymertens}
\idea In $\mathbb F_q[x]$ the prime number theorem is an identity, $\sum_{k\mid m}kN_k=q^m$, so the logarithm of Mertens's product is \emph{exactly} the harmonic sum, up to a tiny
error. Euler's constant enters through $H_n-\log n\to\gamma$, the polynomial counterpart of the integral $\int_0^1\frac{1-e^{-t}}t-\int_1^\infty\frac{e^{-t}}t$ of Problem~\ref{pb:dickmanint}.

\writeup (a) $-\log E_n=\sum_{k\le n}N_k\sum_{j\ge1}\frac{q^{-jk}}j$. The terms with $jk=m\le n$ give, for each $m$, $\frac{q^{-m}}m\sum_{k\mid m}kN_k=\frac1m$; together $H_n$.
The others have $k\le n<jk$, so $j\ge2$. Using $N_k\le\frac{q^k}k$, each is at most $\frac{q^{-(j-1)k}}{jk}$. For fixed $k$ the smallest admissible $j$, say $j_0$, has $(j_0-1)k\ge\frac n2$:
if $k\le\frac n2$ then $(j_0-1)k=\lfloor\frac nk\rfloor k>n-k\ge\frac n2$, and if $k>\frac n2$ then $j_0=2$ and $(j_0-1)k=k>\frac n2$. The series over $j$ has ratio at most $q^{-k}\le\frac12$, so
\[
0\le\varepsilon_n\le\sum_{k\le n}\frac{2q^{-n/2}}{j_0k}\le q^{-n/2}H_n\le(1+\log n)q^{-n/2}.
\]
Hence $E_n=e^{-H_n-\varepsilon_n}\sim e^{-\log n-\gamma}=\frac{e^{-\gamma}}n$.

(b) A reducible polynomial of degree $n$ has a factor of degree at most $\frac n2$, so the polynomials with no such factor are the irreducible ones, a proportion $\frac1n+O(\frac{q^{-n/2}}n)$
(Problem~\ref{pb:randpoly}). The prediction is $E_{\lfloor n/2\rfloor}\sim\frac{e^{-\gamma}}{n/2}=\frac{2e^{-\gamma}}n$. The ratio tends to $2e^{-\gamma}$.

(c) $H_n=\log n+\gamma+\frac1{2n}-\frac1{12n^2}+O(\frac1{n^4})$ (Euler--Maclaurin, Volume~5, Corollary~13.6), so $ne^\gamma E_n=\exp\bigl(-\frac1{2n}+\frac1{12n^2}+O(\frac1{n^4})-\varepsilon_n\bigr)=1-\frac1{2n}+O(\frac1{n^2})$.
At $n=10$, $\exp(-\frac1{20}+\frac1{1200})=0.952017$; for $q=5$, $\varepsilon_{10}\le(1+\log10)5^{-5}$ is small, and the numerical value $0.952015$ agrees.

\tools{Counting irreducible polynomials (Volume~3, Chapter~9); Volume~5, Corollary~13.6; Problem~\ref{pb:randpoly}.}
\leads The same factor $2e^{-\gamma}$ in both worlds shows that Legendre's failure is not a defect of the integers: it is what happens whenever one sieves an interval with ``primes'' up to the
square root of its size. The analogy between $\mathbb Z$ and $\mathbb F_q[x]$, with $\log x$ playing the role of the degree, runs through the whole of number theory; Rosen's book develops it.
\end{solution}

\section*{Sources and further reading}

\begin{itemize}
\item J.~M.~Borwein and P.~B.~Borwein, \emph{Pi and the AGM}, Wiley, 1987. Chapter~1 (the AGM, Gauss's theorem, Legendre's relation and the
algorithm for $\pi$) follows the same route as this chapter; Chapter~7 gives the logarithm by the AGM with error bounds.
\item D.~A.~Cox, \emph{The arithmetic-geometric mean of Gauss}, L'Enseignement Math\'ematique 30 (1984), 275--330. The history, Gauss's diary,
and the complex AGM, whose many limits are explained by modular forms.
\item D.~J.~Newman, \emph{A simplified version of the fast algorithms of Brent and Salamin}, Mathematics of Computation 44 (1985), 207--210.
\item E.~Salamin, \emph{Computation of $\pi$ using arithmetic-geometric mean}, Mathematics of Computation 30 (1976), 565--570; and R.~P.~Brent,
\emph{Fast multiple-precision evaluation of elementary functions}, Journal of the ACM 23 (1976), 242--251.
\item E.~A.~Marchisotto and G.-A.~Zakeri, \emph{An invitation to integration in finite terms}, College Mathematics Journal 25 (1994), 295--308.
\item D.~V.~Chudnovsky and G.~V.~Chudnovsky, \emph{Approximations and complex multiplication according to Ramanujan}, in \emph{Ramanujan
Revisited}, Academic Press, 1988, 375--472.
\end{itemize}
